\documentclass[11pt]{article}
\usepackage[T1]{fontenc}
\usepackage[utf8]{inputenc}
\usepackage{lmodern}
\usepackage{amsmath,amssymb,amsthm,booktabs,longtable,array}
\usepackage{mathtools}
\usepackage{mathrsfs}
\usepackage[a4paper,margin=27mm]{geometry}
\usepackage[hidelinks]{hyperref}
\usepackage{url}
\setlength{\emergencystretch}{3em}
\newcommand{\GT}{\widehat{\mathrm{GT}}}
\newcommand{\GTz}{\widehat{\mathrm{GT}}_0}
\newcommand{\Z}{\mathbf Z}
\newcommand{\Q}{\mathbf Q}
\newcommand{\F}{\mathbf F}
\newcommand{\HH}{\mathcal H}
\DeclareMathOperator{\Gal}{Gal}
\DeclareMathOperator{\Aut}{Aut}
\DeclareMathOperator{\Out}{Out}
\newenvironment{statement}[3]{\paragraph{#1 (#2).}\textbf{#3}\par}{}
\title{DELTA23: two--three frame reciprocity,\quad
  order-five characters, and the remaining pentagon gate}
\author{Research run DELTA23\\GPT-6 Astra Pro, four cooperating workers}
\date{Working manuscript, 20--21 November 2026, Europe/Berlin}
\begin{document}
\maketitle
\begin{abstract}
The main question is whether the supplied continuous character
$\Delta:\GTz\to\F_2$ vanishes on full profinite $\GT$.
This working version does not assert a decision of that question.
It gives explicit Hilbert-symbol formulas and their limitations,
an actual integral order-five frame character, the full image of
the normalized adelic two-frame map, an exact remaining kernel
for relative quadratic characters, and a marked two-word
Atkin--Lehner test.  A new affine free-group evaluation detects
the cubic sign in a group of order $216$ and reduces the main
question to twelve joint finite source residues, six of mixed
parity, in a joint source of order $5184$. A checked six-dimensional
Krammer quotient retains this cubic detector. All its $216$
reduced pentagon assignments have common-word realizations, and
the harmonic $S_4$ phase remains independent at that word level.
We also give an explicit genus-zero quotient diagram for the
cubic detector and a general limitation of fixed-target
distribution tests. A canonical cubic descent identifies the
half-character $\Psi_3/2$ with a descended cubic frame sign.
Ongoing full-GT comparisons are separated
from these completed proofs. All dependence on supplied
lifting and arithmetic inputs is explicit.
\end{abstract}
\tableofcontents

\section{Conventions, status, and source inputs}
Write $G=\GTz$, $\mathcal H=\ker\lambda$, and
$\mathcal H_5=\GT\cap\mathcal H$.  In sections using $H$
for this kernel it means $\mathcal H$; the level-six section
explicitly defines its own subgroup $H\subset\widehat B$.
The modular group is $B=C_2\amalg C_3$, with generators
$s,r$ and $\tau=sr$.  We retain the positive pure frames
$\beta_h,A_h$ of LM7\_PUSH E.65.  With its signed correction
the two parities are $a_2=\epsilon_2(W)$ and
$a_3=\epsilon_3(T)$, and $\delta=a_2+a_3$ on $\mathcal H$.
The full relative extension is $\Delta$.

An equation in a finite matrix group is not a proof that a
candidate lifts to an actual gentle automorphism or full-GT
element.  Conversely, the E.58 and LI.8 elements are used as
actual elements supplied by the attached runs, subject to their
stated lifting input.  The earlier D.50 lifting proof is not
present in the attachment.  It has not been silently upgraded
to an independently proved result of this run.

The following are the current main states.  The machine-readable
questions and check events accompany this paper.
\begin{center}\small
\begin{tabular}{p{.12\linewidth}p{.78\linewidth}}
\toprule
Question&Current result or gate\\\midrule
Item 1&Explicit compressed Hilbert realization; global-product
failure refuted; specified matrix-local and Kubota interpretations
separated.\\
A.84&Open: the six mixed states of D23.P10 must be excluded on
actual full-GT elements. The finite Krammer word test does not
exclude them; its full Nielsen closure is under analysis.\\
A.24&Actual order-five target map and exact Galois scalar
comparison constructed. The further full-GT comparison is open.\\
A.86&Explicit marked involution and two necessary and sufficient
word equations. Every congruence test is trivial; the full
noncongruence comparison remains open.\\
B.40&All characters factoring through the full adelic frame pair
are $0,\Delta$. Remaining gate is the relative elementary-two
image of its kernel, D23.C8.\\\bottomrule
\end{tabular}\end{center}

\subsection{Exact supplied inputs}
The notations R2, R3, R4 in the attached DELTA\_ON\_GT run
refer respectively to LM7\_PUSH E.58, the finite covariance
calculation, and E.65.  The proofs below give their exact
equations whenever a sign or a normalization is used.
Other supplied inputs used here are DELTA\_ON\_GT A.12--A.14,
A.25, A.62, AD.1--AD.3, AF.4--AF.6, LI.8 (integrated A.94),
and FP.3, FP.6, FP.17--FP.21, FP.32--FP.39.  Their source
files and SHA-256 values are listed in the source manifest.
The ledger distinguishes use as background from a proof
actually reread or independently checked here.

\part{Findings with a completed whole-proof check}
\section{Local symbols and the product formula}
% DELTA23 root proof draft. All uses of E.58/E.65 are explicit supplied inputs.
\subsection{What a two--three Hilbert-symbol interpretation proves}

Write $\mathcal G=\widehat{GT}_0$, $\mathcal H=\ker\lambda$, and
$\mathcal H_5=\widehat{GT}\cap\mathcal H$. We retain the signed
frame coordinates of LM7\_PUSH E.65 and the CRT map of
DELTA\_ON\_GT A.6. Thus $a_2=\epsilon_2(W)$,
$a_3=\epsilon_3(T)$ and $\delta=a_2+a_3$ on $\mathcal H$.
The additive character $\Delta$ and the multiplicative Dirichlet
character $\chi_{12}$ have different codomains: the correctly typed
identity is $(-1)^{\delta(h)}=\chi_{12}(u(h))$.

\paragraph{D23.L1 (a well-defined Hilbert realization of the compressed frames).}
For $h\in\mathcal H$, choose any rational number $n$ which is a unit
at $2$ and $3$ and has residues
\[
n\equiv(-1)^{a_2(h)}\pmod4,\qquad
n\equiv(-1)^{a_3(h)}\pmod3.
\]
In particular one may choose the unique $n\in\{1,5,7,11\}$.
Define $L_2(h)=(n,3)_2$ and $L_3(h)=(n,3)_3$, using the quadratic
Hilbert symbols over $\mathbf Q_2$ and $\mathbf Q_3$. These are
independent of the indicated choice of $n$, are continuous characters,
and satisfy
\[
L_2=(-1)^{a_2},\qquad L_3=(-1)^{a_3},\qquad
L_2L_3=(-1)^\delta=\chi_{12}(u).
\tag{D23.L1.1}
\]
For the E.58 element $h_*$, in the supplied convention
$(a_3,a_2)=(1,0)$, one has $n=5$ and $(L_2,L_3)=(1,-1)$.
\begin{proof}
Use the standard local Hilbert formulas, with the precise locators
Voight, \emph{Quaternion algebras}, (12.4.10), 12.4.12(b), and
(12.4.14). For a $2$-adic unit $n$ they give
$(n,3)_2=(-1)^{(n-1)/2}$, where the exponent is read modulo two.
For a $3$-adic unit $n$ they give $(n,3)_3=(n/3)$.
These depend only on $n\bmod4$ and $n\bmod3$, respectively, and
give the first two identities. Their product gives the third by
the supplied definition of $\delta$. The four residues give
$\chi_{12}$ with values $1,-1,-1,1$, in the order $1,5,7,11$.
Continuity and multiplicativity follow from the already proved
frame characters, not from a purported multiplicative choice of
rational representatives. Finally E.58.23 and E.65.2--4 give
$T=-I$ at three and $W=I$, so the last evaluation follows.
\end{proof}

This is a positive answer for an explicitly defined pair of symbols
of the \emph{compressed} frame data. The definition has used the
two frame parities. It is not an independent arithmetic comparison
and does not prove that the two symbols have equal values.

\paragraph{D23.L2 (the proposed failure of global reciprocity is false).}
The rational pair $(n,3)$ in D23.L1 obeys Hilbert reciprocity for
every frame value. For the four positive representatives its complete
ramification table is
\[
\begin{array}{c|cc|c|c}
n&(n,3)_2&(n,3)_3&\text{other ramified finite primes}&(n,3)_\infty\\\hline
1&1&1&\varnothing&1\\
5&1&-1&5&1\\
7&-1&1&7&1\\
11&-1&-1&\varnothing&1
\end{array}
\tag{D23.L2.1}
\]
In particular E.58 does not violate a global product formula: its
extra factor is $(5,3)_5=-1$.
\begin{proof}
At an odd prime not dividing $3n$, both arguments are units and
the Hilbert symbol is one by Voight 12.4.12(a). At a prime dividing
one of $5,7,11$ the symbol is $(3/p)$. The square residues give
$(3/5)=(3/7)=-1$ and $(3/11)=1$ (since $5^2=3\pmod{11}$).
The two local values are D23.L1, and the real symbol is one
because $3>0$. These facts give the entire table. Each row has
an even number of minus signs, in agreement with Hilbert
reciprocity, Voight Proposition 14.2.1 (or Corollary 14.6.2).
\end{proof}

\paragraph{D23.L3 (the missing support hypothesis is exactly the old parity gate).}
For two bits $(b_2,b_3)$, the following are equivalent:
\begin{enumerate}
\item $b_2=b_3$;
\item there is $a\in\mathbf Q^\times$, a unit at $2$ and $3$, with
$(a,3)_2=(-1)^{b_2}$ and $(a,3)_3=(-1)^{b_3}$, and with
$(a,3)_v=1$ at every other place;
\item there is a rational quaternion algebra whose local invariants
at $2,3$ are $b_2/2,b_3/2$ and which is split at every other place.
\end{enumerate}
For the two equal-bit cases the algebras can be chosen as
$M_2(\mathbf Q)$ and the discriminant-six algebra $(2,3)$.
Thus the existence of such a two-place globalization for every
$h\in\mathcal H_5$ is \emph{equivalent} to A.84; it is not an
already proved arithmetic reason for A.84.
\begin{proof}
The implication from the third condition to the first is Hilbert
reciprocity (or the even-ramification part of Voight Main Theorem
14.6.1). The second condition implies the third by taking $(a,3)$.
For equal bits choose $a=1$ or $a=-1$. The local formulas give,
for $a=-1$, minus signs exactly at $2$ and $3$, and the real symbol
is again positive. This proves the first implication to the second.
Finally $(-1,3)$ is isomorphic to $(2,3)$. For an explicit change
of generators in $(-1,3)$, write $i^2=-1$, $j^2=3$, $ij=-ji$.
Then $i'=(1+j)i$ has square
$(1+j)(1-j)i^2=(1-3)(-1)=2$ and anticommutes with $j$.
Since $1+j$ is invertible, $i',j$ generate the original algebra.
This proves the asserted identification without any appeal merely
to coincident discriminants. Applying the equivalence to
$(b_2,b_3)=(a_2(h),a_3(h))$ gives the final assertion.
\end{proof}

\paragraph{D23.L4 (relative pro-three frames cannot recover the cubic bit).}
Let $U=F'F^3$ and let $\Gamma_3=B/O^3(U)$ be the specified
relative maximal pro-three quotient used in E.58. There is no
function of the induced automorphism of $\Gamma_3$ together with
the two retained pure frame values there which equals $a_3$ on
all of $\mathcal H$. In particular the cubic bit cannot be a
normalized local symbol depending only on that data.
\begin{proof}
The identity element has identity action and identity frames, with
$a_3=0$. The actual element supplied by E.58 has identity action
on every relative local factor. By E.58.20 its two frame values
in the factor at three are both one, while E.58.23 gives $a_3=1$.
The input data for the proposed function are equal and the
required outputs differ. This is a contradiction.
\end{proof}

The quotient $\Gamma_3$ in D23.L4 must not be confused with the
matrix target $\operatorname{SL}_2(\mathbf Z_3)$. The latter has
non-three-primary finite subgroups; reduction modulo three already
retains a quaternion group. D23.L4 refutes a precise pro-three
interpretation, not every possible arithmetic local comparison.

\paragraph{D23.L5 (full multipliers and the cyclotomic distinction).}
For $g\in\mathcal G$ choose $\sigma\in G_{\mathbf Q}$ with
$\chi(\sigma)=\lambda(g)$ and put $h=g\sigma^{-1}$.
The product $L_2(h)L_3(h)=(-1)^{\Delta(g)}$ is independent of this
choice, while the individual symbols need not be. Moreover neither
of the two individual bits is the restriction of a Dirichlet
character of the multiplier.
\begin{proof}
Choice independence of the product is E.65.5--6. The source
DELTA\_ON\_GT A.25 constructs $\sigma_0\in G_{\mathbf Q^{\rm ab}}$
which sends $r_0=(3/4)^{1/4}$ to $-r_0$, with
$a_2(\sigma_0)=a_3(\sigma_0)=1$ and $\lambda(\sigma_0)=1$.
For $g=1$, taking $\sigma=1$ gives the two symbols $(1,1)$,
whereas taking $\sigma=\sigma_0$ gives $(-1,-1)$. This proves
the individual ambiguity. A character of the multiplier is
trivial on every multiplier-one element, while each of these
frame bits takes value one at $\sigma_0$. This proves the last
assertion. The cited example concerns a nonabelian quartic-root
extension over $\mathbf Q$, quadratic only after adjoining
$\mathbf Q^{\rm ab}$; it is not the cyclotomic quadratic field
$\mathbf Q(\sqrt3)$.
\end{proof}

\subsection{A precise Kubota test which does not detect the discrepancy}

\paragraph{D23.L6 (the integral splitting-normalized test is trivial).}
Use the actual signed adelic frames $\mathcal T_h,\mathcal W_h$
of DELTA\_ON\_GT AD.2.  For a prime $p$, and a matrix
$g=\left(\begin{smallmatrix}a&b\\c&d\end{smallmatrix}\right)
\in\operatorname{SL}_2(\mathbf Q_p)$, put
\[
 x(g)=\begin{cases}c&c\ne0,\\d&c=0,\end{cases}\qquad
 a_p(g,h)=(x(g),x(h))_p
             (-x(g)^{-1}x(h),x(gh))_p.
\]
This is Kubota's factor set (Kubota 1967, equation (2), printed
p.~114).  Fix the splitting normalization
\[
 s_p(g)=\begin{cases}
 (c,d)_p&cd\ne0\text{ and }v_p(c)\text{ is odd},\\
 1&\text{otherwise},
 \end{cases}
 \qquad
 \kappa_p(g,h)=a_p(g,h)s_p(g)s_p(h)s_p(gh)^{-1}.
\tag{D23.L6.1}
\]
Then the concrete functions of the corrected frame pair
\[
 K_p(h)=\kappa_p(\mathcal T_{h,p},\mathcal W_{h,p}),
 \qquad p=2,3,
\]
are both identically $1$ on the entire gentle multiplier-one
group.  In particular $K_2K_3$ is not $(-1)^\delta$: the
stipulated actual E.58 element has $\delta=1$.

\begin{proof}
The standard local splitting theorem says that (D23.L6.1) is
one on $\Gamma_1(4,\mathbf Z_2)$ at two and on
$\operatorname{SL}_2(\mathbf Z_p)$ at every odd prime.
A primary source giving precisely this normalization is Hida,
\emph{Adjoint L-value as a period integral and the mass formula
of Siegel--Shimura}, version 20 September 2023, Section 1.3,
printed p.~7, the displayed definition of $s_p$ and paragraph
immediately before (1.5); it cites Gelbart, \emph{Weil's
Representation and the Spectrum of the Metaplectic Group},
Proposition 2.8.  We use this published splitting statement,
not a computation of a few residue classes.

AD.2 gives $\mathcal T_{h,2}\equiv\mathcal W_{h,2}\equiv I$
modulo four, so both matrices lie in the indicated group at
two.  At three both are integral matrices of determinant one.
The two conclusions follow directly from the respective
splittings.  E.58.23 and E.65 give $\delta(h_*)=1$, providing
the asserted failure of this proposed comparison.
\end{proof}

This refutes the specified test, not every conceivable local
construction.  A factor set changes by a coboundary when its
section is changed: replacing a section by $b(g)$ times that
section multiplies its factor set by
$b(g)b(h)b(gh)^{-1}$.  Thus an unqualified phrase ``the Kubota
symbol of two frames'' does not specify a function until the
matrix lifts and the section are fixed.  The explicit
normalization above fixes both and has the stated outcome.

\subsection{The two proposed quaternion algebras are different}
\paragraph{D23.L7.}
The rational quaternion algebras $(-1,-3)$ and $(2,3)$ are
not isomorphic.  Their ramification sets are, respectively,
$\{3,\infty\}$ and $\{2,3\}$.  Thus the former is definite
and has finite discriminant three; the latter is indefinite
and has discriminant six.  An identification between them cannot
serve as the missing two--three frame comparison.
\begin{proof}
For $(-1,-3)$ the odd-prime Hilbert formula is one away from
three, and at three it is the quadratic residue symbol
$(-1/3)=-1$.  At two the two odd units have residues $7,5$
modulo eight; the exponent in Voight (12.4.14) is
$(({-1}-1)/2)(({-3}-1)/2)=2$, which is even.
The real symbol is negative because both arguments are negative.
This gives its entire ramification set.  The ramification set
of $(2,3)$ was calculated in D23.L3.  A rational algebra
isomorphism would induce an isomorphism over $\mathbf Q_2$,
where the first algebra is split and the second is division,
which is impossible.
\end{proof}
This statement does not dispute a supersingular endomorphism
interpretation of $(-1,-3)$.  It identifies the extra comparison
which would be needed before that interpretation could bear on
the discriminant-six proposal.

\subsection{Direct local formulas for the actual matrix frames}
\paragraph{D23.L8.}
Write the three-components of the two signed corrected adelic
frames of $h\in\mathcal H$ as
\[
 \mathcal T_{h,3}=aI+bR_0,\quad a^2-ab+b^2=1,
 \qquad
 \mathcal W_{h,3}=cI+dS,\quad c^2+d^2=1.
\]
The two quantities $a+b$ and $c^2-d^2$ belong to
$\mathbf Z_3^\times$, and one has the direct matrix formulas
\[
 (a+b,3)_3=(-1)^{a_3(h)},\qquad
 (c^2-d^2,3)_3=(-1)^{a_2(h)}.
\tag{D23.L8.1}
\]
Their product is $(-1)^{\delta(h)}$.  These formulas use
actual local matrix coefficients, without first naming the bits,
but both are formulas at three.  They do not turn the harmonic
frame bit into a function of the two-adic frames; D23.C7 rules
out that latter factorization.
\begin{proof}
Modulo three the first norm equation is $(a+b)^2=1$.
Also $R_0$ acts as one on its one-dimensional fixed space,
so $a+b\bmod3$ is exactly the signed fixed-line eigenvalue
defining $a_3$.  For the second norm equation, the only pairs
$(c,d)\in\mathbb F_3^2$ with $c^2+d^2=1$ are
$(1,0),(-1,0),(0,1),(0,-1)$.  They represent respectively
$I,-I,S,-S$.  The first two have even $S$-exponent and the
last two have odd $S$-exponent.  The value $c^2-d^2$ is
respectively $1,1,-1,-1$.  This proves both the unit assertions
and the equality of the residue signs with the indicated frame
signs.  For a three-adic unit $u$, the odd-prime Hilbert formula
gives $(u,3)_3=(u/3)$, which in $\mathbb F_3^\times$ is
precisely this sign.  The product assertion follows.
For E.58, the two matrices are $-I,I$ and the values are
therefore $-1,1$, as in D23.L1.
\end{proof}


\section{Integral order-five construction}
% DELTA23: full proof draft, root, 2026-11-21 CET.
\subsection{Integral covariance of the five-strand symplectic representation}

\paragraph{D23.T1 (integral specialization of the rigidity method).}
Let $J$ be the alternating $4\times4$ integral matrix with
$J_{i,i+1}=1$, $J_{i+1,i}=-1$ and all other entries zero.
Put $t_v=I+vv^{\mathsf t}J$ and $t_i=t_{e_i}$.
The assignment $q(\sigma_i)=t_i$ defines a continuous homomorphism
$q:\widehat B_5\to\operatorname{Sp}(J,\widehat{\mathbf Z})$.
For every multiplier-one element $h\in\widehat{GT}$ there is a unique
$V_h\in\operatorname{Sp}(J,\widehat{\mathbf Z})$ such that
\[
V_he_1=e_1,\qquad q\Phi_h=\operatorname{Ad}(V_h)q.
\]
The map $h\mapsto V_h$ is a continuous homomorphism.
This statement concerns the genuine five-strand braid action; it does
not assert that this representation descends to the five-punctured
sphere group.

\begin{proof}
The determinant of $J$ is one. The transvections preserve $J$.
For $c=v^{\mathsf t}Jw$, direct multiplication gives
\[
t_vt_w-t_wt_v=c(vw^{\mathsf t}+wv^{\mathsf t})J,
\qquad
t_vt_wt_v-t_wt_vt_w=(1-c^2)(vv^{\mathsf t}-ww^{\mathsf t})J.
\tag{D23.T1.1}
\]
The entries of $J$ therefore give the braid and distant commutation
relations for the $t_i$. The discrete homomorphism extends to the
profinite completion, since reduction modulo each integer has finite
target.

Fix a prime $\ell$. The standard GT action has
$\Phi_h(\sigma_1)=\sigma_1$ and sends every other $\sigma_i$ to a
specified conjugate of itself. Hence
$q_\ell\Phi_h(\sigma_i)=t_{v_i}$ with $v_i$ primitive over
$\mathbf Z_\ell$ and $v_1=e_1$. The $v_i$ can be chosen continuously
using the specified conjugators in the GT braid formula.

We may first work over the field $\mathbf Q_\ell$. If $t_v,t_w$
commute, the first identity in (D23.T1.1) implies $c=0$.
Indeed, if $c\ne0$, it would give $vw^{\mathsf t}+wv^{\mathsf t}=0$.
Choose a linear functional $L$ with $L(v)=1$ and multiply this matrix
on the left by $L$; then $w=-L(w)v$, which implies $c=0$, a
contradiction. For a pair satisfying the braid relation, the second
identity shows that either $c^2=1$ or $t_v=t_w$.

An equality $t_{v_i}=t_{v_{i+1}}$ is impossible. A neighboring
generator $t_{v_{i+2}}$, when present, commutes with $t_{v_i}$ and
braids with $t_{v_{i+1}}$. Two commuting invertible elements which
braid are equal: $a^2b=b^2a$ implies $a=b$. Equality thus propagates
to the right, and the identical argument propagates it to the left.
All four generators would have equal image. But $\Phi_h$ is an
automorphism of $\widehat B_5$, so $q_\ell\Phi_h$ has exactly the
same image as $q_\ell$, and the latter is nonabelian, as witnessed by
$t_1,t_2$. This is a contradiction.

Consequently $c_i=v_i^{\mathsf t}Jv_{i+1}$ equals $1$ or $-1$,
and the pairings of distant vectors are zero. Define
$\epsilon_1=1$ and $\epsilon_{i+1}=\epsilon_i c_i$ and replace
$v_i$ by $w_i=\epsilon_i v_i$. The transvections do not change and
$w_i^{\mathsf t}Jw_{i+1}=1$. Thus the column matrix
$V_\ell=(w_1,w_2,w_3,w_4)$ satisfies
$V_\ell^{\mathsf t}JV_\ell=J$. Taking determinants shows
$\det(V_\ell)^2=1$, so $V_\ell$ is invertible over
$\mathbf Z_\ell$. It fixes $e_1$ and conjugates every $t_i$ to
$t_{v_i}$.

For uniqueness, the simultaneous centralizer of the $t_i$ in
$\operatorname{GL}_4(\mathbf Q_\ell)$ is the scalar group. To see
this, commuting with $t_i-I$ preserves its image $\mathbf Q_\ell e_i$,
so a centralizing matrix is diagonal. Its commutation with the
nonzero $(i,i+1)$ entries of $t_i-I$ equates all adjacent diagonal
entries. A scalar symplectic matrix is $\pm I$, and fixing $e_1$
chooses $I$.

The product $V_h=(V_\ell)_\ell$ lies in the indicated adelic
symplectic group. The displayed formulas using the original
continuous $v_i$ and their pairings prove its continuity, prime by
prime, including at $2$ and $5$. Both $V_{hk}$ and $V_hV_k$ fix
$e_1$ and intertwine $q\Phi_h\Phi_k$ with $q$. Uniqueness gives
$V_{hk}=V_hV_k$.
\end{proof}

\paragraph{D23.T2 (the exact integral order-five centralizer).}
Put $C=t_1t_2t_3t_4$. Then
\[
C=\begin{pmatrix}0&0&0&1\\-1&0&0&1\\0&-1&0&1\\0&0&-1&1\end{pmatrix},
\quad C^5=-I,\quad \det(XI-C)=X^4-X^3+X^2-X+1.
\]
For every commutative ring $R$, the endomorphism centralizer of
$C$ in $M_4(R)$ is $R[C]\simeq R[\zeta_5]$, with
$\zeta_5$ represented by $-C$. Its symplectic centralizer is
\[
\operatorname{Cent}_{\operatorname{Sp}(J,R)}(C)
=\{z\in R[\zeta_5]:z\overline z=1\},
\qquad \overline{\zeta_5}=\zeta_5^{-1}.
\]
\begin{proof}
Multiplication gives the matrix and polynomial identities. The four
vectors $e_1,Ce_1,C^2e_1,C^3e_1$ are respectively
$e_1,-e_2,e_3,-e_4$, so they form a basis over every $R$.
Any endomorphism commuting with $C$ is determined by its value on
$e_1$; writing that value in this basis shows that the endomorphism
is a unique polynomial of degree at most three in $C$. This proves
the asserted algebra centralizer without any restriction on the
residue characteristic.
For the symplectic adjoint $X^*=J^{-1}X^{\mathsf t}J$ one has
$C^*=C^{-1}$, since $C$ is symplectic. On the commutative algebra
$R[C]$ the adjoint is therefore exactly the displayed involution.
The symplectic condition is $z^*z=1$. It already implies that $z$
is a unit, and is precisely the norm-one equation.
\end{proof}

% DELTA23, agent /root/five_level_six.  New algebraic results.
% Model: GPT-6 Astra Pro. Effort inherited, not exposed in the runtime.
% Run start: 2026-11-20 22:59 UTC.  This fragment is not a final certification.

\subsection*{D23.F1. A rigid order-five frame with one distinguished pole}
Let $L_6=(\widehat\Gamma_{0,[6]})_6$ be the subgroup fixing the
sixth label, and use the standard identification
\[
 L_6=\widehat B_5/\overline{\langle z_5\rangle},\qquad
 z_5=(\sigma_1\sigma_2\sigma_3\sigma_4)^5.
\]
Put $a=\sigma_1\sigma_2\sigma_3\sigma_4$ and write $\bar a$
for its image in $L_6$.  Let $\mathcal H_5$ denote the
multiplier-one subgroup of full profinite $\widehat{GT}$.
There is a unique continuous map
\[
 A_5:\mathcal H_5\longrightarrow\widehat{PB}_5
\]
such that
\[
 \Phi_h(a)=A_5(h)aA_5(h)^{-1},\qquad
 \ell_{12}(A_5(h))=0,
 \tag{D23.F1.1}
\]
where $\ell_{12}$ is the integral linking coordinate for the first
two strands, normalized by $\ell_{12}(z_5)=1$.
Its composition rule is
\[
 A_5(hk)=\Phi_h(A_5(k))A_5(h).
 \tag{D23.F1.2}
\]
The image of $A_5(h)$ in $P_5$ after forgetting the distinguished
sixth point is the unique order-five frame of
DELTA\_ON\_GT, FP.6.

\begin{proof}
First $C_{P_6}(\bar a)=1$.  Forget the sixth point.  The image of a
pure element centralizing $\bar a$ lies in $C_{P_5}(a_5)=1$,
by DELTA\_ON\_GT FP.3.  Hence it lies in the free kernel
$\widehat F_4$ of the pure Birman sequence.  Geometrically the
remaining punctures are $\mu_5$, the sixth point is $\infty$, and
the action is rotation $z\mapsto\zeta_5 z$.  The quotient orbifold
of $\mathbb P^1\setminus\mu_5$ has one puncture and two cone
points of order five.  Thus
\[
 \widehat F_4\rtimes\langle\bar a\rangle
       \simeq C_5\amalg C_5,
\]
with the section at $\infty$ a factor.  The factor-centralizer
theorem (Herfort--Ribes 1985, Theorem B', Section 4, printed
p.~159, also used in FP.3) makes its centralizer meet
$\widehat F_4$ trivially.  This proves the assertion.  The
profinite Birman exactness can be justified without assuming
the five-point base group is free. After fixing three labels
at $0,1,\infty$, the discrete forgetful map is
$\operatorname{Conf}_3(\mathbb C\setminus\{0,1\})\to
\operatorname{Conf}_2(\mathbb C\setminus\{0,1\})$.
It has a continuous section choosing its extra real point
farther from the origin than the absolute values of all the
existing points plus two. Its fundamental-group extension
therefore splits and has finitely generated free kernel.
Cofinal characteristic finite-index subgroups of that kernel
are invariant under the base action, and the induced action
on each finite quotient has finite image. Taking a finite-index
subgroup of the base that kills that image proves that the
topology induced on the kernel is its full profinite topology.
The split finite quotients also separate the base coordinate,
so completion gives the asserted exact sequence. The same
finite-extension argument applies to the subgroup over the
order-five rotation.

Collas 2012, Theorem B, applies to the order-five element $\bar a$
in the full six-point group.  For $h\in\mathcal H_5$ it gives
$\Phi_h(\bar a)=g\bar a g^{-1}$.  The standard GT action is the
identity on the permutation quotient.  The permutation of $g$
therefore centralizes a five-cycle with one fixed label.  It
fixes that unique sixth label and is a power of the five-cycle.
Multiplying $g$ on the right by an appropriate power of $\bar a$
gives a pure conjugator $\alpha_5(h)\in P_6$.
The preceding centralizer calculation proves uniqueness.
The continuous injective orbit map from the compact group $P_6$
to the Hausdorff group $G_6$ is a homeomorphism onto its image;
therefore this frame is continuous.

The central extension of pure groups splits by $\ell_{12}$.
For clarity, on ordered configurations one has the algebraic
isomorphism
\[
 \operatorname{Conf}_5(\mathbb A^1)
 \simeq\mathbb A^1\times\mathbb G_m\times\mathcal M_{0,6},
\quad
 (z_i)\longmapsto
 \left(z_1,z_2-z_1,
       \frac{z_3-z_1}{z_2-z_1},\ldots,
       \frac{z_5-z_1}{z_2-z_1}\right).
\]
Here the last factor has the first two points normalized to $0,1$
and the distinguished sixth point to $\infty$.
The $\mathbb G_m$ loop is the full twist, and its winding number
is $\ell_{12}$.  Hence
\[
 \widehat{PB}_5\simeq
 \overline{\langle z_5\rangle}\times\ker\ell_{12},
 \qquad \ker\ell_{12}\xrightarrow{\sim}P_6.
\]
Lift $\alpha_5(h)$ through this last isomorphism; call the lift
$A_5(h)$.  The frame equation in $L_6$ lifts to
\[
 \Phi_h(a)=A_5(h)aA_5(h)^{-1}z_5^e
\]
for some $e\in\widehat{\mathbb Z}$.  The exponent-sum map on
$\widehat B_5$ sends $a$ to four and $z_5$ to twenty, and GT
acts on this map by its multiplier, which is one.  Thus $20e=0$.
Multiplication by twenty on $\widehat{\mathbb Z}$ is injective,
so $e=0$.  This proves (D23.F1.1).

At multiplier one, $\Phi_h$ is the identity on the abelianization
of the pure braid group.  Indeed, its standard generators are
conjugates of half-twist squares; their images are conjugates
of the same squares, with the same strand-pair permutation.
Their linking coordinates are consequently unchanged.
It follows that $\Phi_h(A_5(k))A_5(h)$ has linking coordinate
zero and is a conjugator for $\Phi_{hk}(a)$.  Uniqueness gives
(D23.F1.2).  Finally the forgetful map takes its frame equation
to the five-point frame equation, and pure uniqueness there
identifies the two frames.
\end{proof}

\subsection*{D23.F2. The actual order-five norm-one frame character}
Use the symplectic braid representation and normalized covariance
$q\Phi_h=\operatorname{Ad}(V_h)q$ of D23.T1--D23.T2.
Set $C=q(a)$ and
\[
 \mathcal C_5(h)=q(A_5(h))^{-1}V_h.
 \tag{D23.F2.1}
\]
Then $\mathcal C_5$ is a continuous homomorphism on
$\mathcal H_5$, taking values in
\[
 \operatorname{Cent}_{\operatorname{Sp}_4(\widehat{\mathbb Z})}(C)
  \simeq
 \{u\in(\widehat{\mathbb Z}[\zeta_5])^\times:
                    u\overline u=1\}.
 \tag{D23.F2.2}
\]
The torus is the order-five target of DELTA\_ON\_GT FP.19;
here it receives an actual frame character.  The source is the
six-point stabilizer, so this does not assert that the
symplectic representation descends to the five-point sphere
group refuted in FP.21.

\begin{proof}
Apply $q$ to (D23.F1.1) and use covariance:
$V_h C V_h^{-1}=q(A_5(h))Cq(A_5(h))^{-1}$.
This proves centralization.  The precise integral centralizer
is D23.T2: $C^5=-I$, $-C$ has polynomial $\Phi_5$, is cyclic
over every coefficient ring, and its symplectic adjoint is its
inverse.  Thus commuting symplectic matrices are exactly the
displayed norm-one units.

For composition, write $Q_h=q(A_5(h))$.  Equation (D23.F1.2)
gives $Q_{hk}=V_hQ_kV_h^{-1}Q_h$, while D23.T1 gives
$V_{hk}=V_hV_k$.  Hence
\[
 Q_{hk}^{-1}V_{hk}
 =Q_h^{-1}V_hQ_k^{-1}V_k
 =\mathcal C_5(h)\mathcal C_5(k).
\]
Continuity follows from the continuity of all its factors.
\end{proof}

\paragraph{Current arithmetic comparison gate.}
The construction of $\mathcal C_5$ is algebraic and unconditional
on full $\mathcal H_5$ once the stated published torsion and
standard GT-action inputs are granted.  Equality with the
Jacobian Tate action requires accounting for the normalization
$z_1=0,z_2=1$.  At the ordered CM point
$(1,\zeta_5,\ldots,\zeta_5^4)$ put $d=\zeta_5-1$ and
$t=(x-1)/d$.  The normalized monic hyperelliptic equation is
\[
 y_0^2=\prod_{j=0}^{4}
 \left(t-\frac{\zeta_5^j-1}{d}\right),
 \qquad y=d^{5/2}y_0.
\]
Thus this normalized curve is a quadratic twist of
$y^2=x^5-1$ by $d^5$, whose Kummer class is the class of $d$.
The exact scalar comparison is not to be suppressed.

% DELTA23 /root/five_level_six; proof draft for independent checking.

\subsection*{D23.F3. An original fifth-fibre Tate character on the gentle group}
Let $F=\widehat F_2=\langle x,y\rangle$ with $z=(xy)^{-1}$,
and let $\mathcal H=\ker\lambda\subset\widehat{GT}_0$ act by
$\alpha_h(x)=x$, $\alpha_h(y)=f_h^{-1}yf_h$.
Consider the cyclic quotient
\[
 F\twoheadrightarrow C_5\times C_2,
 \qquad x\mapsto(1,0),\quad y\mapsto(0,1),
\]
with kernel $K_{10}$, and let $N$ be the closed normal subgroup
of $F$ generated by $x^5,y^2,(xy)^{10}$.
Then $N\subset K_{10}$ and
\[
 \Pi_5=K_{10}/N
\]
is the completed fundamental group of the compact genus-two
curve $C_5:y^2=x^5-1$.  Every $h\in\mathcal H$ induces an
automorphism of $\Pi_5$ and a continuous character
\[
 \Omega_5:\mathcal H\longrightarrow
 \{u\in(\widehat{\mathbb Z}[\zeta_5])^\times:
                         u\overline u=1\}.
 \tag{D23.F3.1}
\]
On $G_{\mathbb Q^{\rm ab}}$ this is the full Tate representation
of $J_5=\operatorname{Jac}(C_5)$, in the fixed covering marking.
If $f_h\in F''$, then $\Omega_5(h)=1$.

\begin{proof}
The algebraic map
\[
 C_5\longrightarrow\mathbb P^1_t,
 \qquad t=x^5=y^2+1
\]
is Galois of degree ten over $\mathbb Q(\zeta_5)$, with deck
group $C_5\times C_2$.  Its ramification orders over $0,1,\infty$
are $5,2,10$.  Removing these fibres gives the covering of
$\mathbb P^1\setminus\{0,1,\infty\}$ specified by $K_{10}$;
filling them gives exactly the normal closure of the displayed
peripheral powers.  This identifies $\Pi_5$.  Riemann--Hurwitz
gives
\[
 2g-2=10\left(-2+(1-1/5)+(1-1/2)+(1-1/10)\right)=2,
\]
so $g=2$.

The gentle action is peripheral: each of $x,y,z$ is sent to a
conjugate of itself at multiplier one.  It is the identity on
the abelian quotient $C_5\times C_2$.  Hence it preserves
$K_{10}$ and $N$, with the same true statement for its inverse;
the resulting map on $\Pi_5$ is an automorphism.  It commutes
with the order-five deck action on $\Pi_5^{\rm ab}$, since it
fixes $x$ literally.

The induced map on the four-dimensional homology is symplectic.
Here is a direct boundary argument that specifies its multiplier.
The open cover has eight punctures: two over zero, five over one,
and one over infinity.  The two over zero are fixed because $x$
is fixed.  The five over one are fixed individually because the
conjugator $f_h$ lies in $F'\subset K_{10}$, so it has trivial
deck coordinate.  The puncture over infinity is unique.
All eight peripheral loops are therefore preserved up to conjugacy,
with exponent one.  Fill the first seven punctures but leave
the one over infinity.  The resulting completed free group has
rank four; in a geometric free basis its remaining boundary is
\[
 [a_1,b_1][a_2,b_2].
\]
Our induced automorphism preserves its conjugacy class.  In
the quotient of the lower central series of degrees two and
three, conjugation disappears and this boundary gives
$a_1\wedge b_1+a_2\wedge b_2$.  Preservation of that element
is precisely the symplectic equation.  This proof works after
every finite reduction, hence over $\widehat{\mathbb Z}$.

The order-five rotation on $C_5$ has eigenvalues
$\zeta_5,\zeta_5^2$ on $dx/y,x\,dx/y$, and the conjugate
eigenvalues on the antiholomorphic part.  Its characteristic
polynomial on integral first homology is $\Phi_5$.
Consequently this lattice is a torsion-free rank-one module over
$\mathbb Z[\zeta_5]$, a Dedekind ring.  Such a module is an
invertible ideal and its endomorphism ring is exactly
$\mathbb Z[\zeta_5]$; after completion the same statement
holds prime by prime.  The symplectic adjoint of rotation is
its inverse, so the adjoint on that endomorphism ring is complex
conjugation.  The commuting symplectic automorphisms are exactly
the norm-one units in (D23.F3.1).  Restriction and abelianization
respect composition, establishing the asserted character and
its continuity.

For the arithmetic identification use the tangential point
at $t=0$ directed toward $1$, lifted to $x=0,y=i$ with leading
coefficient one in the local parameter $x=t^{1/5}$.
All its constant choices lie in $\mathbb Q^{\rm ab}$.
The standard Galois inclusion in $\widehat{GT}_0$ is the action
on this marked base fundamental group.  Functoriality of the
specified algebraic cover, followed by filling its punctures,
therefore identifies its restricted homology action with the
Tate action of $J_5$.  The complex/etale comparison used here is
the Riemann existence theorem; for an exact modern locator see
Milne, \emph{Lectures on Etale Cohomology}, v2.21, Theorem 3.4.
No claim about a five-point braid frame enters this identification.

Finally, $F'\subset K_{10}$ implies $F''\subset K_{10}'$.
If $f_h\in F''$, the action on $F/K_{10}'$ fixes both $x$ and
$y$ and hence is the identity.  It is therefore the identity on
$K_{10}^{\rm ab}$ and on its quotient $\Pi_5^{\rm ab}$.
\end{proof}

\subsection*{D23.F4. A scalar cyclotomic-unit character for the normalization}
Put $K_5^{\rm cyc}=\ker(F\to C_5)$, with $x\mapsto1$,
$y\mapsto0$.  Its standard Schreier generators are
\[
 X=x^5,\qquad Y_j=x^j y x^{-j}\quad(0\leq j<5).
\]
For $h\in\mathcal H$ define
\[
 \kappa_5(h)=\operatorname{coeff}_{Y_1}
       (f_h\bmod (K_5^{\rm cyc})')\pmod2.
 \tag{D23.F4.1}
\]
This is a continuous homomorphism to $\mathbb F_2$.
With the cyclic covering $t=x^5$ and the usual radial path
from zero to one, its restriction to $G_{\mathbb Q^{\rm ab}}$
is the Kummer character of $\zeta_5-1$.

\begin{proof}
Every $f_h$ belongs to $F'\subset K_5^{\rm cyc}$.  The action
fixes $X$ and conjugates each $Y_j$ by an element of
$K_5^{\rm cyc}$.  It is therefore the identity on the
abelianization of that group.  The composition rule
$f_{hk}=f_h\alpha_h(f_k)$ now makes every integral coefficient
in that abelianization additive; reduction modulo two proves
the first assertion.

For the arithmetic coefficient, pull back the double cover of
the $x$-line given by
\[
 v^2=1-x/\zeta_5.
\]
Its sign character on the cyclic covering fundamental group
is one on $Y_1$, zero on the other $Y_j$, and zero on $X$.
The radial path from $x=0$ to $x=1$ lifts from $v=1$ to a
chosen square root of $1-\zeta_5^{-1}$.  The usual Galois
path cocycle is the lift of $f_h$, up to inversion according
to the convention; inversion has no effect on its sign.
Thus (D23.F4.1) is the quadratic Kummer character of
$1-\zeta_5^{-1}$.  Its ratio with $\zeta_5-1$ is $\zeta_5^{-1}$,
which has a square root in $\mathbb Q^{\rm ab}$.  The two
quadratic characters on $G_{\mathbb Q^{\rm ab}}$ coincide.
\end{proof}

\subsection*{D23.F5. A direct common fifth-fibre bit is impossible}
Let $k=\mathbb Q^{\rm ab}$ and retain the actual torsion fields
$L_3,L_4,L_5$ of DELTA\_ON\_GT AF.4--AF.6.  No continuous
quadratic character of $\operatorname{Gal}(L_5/k)$ pulls back
on $G_k$ to the common arithmetic bit
\[
 a_2|_{G_k}=a_3|_{G_k}
        =\kappa_{\sqrt[4]{3}}.
 \tag{D23.F5.1}
\]
Thus a proposed proof identifying both bits with one quadratic
character of the genuine $J_5$ Tate representation is false
already on the Galois subgroup.

\begin{proof}
The field of the nonzero character on the right is
$k(\sqrt[4]3)=L_3\cap L_4$, by E.65 and AF.6.
If it factored through $\operatorname{Gal}(L_5/k)$, Galois
correspondence would give $k(\sqrt[4]3)\subset L_5$.
But AF.6 states $L_3\cap L_4\cap L_5=k$, a contradiction.
This applies to the full Tate image, so passing to a finite
level cannot repair the contradiction.
\end{proof}

\paragraph{Scope.}
D23.F3 and D23.F4 construct actual homomorphisms on the gentle
group, with arithmetic interpretations.  D23.F1--D23.F2
construct a different, frame-induced character on full GT.
An equality between the two constructions is not part of
these statements.  The explicit quadratic twist above shows
which scalar must be included in a comparison at the chosen
normalization; it does not make such a comparison on all GT
automatic.

% Arithmetic comparison. Independent whole-proof reading by /root completed;
% explicit positive-rotation marking added after that reader's requested check.
% No all-GT equality is asserted.

\subsection*{D23.F9. The exact scalar on the Galois subgroup}
Fix the usual ordered tangential braid base point
\[
 b_5=(0,t^4,t^3,t^2,t)
\]
and the monic hyperelliptic family
$Y^2=\prod_{j=1}^5(X-v_j)$ over the ordered configuration
space.  Choose the geometric path to the ordered CM
configuration $(1,\zeta_5,\zeta_5^2,\zeta_5^3,\zeta_5^4)$
that identifies its positive order-five rotation with the
standard class of $a=\sigma_1\sigma_2\sigma_3\sigma_4$.
The marking on the original character $\Omega_5$ is transported
along this path, with the order-five rotation represented by
$-q(a)$.  Then for $\sigma\in G_{\mathbb Q^{\rm ab}}$,
\[
 \mathcal C_5(\sigma)
       =(-1)^{\kappa_5(\sigma)}\Omega_5(\sigma).
 \tag{D23.F9.1}
\]
Consequently the character
\[
 \mathcal E_5(h)=(-1)^{\kappa_5(h)}
                  \mathcal C_5(h)\Omega_5(h)^{-1}
 \quad(h\in\mathcal H_5)
\]
is trivial on $G_{\mathbb Q^{\rm ab}}$.  This is a relative
character on the multiplier-one full GT group; no claim of its
triviality away from Galois or extension to the full gentle
group is made.

\begin{proof}
Let $M_\sigma$ be the arithmetic Tate action of the displayed
family at $b_5$. The exact base point and the arithmetic
Drinfeld action for every number of strands are recorded in
Nakamura, \emph{On arithmetic monodromy representations of
Eisenstein type in fundamental groups of once punctured
elliptic curves}, RIMS-1691 (2010), Section 5.4,
equation (5.4.2), printed p.~28. Arithmetic/topological
functoriality gives
\[
 q\Phi_\sigma=\operatorname{Ad}(M_\sigma)q.
\]
For multiplier one this is a symplectic action.  In addition
$M_\sigma e_1=e_1$ in the chain basis used by D23.T1.
To check the sign in this last assertion, let the first two
roots coalesce while the others remain at $t^3,t^2,t$.
The unit multiplying the double root at zero has leading
coefficient $-t^6$ and square root $it^3$.  Thus the node
splits over $\mathbb Q^{\rm ab}((t))$, with branches individually
fixed by $G_{\mathbb Q^{\rm ab}}$; its oriented vanishing
cycle is the first chain vector $e_1$.  The local annulus
has inertia $\widehat{\mathbb Z}(1)$, on which Galois acts
by the cyclotomic character, which is one here.  Hence the
cycle is fixed with its sign.  The uniqueness statement of
D23.T1 therefore gives $M_\sigma=V_\sigma$.

We specify the periodic-braid sign before comparing fibres.
Rotate the five branch points counterclockwise through
$2\pi/5$. Along the angular parameter $0\leq\theta\leq2\pi/5$
the isomorphisms from the initial curve to the rotating
family are
\[
 (X,Y)\longmapsto(e^{i\theta}X,e^{5i\theta/2}Y).
\]
Their endpoint is $(X,Y)\mapsto(\zeta_5X,-Y)$.
The positive periodic braid of this isotopy has the standard
class of $a$; choose the comparison path to identify these
positive classes. In the disc braid group both lifts have
exponent sum four, so a possible difference $z_5^e$ has
$20e=0$ and is trivial. The homological action of the
positive periodic braid is therefore exactly $q(a)$.
The order-five CM deck map $(X,Y)\mapsto(\zeta_5X,Y)$
differs from it by the hyperelliptic involution, whose
action on homology is $-I$. Its action is consequently
$-q(a)$, as required. This argument fixes the sign and
the direction of the cyclotomic marking, rather than only
the projective action.

Normalize the first two configuration coordinates to $0,1$.
As shown in D23.F1, its fundamental group is exactly
$\ker\ell_{12}\subset\widehat{PB}_5$.
At the base point the normalizing difference is $t^4$;
its fifth power has the fixed square root $t^{10}$, so this
normalization introduces no scalar on
$G_{\mathbb Q^{\rm ab}}$.

Let $p$ be the normalized CM point and transport its arithmetic
action to this normalized tangential base point along the fixed
path.  Its order-five automorphism is defined over
$\mathbb Q(\zeta_5)$ and is fixed by $\sigma$.
Consequently the path cocycle is a pure conjugator of $\bar a$.
By the pure centralizer calculation in D23.F1 it is precisely
$\alpha_5(\sigma)$.  Since the path lies in the normalized
configuration space, its lift in the pure braid group has
$\ell_{12}=0$ and is exactly $A_5(\sigma)$.
The standard base-point change formula for a local system now
identifies the Tate action at $p$, transported to the fixed
homology lattice, with
\[
 q(A_5(\sigma))^{-1}M_\sigma
       =\mathcal C_5(\sigma).
\]

Finally set $d=\zeta_5-1$.  The coordinate change
$X=1+dt$ in the original CM equation gives
\[
 X^5-1=d^5\prod_{j=0}^4
             \left(t-\frac{\zeta_5^j-1}{d}\right).
\]
The normalized monic family and the original curve are
therefore related by $Y=d^{5/2}Y_0$.
The only nontrivial scalar in their Tate actions is
$\sigma(\sqrt d)/\sqrt d$, because $d^2$ is already in
$\mathbb Q(\zeta_5)$.  D23.F4 identifies that scalar as
$(-1)^{\kappa_5(\sigma)}$.  The original Tate action is
$\Omega_5(\sigma)$ by D23.F3, giving (D23.F9.1).
All three factors are homomorphisms into the commutative
norm-one torus, so their product defines the stated relative
character.
\end{proof}

\paragraph{Scope of the comparison.}
The formula retains the specified tangential point, ordered
path, vanishing-cycle sign and base-point-change convention.
In particular the order-five deck rotation in $\Omega_5$
is identified with $-q(a)$. No equality between the three
characters on all of the full GT group is inferred from
their arithmetic equality.


\section{Level-six marking and exact comparison}
% DELTA23 /root/five_level_six. Exact marking and decision criterion.

\subsection*{D23.F6. An explicit marked Atkin--Lehner involution}
Use the literal modular matrices
\[
 s=\begin{pmatrix}0&1\\-1&0\end{pmatrix},\quad
 r=\begin{pmatrix}0&-1\\1&-1\end{pmatrix},\quad
 \tau=sr=\begin{pmatrix}1&-1\\0&1\end{pmatrix}
\]
in $\operatorname{PSL}_2(\mathbb Z)$, and put
$H=\overline{\Gamma_0(6)/\{\pm I\}}\subset\widehat B$.
Then $H$ is the free profinite group on
\[
 a=\tau,\qquad b=r\tau^5r,\qquad
 c=r\tau^2r\tau r\tau^2r.
 \tag{D23.F6.1}
\]
Their projective matrix representatives are respectively
\[
 A=\begin{pmatrix}1&-1\\0&1\end{pmatrix},\quad
 B=\begin{pmatrix}1&-1\\6&-5\end{pmatrix},\quad
 C=\begin{pmatrix}5&-3\\12&-7\end{pmatrix}.
\]
The Atkin--Lehner matrix
\[
 W_3=\begin{pmatrix}3&-2\\6&-3\end{pmatrix},\qquad
 \det W_3=3,\quad W_3^2=-3I
\]
induces the literal involution
\[
 w(a)=c^{-1},\qquad w(b)=b^{-1},\qquad w(c)=a^{-1}
 \tag{D23.F6.2}
\]
on this marked free group. Its outer class is the $w_3$ used
in DELTA\_ON\_GT FP.17.

\begin{proof}
Here is an explicit Bass--Serre/Schreier certificate for the
free-basis assertion.  The twelve right cosets of $H$ are
the projective bottom rows modulo two and three.  Numbering
them starting with $H$ as zero, the right actions of the
generators are
\[
 s=(0\ 1)(2\ 3)(4\ 5)(6\ 8)(7\ 9)(10\ 11),
\]
\[
 r=(0\ 2\ 1)(3\ 4\ 6)(5\ 7\ 10)(8\ 11\ 9).
\]
These permutations can be checked by multiplying rows by
the displayed two matrices in $\mathbb P^1(\mathbb F_2)$
and $\mathbb P^1(\mathbb F_3)$; their stabilizer at zero
is exactly the condition that the lower-left entry vanish
modulo six.  Use two consecutive edges in each of the four
displayed $r$-cycles, starting at $0,3,5,8$, and the three
$s$-edges $(2,3),(4,5),(6,8)$ as a tree.  After collapsing
that tree and the four $r^3$ relations, precisely the three
$s$-edges $(0,1),(7,9),(10,11)$ remain.  Their oriented
Schreier loops are exactly (D23.F6.1).  The $s^2$ relations
only reverse the chosen orientation of each such edge.
Thus these loops form a discrete free basis.  Completion of
this finite-index subgroup has the induced topology from
$B$, and gives the asserted profinite free basis.

Direct multiplication yields the three displayed matrices.
It also gives, in projective matrices,
\[
 W_3 A W_3^{-1}=C^{-1},\qquad
 W_3 B W_3^{-1}=B^{-1},\qquad
 W_3 C W_3^{-1}=A^{-1}.
\]
The free-basis result now proves (D23.F6.2) as full group
equalities, not only in a finite quotient.  In particular
$W_3$ normalizes the modular subgroup.  Its shape
$\left(\begin{smallmatrix}3u&v\\6w&3z\end{smallmatrix}\right)$
and determinant three give the usual level-six Atkin--Lehner
coset for the exact divisor three.  Replacing a representative
of that coset by a member of $H$ changes the induced map only
by an inner automorphism, so the outer class is the prescribed
one.  The supplied exact program \texttt{level6\_words.py}
independently performs the row enumeration, Schreier rewriting,
matrix conjugations, and square check.
\end{proof}

\subsection*{D23.F7. Two exact word tests for the level-six comparison}
For $h\in\mathcal H$ let $\Phi_h$ be the normalized modular
automorphism fixing $\tau$, and let $\rho_h$ be E.65's
integral harmonic coordinate.  Set
\[
 \psi_h=\operatorname{Ad}(a^{8\rho_h})\Phi_h\big|_H.
 \tag{D23.F7.1}
\]
This represents $\rho_6(h)$ of FP.17 and fixes $a$.
With $y=s\tau^{-2}s$ put
\[
 k_h=\psi_h(y^{-1})y\in H,\qquad
 d_h=\operatorname{coeff}_a(k_h^{\rm ab})-
                     \operatorname{coeff}_c(k_h^{\rm ab}),
 \qquad v_h=k_hc^{d_h}.
 \tag{D23.F7.2}
\]
In this formula $\psi_h$ on $y^{-1}$ means the same conjugated
automorphism of all of $\widehat B$ before restriction.
Then the exact criterion for A.86 is
\[
 AL_{6,3}(h)=1\quad\Longleftrightarrow\quad
 \begin{cases}
 v_h w(v_h)=1,\\
 \psi_h(b)^{-1}=v_h w(\psi_h(b))v_h^{-1}.
 \end{cases}
 \tag{D23.F7.3}
\]
The balanced frame $v_h$ does not depend on which conjugator
$k_h\in H$ is used for $\psi_h(c)=k_hck_h^{-1}$.
There is no remaining unspecified inner conjugator in this test.

\begin{proof}
E.65.9--10 proves
$q_{12}\Phi_h=\operatorname{Ad}(U^{8\rho_h})q_{12}$
on the actual modular quotient, with $U=q(\tau^{-1})$.
Thus (D23.F7.1), defined first on $\widehat B$, is the identity
on its level-six quotient and preserves $H$.  It fixes $a$
since both factors do.  It represents the restriction outer
action of FP.17 because its correction is inner in
$\widehat B$.

In the literal matrices one has
\[
 g=y^{-1}=\begin{pmatrix}1&0\\2&1\end{pmatrix},
 \qquad c=ga^{-3}g^{-1}.
\]
Consequently $k_h=\psi_h(g)g^{-1}$ conjugates $c$ to
$\psi_h(c)$.  Because $\psi_h$ is the identity modulo six,
$k_h$ belongs to the principal level-six subgroup, hence to
$H$.  In terms of the original parameter this is the explicit word
\[
 k_h=a^{8\rho_h}f_h^{-1}y^{-1}f_h a^{-8\rho_h}y,
 \quad
 \psi_h(b)=a^{8\rho_h}f_h^{-1}r a^5 f_h^{-1}r a^{-8\rho_h}.
 \tag{D23.F7.4}
\]

For completeness, the following elementary free-group argument
eliminates the inner ambiguity.  Write $\psi=\psi_h$, $k=k_h$.
The condition $[\psi,w]=1$ in $\operatorname{Out}(H)$ says
that for some $v\in H$,
$\psi w=\operatorname{Ad}(v)w\psi$.
Apply this equality to $a$.  It gives
$k c^{-1}k^{-1}=v c^{-1}v^{-1}$.  Since a free generator
has centralizer its procyclic subgroup, $v=kc^e$ for a
unique $e\in\widehat{\mathbb Z}$.
Applying it to $c$ gives
\[
 k c^e w(k)\in\overline{\langle a\rangle}.
\]
If $k^{\rm ab}=p a+q b+t c$, then (D23.F6.2) shows that
the abelianization on the left is
$(p-t)a+(t+e-p)c$.
Thus necessarily $e=p-t=d_h$, and the displayed element
must be exactly $a^{d_h}$.  This is equivalent to
$v_hw(v_h)=1$.  The equality on the final free generator
$b$ is exactly the second line of (D23.F7.3).
Conversely those two equalities give the required identity
on all three free generators and hence the equality of
continuous automorphisms.  This proves the criterion.

Changing $k$ to $kc^j$ changes $d_h$ to $d_h-j$, so
$kc^j c^{d_h-j}=kc^{d_h}$.  This proves the asserted
independence.  The centralizer fact used here is the
free-factor centralizer theorem for
$H=\overline{\langle c\rangle}\amalg\widehat F_2$
(and similarly for $a$); it is the same Herfort--Ribes
Theorem B' used in D23.F1.
\end{proof}

\paragraph{What the criterion does and does not decide.}
For the element $h_*$ of E.58 one has $\rho_{h_*}=0$, so
\[
 k_{h_*}=f_{h_*}^{-1}y^{-1}f_{h_*}y,
 \qquad \psi_{h_*}(b)=f_{h_*}^{-1}r a^5 f_{h_*}^{-1}r.
\]
E.58 specifies a relative-pronilpotent lift and finite frame
coordinates, rather than these two full free-group values.
No equality or inequality in (D23.F7.3) on that element is
claimed here.  The computation supplies a marked basis and
an exact normalization for the test that was absent in FP.17;
the asserted equivalence with $\Delta=0$ remains a separate
decision.

\subsection*{D23.F8. The Atkin--Lehner matrices do not generate the division algebra of discriminant six}
An algebra generated over $\mathbb Q$ by rational
$2\times2$ Atkin--Lehner matrices is a subalgebra of
$M_2(\mathbb Q)$.  It cannot be isomorphic to the division
quaternion algebra $(2,3)$ of discriminant six.  In particular
the literal matrices used to mark A.86 do not provide an
identification with that quaternion algebra.

\begin{proof}
The Hilbert-symbol calculation $(2,3)_2=(2,3)_3=-1$,
with all other local symbols positive, is the standard
local calculation already made in D23.L3.  It is therefore
a division quaternion algebra over $\mathbb Q$.
An injective homomorphism from that four-dimensional algebra
to $M_2(\mathbb Q)$ would be an isomorphism by dimension.
But $M_2(\mathbb Q)$ has nonzero noninvertible elements and
a division algebra does not, a contradiction.  Any nonzero
unital homomorphism from a division algebra is injective,
so even that weaker possibility is excluded.
This obstruction concerns a literal matrix-algebra
identification; it does not exclude a separately constructed
Brauer class or a Shimura-curve comparison.
\end{proof}

\subsection*{D23.F10. The extra-CM element passes every pro-two level-six test}
Let $h_*$ be the actual element of LM7\_PUSH E.58, and let
$H^{(2)}$ be the maximal pro-two quotient of the marked
level-six subgroup $H$.  The automorphism $\rho_6(h_*)$
is the identity on $H^{(2)}$. Consequently each of
$AL_{6,2}(h_*),AL_{6,3}(h_*),AL_{6,6}(h_*)$ is the identity
on that entire quotient, although $\Delta(h_*)=1$.
No test factoring through the induced action on
$H^{(2)}$ can establish $AL_{6,3}=1\Longleftrightarrow\Delta=0$.

\begin{proof}
Put $U=F'F^3$, as in E.58.  That theorem supplies
$\Phi_{h_*}=1$ on $\widehat B/R_{\rm nil}(U)$ and
$\rho_{h_*}=0$.  The latter makes the correction in
(D23.F7.1) trivial.
The finite source calculation E.58.11--13, reduced modulo
six, gives
\[
 q_6(U)=\{1\}\times V_4
 \quad\hbox{in }S_3\times A_4.
\]
On the other hand $q_6(H)=C_2\times C_3$ by FP.17.
Their intersection is trivial. Thus
$V=U\cap H=\ker(q_6|_U)$ is normal of index four in $U$,
with quotient $V_4$.

For a normal subgroup $V$ of a profinite group $U$ with
finite two-group quotient, one has $O^2(V)=O^2(U)$.
Indeed, if $N$ is an open normal subgroup of $V$ of
two-power index, intersect its finitely many $U$-conjugates.
The resulting $N_0$ is normal in $U$; the quotient
$V/N_0$ embeds into a finite product of two-groups, and
$U/N_0$ is an extension of that two-group by $U/V$.
Hence $O^2(U)\subset N_0\subset N$. Intersecting over $N$
gives $O^2(U)\subset O^2(V)$.  Conversely restriction of
every finite two-group quotient of $U$ to $V$ gives
$O^2(V)\subset O^2(U)$.
Since every finite nilpotent residual lies in every
pro-prime residual, and every homomorphism from $H$ to a
finite two-group restricts to such a homomorphism on $V$,
we obtain
\[
 R_{\rm nil}(U)\subset O^2(U)=O^2(V)\subset O^2(H).
\]
For each $h\in H$ the difference
$\Phi_{h_*}(h)h^{-1}$ belongs to $R_{\rm nil}(U)$,
so its image in $H^{(2)}$ is one.  Thus the induced
automorphism is the identity.  The characteristic subgroup
$O^2(H)$ is preserved by each $w_e$, so their crossed
commutators with this automorphism are identity too.
The sign $\Delta(h_*)=1$ is E.58.23 and E.65.4.
\end{proof}

\subsection*{D23.F12. One actual mixed-sign element is invisible to every pronilpotent level-six test}
Assume the supplied relative peripheral lifting theorem
LM7\_PUSH E.58, in its stated version for every open normal
subgroup.  There is an actual element $h_6\in\mathcal H$
with
\[
 \Delta(h_6)=1,\qquad\rho_{h_6}=0,
 \qquad \rho_6(h_6)|_{H/R_{\rm nil}(H)}=1.
 \tag{D23.F12.1}
\]
Consequently all three level-six Atkin--Lehner defects of
this one element are the identity on every finite nilpotent
quotient of $H$, and on its entire maximal pronilpotent
quotient.  This is an actual gentle element, rather than
a collection of unrelated finite shadows.

\begin{proof}
Let
\[
 U=N_6=\ker(\widehat B\longrightarrow
              Q_6=\operatorname{SL}_2(\mathbb Z/6)/\{\pm I\}).
\]
Then $U$ is open normal and $U\subset H\cap F$.
Let $q=q_{12}$ be the exact projective level-twelve quotient
of E.58. Its image $q(U)=\ker(Q_{12}\to Q_6)$ is an
elementary abelian two-group of order eight.  Indeed every
class in this kernel has a unique representative congruent
to $I$ modulo six, of the form $I+6M$ with
$M\in M_2(\mathbb F_2)$ and $\operatorname{tr}M=0$.
Multiplication is addition of $M$, proving the assertion.
The central element $z_0=[5I]=[7I]$ belongs to this kernel
and corresponds to $M=I$.

Thus $q$ factors through
$\Gamma_2=\widehat B/O^2(U)$, and the kernel of
$\Gamma_2\twoheadrightarrow Q_{12}$ is pro-two.
The element $r$ has order three in both groups.
The coprime centralizer-lifting argument in the proof of
E.58.19 gives a lift
\[
 c_2\in C_{\Gamma_2}(r)\cap U^{(2)},
 \qquad q(c_2)=z_0.
\]
For specificity, that argument lifts the centralizer
element $z_0$, then conjugates its lifted conjugate of
$\langle r\rangle$ back by the pro-two kernel using
finite Schur--Zassenhaus complement conjugacy and compactness.
Since $z_0$ is identity in the retained base $Q_6$, the
corrected lift lies in $U^{(2)}$.

Use E.58.2 to identify the relative pronilpotent quotient
with the fibre product of the relative pro-prime quotients
over $Q_6$.  Give its identity automorphism the specified
peripheral conjugators
\[
 b_s=1,\qquad b_\tau=1,\qquad
 b_r=(c_2\text{ at two},\ 1\text{ at every other prime}).
\]
Their base coordinates agree, and $b_r$ really centralizes
$r$ in that quotient.  Apply the supplied relative lifting
theorem, and normalize the full lift to fix $\tau$ exactly.
This produces an actual automorphism $\Phi$ of
$\widehat B$, identity on
$\widehat B/R_{\rm nil}(U)$, with unique pure frames
\[
 \beta\in R_{\rm nil}(U),\qquad A\in U,
 \qquad q(\beta)=1,\quad q(A)=z_0.
\]
The literal parameter $f=s\Phi(s)$ lies in
$R_{\rm nil}(U)\subset U'\subset F'$.
The torsion and product calculation E.58.21, applied to
this normalized lift, gives the two gentle equations and
the actual free-group automorphism.  Hence this is an
element $h_6\in\mathcal H$.  Because $\beta\in U'\subset F'$,
its harmonic coordinate is zero.  The same signed
calculation as E.58.23 gives $(a_3,a_2)=(1,0)$ and
$\Delta(h_6)=1$.

Every continuous map from $H$ to a finite nilpotent group
restricts to such a map on $U$, since $U\subset H$.
Therefore
\[
 R_{\rm nil}(U)\subset R_{\rm nil}(H).
\]
The identity of $\Phi$ modulo the former residual proves
the last assertion in (D23.F12.1).  Each Atkin--Lehner
automorphism preserves the characteristic subgroup
$R_{\rm nil}(H)$, so its crossed commutator with this
identity action is trivial on the whole quotient.
All finite nilpotent tests follow at once from this one
full quotient identity.
\end{proof}

\paragraph{The remaining level-six gate after D23.F12.}
To distinguish this witness from the kernel of $\Delta$
using $AL_{6,3}$, a test must retain a genuinely
nonnilpotent quotient of the free level-six group.
No claim that $AL_{6,3}(h_6)$ itself is trivial is made.

\subsection{Every congruence Atkin--Lehner test is trivial}

\paragraph{D23.F18.}
Let $H$ be the marked profinite level-six subgroup of
D23.F6 and let $H^{\rm cong}$ be its full congruence
completion, namely the closure of
$\Gamma_0(6)/\{\pm I\}$ in
$\operatorname{SL}_2(\widehat{\mathbf Z})/\{\pm I\}$.
For every multiplier-one gentle element $h$, its
induced outer action on $H^{\rm cong}$ is trivial.
The three Atkin--Lehner involutions extend to this
completion. Consequently every outer defect
$AL_{6,e}(h)$, $e=2,3,6$, is trivial on the entire
congruence completion and on every finite quotient
of it to which the Atkin--Lehner action descends.

In particular the congruence version of
$AL_{6,3}=1\Longleftrightarrow\Delta=0$ is false:
the supplied E.58 witness has $\Delta=1$ and still
has trivial congruence defect. Any decision of the
full A.86 comparison must use noncongruence
information, not just a finer congruence level.

\begin{proof}
Write $q$ for the adelic congruence representation,
projectivized by the global sign when its source
is $\widehat B$.  The signed covariance theorem
DELTA\_ON\_GT AD.1 gives a uniquely defined
$t_h\in\widehat{\mathbf Z}$ with
\[
 q\Phi_h=\operatorname{Ad}(U^{t_h})q,
 \qquad U=\begin{pmatrix}1&1\\0&1\end{pmatrix}.
\]
The marked cusp generator of $H$ is $a=\tau$,
with $q(a)=U^{-1}$.  Thus $U^t$ belongs to
$H^{\rm cong}$ for every $t\in\widehat{\mathbf Z}$.
This can also be seen directly from its lower-left
entry, which is zero.

For the representative used in D23.F7,
$\psi_h=\operatorname{Inn}(a^{8\rho_h})\Phi_h$,
the congruence covariance is
\[
 q\psi_h
  =\operatorname{Ad}(U^{t_h-8\rho_h})q.
\tag{D23.F18.1}
\]
The conjugating element is itself in $H^{\rm cong}$.
Hence the induced automorphism of that completion
is inner. In particular the congruence kernel of
$H$ is preserved: equation (D23.F18.1) proves this,
so no kernel invariance is assumed in forming
the quotient action.

For precision about bad primes, the closure in
question is exactly the projective image of
\[
 \{g\in\operatorname{SL}_2(\widehat{\mathbf Z}):
                         g_{21}\in6\widehat{\mathbf Z}\}.
\]
Indeed reductions of $\operatorname{SL}_2(\mathbf Z)$
surject onto every finite congruence quotient,
and lifting at a modulus divisible by six retains
the displayed lower-left condition. Thus it
is the inverse image of the usual level-six
subgroup and not an unspecified larger adelic
group.

Every rational Atkin--Lehner matrix normalizing
$\Gamma_0(6)$ acts continuously by conjugation
on its local closures, with a continuous inverse,
and therefore on the full congruence completion.
For the specified $W_3$ this is completely
explicit. If
\[
 W_3=\begin{pmatrix}3&-2\\6&-3\end{pmatrix},
 \qquad g=\begin{pmatrix}a&b\\6c&d\end{pmatrix},
\]
then
\[
 W_3gW_3^{-1}=
 \begin{pmatrix}
 -3a+12c-6b+4d&2a-8c+3b-2d\\
 6(-a+3c-2b+d)&4a-12c+6b-3d
 \end{pmatrix}.
\]
These integral formulas make sense over
$\widehat{\mathbf Z}$, preserve determinant one
and the level-six condition, and their square
is the identity because $W_3^2=-3I$.
The same closure argument applies to $W_2,W_6$;
their classes square to inner automorphisms as
usual for Atkin--Lehner involutions.

In the outer automorphism group of
$H^{\rm cong}$ the class of (D23.F18.1) is
therefore identity. Its crossed commutator with
any Atkin--Lehner class is identity as well.
If a chosen finite normal congruence kernel is
not initially Atkin--Lehner stable, intersect
its finitely many translates under the
Atkin--Lehner group; their squares act inner
and hence fix every normal subgroup. The
result is a finite-index invariant congruence
kernel, and the same inner-action argument
applies there. This explains the precise
finite-quotient formulation of the statement.

Finally E.58.23 and E.65 supply an actual gentle
element with $\Delta=1$. The just-proved
congruence conclusion applies to it, proving
the stated refutation of the congruence version.
No conclusion that its full noncongruence
Atkin--Lehner defect is trivial is drawn.
\end{proof}


\section{Relative characters and the full adelic image}
% DELTA23 characters workstream. Author: GPT-6 Astra Pro (/root/characters).
% First proof reading completed on 2026-11-21 Europe/Berlin.
% This file uses the input notation G=GT_0-hat, H=ker(lambda).

\subsection{A full relative character from the harmonic descent defect}

\paragraph{D23.C1 (proved from the cited inputs).}
Let $G=\widehat{\mathrm{GT}}_0$, $H=\ker\lambda$, and let
$D=D_P:G\to G$ be the endomorphism of
DELTA\_ON\_GT A.14, in its literal current marking.  Write
$a_2=\epsilon_2W$ and $a_3=\epsilon_3T$ on $H$.
The character
\[
 d_H(h)=a_2(h)+a_2(Dh)\quad(h\in H)
\]
has a unique continuous extension
\[
 d:G\longrightarrow\mathbb F_2,
 \qquad d|_{G_{\mathbb Q}}=0.
\]
For $g\in G$, its choice-independent formula is
\[
 d(g)=d_H(g\sigma^{-1}),\qquad
 \sigma\in G_{\mathbb Q},\quad\chi(\sigma)=\lambda(g).
\]
On $H$, this is the mixed duplication defect $d_2$ of
DELTA\_ON\_GT A.103 and A.107.  Its nonvanishing is not asserted.

\emph{Proof.}
The two individual characters $a_2,a_3$ are invariant under
conjugation by all of $G$: this is the separate normalizer
calculation in LM7\_PUSH E.65, equations (E.65.16)--(E.65.17),
not just invariance of their sum.  Since $D$ is a homomorphism
preserving $\lambda$, it maps $H$ to $H$, and
\[
 a_2(D(ghg^{-1}))=a_2(Dg\,Dh\,(Dg)^{-1})=a_2(Dh).
\]
Thus $d_H$ is a continuous, $G$-invariant character of $H$.
The arithmetic marking comparison DELTA\_ON\_GT A.62 says
$D(\tau)=c\tau c$ for $\tau\in G_{\mathbb Q}$, where
$c=(-1,1)$.  If also $\tau\in H$, conjugation invariance gives
$a_2(D\tau)=a_2(\tau)$; hence $d_H(\tau)=0$.

Every element of $G$ has the form $h\sigma$ with $h\in H$ and
$\sigma\in G_{\mathbb Q}$, because the cyclotomic character is
surjective.  If $h\sigma=h'\sigma'$, then
$h'^{-1}h=\sigma'\sigma^{-1}\in H\cap G_{\mathbb Q}$,
on which $d_H$ is zero.  Therefore $d(h\sigma)=d_H(h)$ is
well defined.  Multiplication in this form gives
\[
 (h\sigma)(k\tau)=h(\sigma k\sigma^{-1})(\sigma\tau),
\]
so conjugation invariance proves additivity.  The map
$H\times G_{\mathbb Q}\to G$, $(h,\sigma)\mapsto h\sigma$,
is a continuous surjection from a compact space to a Hausdorff
space.  The continuous function $(h,\sigma)\mapsto d_H(h)$ is
constant on its fibres, so the resulting $d$ is continuous.
The same decomposition proves uniqueness.  A.107, equation
(MD.8b), identifies its restriction with $d_2$.  No lifting
argument or pentagon relation enters this proof.\hfill$\square$

\paragraph{D23.C2 (proved; exact restricted classification).}
For an integer $n\geq0$, define the finite parity map
\[
 p_n:H\longrightarrow(\mathbb F_2^2)^{n+1},\qquad
 p_n(h)=\bigl((a_2(D^jh),a_3(D^jh))\bigr)_{0\leq j\leq n}.
\]
Among the continuous characters $G\to\mathbb F_2$ trivial on
$G_{\mathbb Q}$ whose restriction to $H$ factors through $p_n$,
all are precisely the span of
\[
 \Delta\circ D^j\quad(0\leq j\leq n),\qquad
 d\circ D^j\quad(0\leq j<n).
\]
This gives generators, not a claim that they are independent.
In particular, allowing finite strings of descendant frame parities
gives the $\mathbb F_2[D^*]$-module generated by $\Delta$ and $d$,
where $D^*\eta=\eta\circ D$.  It does not classify characters
without this factorization hypothesis.

\emph{Proof.}
Let $S_n=p_n(H)$, an $\mathbb F_2$-subspace of the displayed
finite vector space.  By DELTA\_ON\_GT B.14 the pair $(a_2,a_3)$
on $G_{\mathbb Q}\cap H$ has image $\{(0,0),(1,1)\}$.
The arithmetic formula $D\tau=c\tau c$, together with separate
conjugation invariance of the two parities, shows that the image
of $G_{\mathbb Q}\cap H$ under $p_n$ is exactly the diagonal line
spanned by the vector with all $2n+2$ entries equal to one.

A character factoring through $p_n$ is a linear functional on
$S_n$.  Extend it to the ambient vector space by extending a
basis of $S_n$ to a basis of that space.  Since the diagonal
vector lies in $S_n$ and its functional value is zero, the
extended functional has an even sum of coordinate coefficients.
The even-sum coefficient vectors are spanned by the following
$2n+1$ edge vectors: the sums of the two coordinates at a fixed
level $j$, and the sums of the harmonic coordinates at consecutive
levels $j,j+1$.  To check this without any rank assumption on
$S_n$, view the $2n+2$ coordinate positions as vertices: these
edges form a tree.  Removing a leaf and its incident edge eliminates
one specified coefficient; induction leaves the zero coefficient
at the final vertex because the total sum is even.  The corresponding
functions on $H$ are respectively
\[
 (a_2+a_3)\circ D^j=\Delta|_H\circ D^j,
 \qquad (a_2+a_2\circ D)\circ D^j=d_H\circ D^j.
\]
They extend to the global relative characters in the statement by
E.65 and D23.C1.  Two global characters trivial on $G_{\mathbb Q}$
that agree on $H$ agree on $G=HG_{\mathbb Q}$.  This proves both
containments and the asserted exact restricted classification.
\hfill$\square$

\paragraph{D23.C3 (proved; a separation and the harmonic lifting obstruction).}
Assume the actual element $k$ of DELTA\_ON\_GT LI.8
(integrated statement A.94) with
$Dk=1$, $\rho_k=0$, $q_{12}(\beta_k)=1$, and $\Delta(k)=1$.
Then
\[
 \Delta(k)=1,\qquad
 (\Delta\circ D^j)(k)=0\ (j\geq1),\qquad
 (d\circ D^j)(k)=0\ (j\geq0).
\]
In particular, if $d$ is nonzero it is independent of $\Delta$.
More generally $\Delta$ is outside the span of every positive
iterate of $\Delta$ and every iterate of $d$.

The elementary coprime-centralizer construction used in E.58
cannot, with identity relative action and a seed lying in the
relative kernel, be turned into a witness with
$q_{12}(\beta)=W=\pm S$ at three by replacing its prime two
component with a prime three component.

\emph{Proof.}
The retained frame and zero shear in LI.8 give $a_2(k)=0$.
Since $Dk=1$, the definition in D23.C1 yields $d(k)=0$ and
all positive iterates vanish.  Evaluating any proposed linear
relation at $k$ proves the separation claims.

For the last assertion, the relative lift in E.58 first chooses an
open normal subgroup $U$ and uses a seed whose finite frame value
belongs to $q(U)$, with $q(U)$ a $p$-group.  A harmonic value
$\pm S\in\mathrm{SL}_2(\mathbb F_3)$ has order four, so such
a $q(U)$ must be a two-group: it cannot be a three-group or any
odd-primary group.  But the torsion element being centralized is
$s$, of order two.  The hypothesis that its order be prime to the
pro-$p$ kernel, required for the coprime complement-conjugacy
argument, fails at this forced prime $p=2$.  Therefore that specific
construction does not establish the required centralizer lift.
This proves an obstruction to the proposed proof, not a theorem
that no other harmonic lift exists.\hfill$\square$

% DELTA23 characters workstream, 2026-11-21 Europe/Berlin.
% D23.C4 is an arithmetic derivation. D23.C5 was expanded after a root
% review to separate ordinary Legendre paths from the projective orbifold.
% Root independently read and passed C4--C7 on 2026-11-21, 00:29 CET;
% see the root check register for its exact event time.

\subsection{The arithmetic size of the normalized adelic frame target}

Put $O_3=\mathbb Z[\zeta_3]$, $O_4=\mathbb Z[i]$,
$K_j=O_j\otimes\mathbb Q$, and write a bar for the nontrivial
automorphism of $K_j$.  Set
\[
 N_j=\{u\in(O_j\otimes\widehat{\mathbb Z})^\times:u\bar u=1\},
 \qquad N_j(4)=\{u\in N_j:u\equiv1\pmod4\}.
\]
Here the congruence modulo four concerns the component over two.
The integral matrix embeddings are the ones in DELTA\_ON\_GT
AD.2 (integrated A.72), so $N_3(4)$ and $N_4(4)$ are precisely
its two normalized targets.

\paragraph{D23.C4 (arithmetic derivation; first proof reading passed).}
For $E_3:y^2=x^3+1$ and $E_4:y^2=x^3-x$, choose CM-compatible
bases of their full Tate modules.  Let $\rho_j$ be their adelic
representations and set $G^{\mathrm{cyc}}=G_{\mathbb Q^{\rm ab}}$.
Then
\[
 \rho_4(G^{\mathrm{cyc}})=N_4(4).
\]
For $r_0=(3/4)^{1/4}$ and
$\kappa(\sigma)=\sigma(r_0)/r_0\in\{\pm1\}$ on
$G^{\mathrm{cyc}}$, the normalized representation
\[
 \rho_3^\circ(\sigma)=\kappa(\sigma)\rho_3(\sigma)
\]
has image $N_3(4)$ and has the same kernel as $\rho_3$ on
$G^{\mathrm{cyc}}$.  Both full division fields are
$\mathbb Q(E_j[\mathrm{tors}])=K_j^{\rm ab}$.

\emph{Proof.}
We use the main theorem of complex multiplication in the torsion-point
form: Milne, \emph{Modular Functions and Modular Forms}, version 1.10,
Theorem 12.27, printed p.124
(\url{https://www.jmilne.org/math/CourseNotes/MF110.pdf}).
For a class-number-one imaginary quadratic ring $O$, it states that
the image of the full torsion representation after quotient by the
geometric automorphism group $O^\times$ is all of
$\widehat O^\times/O^\times$, with the usual ray-class interpretation.
It also implies that adjoining the Weber values of all torsion points
generates $K^{\rm ab}$.  All torsion coordinates are abelian over $K$
because the CM action is defined there and its Tate-module centralizer
is the commutative group $\widehat O^\times$.  Hence the full division
field over $K$ is $K^{\rm ab}$.

Let $\pi_4=1+i$ and $\pi_3=\sqrt{-3}$.  The entire subgroup
$E_4[\pi_4^3]$ is $K_4$-rational: besides the four two-torsion
points, the four points with abscissae $i,-i$ have ordinates
$\pm(1-i),\pm(1+i)$, respectively, and are the remaining points
of this kernel.  This last kernel identification can be checked
using the automorphism $[i](x,y)=(-x,iy)$ and the chord-and-tangent
addition law: applying $1+i$ to any of those four points gives
one of the nonzero two-torsion points in $E_4[\pi_4^2]$ which
maps to $(0,0)$ under $1+i$.
Thus $\rho_4(G_{K_4})$ is contained in
\[
 C_4=\{u\in\widehat O_4^\times:u\equiv1\pmod{\pi_4^3}\}.
\]
The reduction $\mu_4\to(O_4/\pi_4^3)^\times$ is an isomorphism:
the target has four elements, and the four roots of unity are distinct
modulo $\pi_4^3$.  Consequently $C_4\to\widehat O_4^\times/\mu_4$
is an isomorphism.  The main theorem therefore gives equality
$\rho_4(G_{K_4})=C_4$.

For $E_3$, all two-torsion points have abscissae
$-1,-\zeta_3,-\zeta_3^2$, and the kernel of $1-\zeta_3$ is
$\{O,(0,1),(0,-1)\}$.  The coprime kernels show that the whole
$E_3[2\pi_3]$ is $K_3$-rational.  The reduction
\[
 \mu_6\longrightarrow (O_3/(2\pi_3))^\times
       \simeq\mathbb F_4^\times\times\mathbb F_3^\times
\]
is an isomorphism: the cubic roots give the first factor and the
sign gives the second.  The same argument proves
\[
 \rho_3(G_{K_3})=C_3
   =\{u\in\widehat O_3^\times:u\equiv1\pmod{2\pi_3}\}.
\]
The CM fields themselves occur in the torsion fields (use the
three roots of the two-division polynomial for $E_3$, and the
Weil pairing on four-torsion for $E_4$), proving the asserted
identities over $\mathbb Q$ as well.

Complex conjugation acts on each Cartan by $u\mapsto\bar u$.
Since $C_j$ is abelian and the quotient of the full Galois image
by $C_j$ has order two, its closed derived subgroup is
$\{u/\bar u:u\in C_j\}$.  This image is compact, so no further
closure changes it.  Kronecker--Weber identifies
$G^{\mathrm{cyc}}$ with the closed derived subgroup of
$G_{\mathbb Q}$; taking the image of a closed derived subgroup
under a compact surjection gives the closed derived subgroup of
the target.  Therefore
\[
 \rho_j(G^{\mathrm{cyc}})=\{u/\bar u:u\in C_j\}.
\]

At every unramified finite place, the map $u\mapsto u/\bar u$
from local units is onto the norm-one local units.  For a split
place this is the explicit map $(a,b)\mapsto(a/b,b/a)$.
For an unramified quadratic local field extension, Hilbert 90
provides a preimage, and multiplying by a base-field uniformizer
power makes it a unit because the valuations have the same value
group.  At the inert prime two of $K_3$, a norm-one unit equal to
one modulo two has a Hilbert-90 unit preimage whose residue is
fixed by the residue-field involution.  That residue is in
$\mathbb F_2^\times=\{1\}$, so the preimage already lies in
$1+2O_{3,2}$.  Conversely such a preimage has a ratio equal to
one modulo two.

At the ramified prime three of $K_3$, a norm-one unit $v$ has
residue $\pm1$ modulo $\pi_3$.  If $v\equiv1$, a Hilbert-90
preimage has even $\pi_3$-valuation: a preimage of odd valuation
has ratio with residue $\pi_3/\bar\pi_3=-1$.  Multiplying by
a rational power of three makes the preimage a unit; multiplying
by $1$ or $-1$ makes its residue one.  Conversely a unit preimage
has ratio congruent to one.  Hence
\[
 \rho_3(G^{\mathrm{cyc}})
 =\{u\in N_3:u_2\equiv1\pmod2,
                     \ u_3\equiv1\pmod{\pi_3}\}=:M_3.
\]

At the ramified prime two of $K_4$, a unit
$u=a+bi\equiv1\pmod{\pi_4^3}$ has even $b$, and therefore
$u/\bar u\equiv1\pmod4$.  Conversely let $v\in N_{4,2}$
be one modulo four.  Put $b_0=(1+v)/2$.  It is a unit, is one
modulo two, and $b_0/\bar b_0=v$.  Among the four residue classes
of units modulo $\pi_4^3$, the two equal to one modulo two are
$1$ and $-1$.  Thus one of $b_0,-b_0$ belongs to
$1+\pi_4^3O_{4,2}$ and has the same ratio.  This proves
$\rho_4(G^{\mathrm{cyc}})=N_4(4)$.

The norm-one subgroup of $(O_3/4)^\times$ which is one modulo
two is $\{1,-1\}$.  For example write a representative as
$1+2a$; its norm is one modulo four exactly when the trace of
$a\in\mathbb F_4$ is zero, i.e. when $a=0$ or $1$.
By DELTA\_ON\_GT A.13, or its explicit four-division formulas,
the sign of $\rho_3(\sigma)$ modulo four is exactly
$\kappa(\sigma)$.  Thus $\rho_3^\circ$ takes values in $N_3(4)$.

The map $M_3\to N_3(4)$ which multiplies by this modulo-four
sign is an isomorphism.  Its inverse sends $v\in N_3(4)$ to
$\varepsilon(v_3)v$, where $\varepsilon(v_3)\in\{\pm1\}$ is
the residue of its three-component modulo $\pi_3$.  Indeed this
product has three-component equal to one modulo $\pi_3$, and
two-component equal to that sign modulo four.  Both compositions
are identity.  This proves surjectivity and the assertion about
the kernel of the normalized representation.\hfill$\square$

\paragraph{D23.C5 (proved; uniform signed integral path comparison).}
In the positive normalization of LM7\_PUSH E.65 and
DELTA\_ON\_GT AD.2, for every
$\sigma\in G_{\mathbb Q^{\rm ab}}$ one has, in CM-compatible
transported bases,
\[
 \mathcal W_\sigma=\rho_4(\sigma),\qquad
 \mathcal T_\sigma=\kappa(\sigma)\rho_3(\sigma).
\]
These are full adelic identities, not merely identities modulo twelve.

\emph{Proof, with the path marking and signed normalization made explicit.}
Use the Legendre family
$\mathscr E_t:Y^2=X(X-1)(X-t)$ on
$Y=\mathbb P^1\setminus\{0,1,\infty\}$, with the fixed positive
peripheral generators and the rational tangent $\overrightarrow{01}$.
The embedded Galois action on $F$ is this normalized tangential
action.  Its geometric monodromy in a symplectic
basis is $x\mapsto U^2$, $y\mapsto L^2$: a positive loop about
either collision of two branch points is the square of the half
twist, whose lift to the double cover is the corresponding Dehn
twist, and the two vanishing cycles have intersection one.  This
gives the signed free-subgroup representation used in AD.2.

At the tangential cusp over $t=0$, the Tate degeneration has a
primitive rank-one vanishing submodule and a rank-one quotient.
The special cubic is $Y^2=X^2(X-1)$; its two branches at the
node have slopes $\pm i$.  Thus its only splitting character is
the quadratic character of $\mathbb Q(i)$, and this is trivial on
$G_{\mathbb Q^{\rm ab}}$.  The vanishing line has cyclotomic
action and the quotient has trivial action.  Both therefore have
action one on that subgroup, including integrally at two.
The actual arithmetic matrix $V^{\rm geom}_\sigma$ at the cusp
is consequently upper unipotent with lower-right entry one.
Functoriality gives
$q\Phi_\sigma=\operatorname{Ad}(V^{\rm geom}_\sigma)q$ on $F$.
This already identifies it with $\mathcal V_\sigma$: an upper
unipotent matrix commuting with $U^2,L^2$ is the identity over
every $\mathbb Z_\ell$, including $\ell=2$, because those rings
are domains.  The normalized intertwiner of AD.1 has the same
covariance on $F$, so $V^{\rm geom}_\sigma=\mathcal V_\sigma$.

Here is the independent orbifold identification of the paths.
The quotient stack $[Y/S_3]$ has geometric fundamental group
$B=C_2\amalg C_3$ and pure subgroup $F$.  The transposition
interchanging $0$ and $1$ fixes $p=1/2$; its isotropy is the
class of $s$.  Choose an ordinary geometric path $\xi$ in $Y$
from the tangential base point to $p$ so that the transported
orbifold isotropy is the named $s$.  Such a choice exists: all
order-two factor subgroups are conjugate in $B$, and conjugators
preserving the prescribed transposition in $S_3$ lie in
$F\langle s\rangle$, so the conjugation may be made pure.
For the order-three point one uses one of
$p=(1\pm\sqrt{-3})/2$ and its positive local isotropy, and makes
the same choice with the named $r$.  These two points account
for the two possible pure conjugacy classes; select the one
corresponding to $r$.  Both points and their coordinate-permutation
isotropies are fixed by $G_{\mathbb Q^{\rm ab}}$.

Paths are composed as maps of fibre functors: $\xi$ goes from
the base point to $p$, so
\[
 c_\sigma=(\sigma\xi)^{-1}\xi\in F
\]
is a loop at the base point.  If $s_p$ denotes the isotropy at
$p$ and $s=\xi^{-1}s_p\xi$, then
\[
 \Phi_\sigma(s)=(\sigma\xi)^{-1}s_p(\sigma\xi)
                  =c_\sigma s c_\sigma^{-1}.
\]
The same formula holds with $r$.  The actual arithmetic action
on this orbifold is the source's normalized extension: it agrees
on $F$ with the given tangential action, fixes $\tau$ when the
multiplier is one, and its extension is unique since the
centralizer of $F$ in $B$ is trivial.  Uniqueness of the pure
torsion frames therefore gives $c_\sigma=\beta_\sigma$ or
$c_\sigma=A_\sigma$, respectively.

For the remaining linear calculation use the Tate sheaf only on
$Y$, where the Legendre family is an actual family; no descent of
that sheaf through the projective modular stack is asserted.
If $\rho_p(\sigma)$ is its arithmetic matrix at the ordinary
point $p$, functoriality on the path gives
\[
 q(\sigma\xi)=\rho_p(\sigma)q(\xi)\mathcal V_\sigma^{-1}.
\]
Define $M_p=q(\xi)^{-1}\rho_p(\sigma)q(\xi)$ in the
transported integral basis.  Taking the inverse and multiplying
by $q(\xi)$ shows, with the order fixed,
\[
 q(c_\sigma)=\mathcal V_\sigma M_p^{-1}.
\]
Consequently
\[
 q(\beta_\sigma)M_p(\sigma)=\mathcal V_\sigma
 \quad\hbox{or}\quad
 q(A_\sigma)M_p(\sigma)=\mathcal V_\sigma.
\]
The matrices $M_p$ are exactly $\mathcal W_\sigma$ and
$\mathcal T_\sigma$.  The path cocycle is an actual element of
$F$, so its signed value is the already specified representation
$x\mapsto U^2,y\mapsto L^2$.  There is no independent scalar
lift to choose.  The calculation holds on every finite torsion
sheaf with compatible transition maps, hence on the full integral
Tate module at every prime.

The harmonic fixed point for the interchange is $t=1/2$.
Putting $X=X'+1/2$ in its equation gives
$Y^2=X'^3-X'/4$, the quadratic twist of $E_4$ by $1/2$.
Its twisting square root belongs to $\mathbb Q^{\rm ab}$, so
its representation on $G_{\mathbb Q^{\rm ab}}$ is $\rho_4$.
For the equianharmonic point $t=(1+\sqrt{-3})/2$, since
$t^2-t+1=0$, the substitution
$X=X'+(1+t)/3$ gives
\[
 Y^2=X'^3+\frac{\sqrt{-3}}9
      =X'^3+\left(\frac1{\sqrt{-3}}\right)^3.
\]
This is the quadratic twist of $E_3$ by $d=1/\sqrt{-3}$.
If a square root of $d$ is chosen, its quotient by $r_0$ lies
in $\mathbb Q^{\rm ab}$ (its square is $-2i/3$, and one may
use $\sqrt{2/3}\,\zeta_8^{-1}$).  Thus its twisting character
on $G_{\mathbb Q^{\rm ab}}$ is $\kappa$, as claimed.
The conjugate equianharmonic point changes the twist by $-1$,
whose square root lies in $\mathbb Q^{\rm ab}$, and yields the
same conclusion.  Choose the transported CM identifications
so that their centralizers are those of $S$ and $R_0$, in the
signed convention of AD.2.  This gives the asserted exact
integral matrix identities in those bases.\hfill$\square$

\paragraph{D23.C6 (proved from D23.C4--C5 and the stipulated inputs).}
The signed adelic frame pair has exact images
\[
 (\mathcal T,\mathcal W)(H)=N_3(4)\times N_4(4),
\]
\[
 (\mathcal T,\mathcal W)(G_{\mathbb Q^{\rm ab}})
 =\{(t,w):\epsilon_3(t\bmod3)=\epsilon_2(w\bmod3)\}.
\]
Every continuous relative abelian character of $G$ whose
restriction to $H$ factors through the full adelic frame pair
factors through $\Delta$.  In particular its quadratic possibilities
are exactly $0$ and $\Delta$.

\emph{Proof.}
D23.C4 and D23.C5 show that the arithmetic pair projects onto
both normalized tori, and that their kernels on
$G_{\mathbb Q^{\rm ab}}$ cut out $K_3^{\rm ab}$ and
$K_4^{\rm ab}$.  DELTA\_ON\_GT A.12 gives their exact
intersection
\[
 K_3^{\rm ab}\cap K_4^{\rm ab}
       =\mathbb Q^{\rm ab}(3^{1/4}),
 \qquad [\mathbb Q^{\rm ab}(3^{1/4}):\mathbb Q^{\rm ab}]=2.
\]
The compact Galois image is therefore the fibre product over a
quotient of order two, and has index two in the torus product.
E.65's already checked comparison modulo twelve puts it in the
equal-parity subgroup displayed above.  That subgroup itself
has index two, so they are equal.

AD.2 puts the full gentle frame image inside the normalized torus
product.  It contains the arithmetic index-two subgroup and the
E.58 element, whose parity pair is $(a_3,a_2)=(1,0)$.
A subgroup containing an index-two subgroup and one element
outside it is the entire group.  This proves the first equality.
The quotient of the full frame image by the arithmetic image is
thus $C_2$, with quotient character $a_2+a_3$.  Any continuous
abelian character satisfying the factorization and arithmetic
vanishing hypotheses descends to that quotient.  Finally
$G=HG_{\mathbb Q}$ makes its full extension unique, and E.65
identifies the nonzero quadratic one with $\Delta$.
\hfill$\square$

\paragraph{D23.C7 (proved; precise failure of a place-two frame factorization).}
There is an actual $\sigma\in G_{\mathbb Q^{\rm ab}}$ with
\[
 a_2(\sigma)=a_3(\sigma)=1,
 \qquad
 q_2(A_\sigma)=q_2(\beta_\sigma)
     =\mathcal V_{\sigma,2}
     =\mathcal T_{\sigma,2}
     =\mathcal W_{\sigma,2}=I.
\]
Here $q_2$ means the full $2$-adic representation, not reduction
modulo two.  Thus neither parity is a function solely of the pair
of raw signed $2$-adic frame matrices, or of the corrected pair,
even if the $2$-adic intertwiner is also supplied.  This is a
statement about those precise congruence data; it makes no claim
about every possible meaning of a ``local'' geometric invariant.

There is also an actual $h\in H$
with $\Delta(h)=1$, $a_3(h)=0$, and $a_2(h)=1$, whose entire
corrected adelic pair has
\[
 \mathcal T_h=I,\qquad
 \mathcal W_{h,3}=S,\qquad
 \mathcal W_{h,\ell}=I\quad(\ell\ne3).
\]
In particular its corrected $2$-adic pair is the identity pair.

\emph{Proof.}
Fix an embedding into a decomposition group at three and write
$I_3$ for inertia.  Its wild subgroup is pro-three and its tame
quotient is $\prod_{\ell\ne3}\mathbb Z_\ell(1)$.  A pro-two
Sylow subgroup $P$ therefore maps isomorphically to the
$\mathbb Z_2(1)$ factor: its intersection with the wild subgroup
is trivial, and the image of a Sylow subgroup under a compact
surjection is a Sylow subgroup.  These assertions can also be
obtained by taking compatible Sylow subgroups in the finite
inertia quotients.  Choose a topological generator $\tau$ of $P$.
The extension generated by $r_0$ is tamely ramified of degree
four after adjoining $\mu_4$, which is unramified at three;
indeed $(2r_0)^4=12$ is Eisenstein at three.  It follows that
$\tau(r_0)=i^{\pm1}r_0$.  Put $\sigma=\tau^2$.  Then
$\sigma(r_0)=-r_0$.

For every $\ell\ne3$, the $\ell$-adic cyclotomic character is
unramified at three.  The image of the pro-two group $P$ under
$\chi_3$ is contained in $\{1,-1\}\subset\mathbb Z_3^\times$,
so its square is one.  Consequently $\chi(\sigma)=1$ and
$\sigma\in G_{\mathbb Q^{\rm ab}}$.  R4's exact finite arithmetic
comparison identifies both parities with the sign on $r_0$;
therefore $a_2(\sigma)=a_3(\sigma)=1$.

The two Legendre fibres used in the path proof of D23.C5 have
good reduction at three.  The harmonic fibre at $t=1/2$ has
discriminant $16t^2(1-t)^2=1$.  For either equianharmonic fibre,
$t(1-t)=1$, so the discriminant is $16$, a unit at three.
The latter fibre is defined over $K_3$, which $\sigma$ fixes;
$\sigma\in I_3\cap G_{K_3}$ is in its local inertia group.
The good-reduction implication of Serre--Tate,
\emph{Good reduction of abelian varieties}, Theorem 1,
\emph{Annals of Mathematics} 88 (1968), p.493, makes both full
$2$-adic Tate actions of $\sigma$ the identity.  The integral
path identity in D23.C5 therefore gives
$\mathcal T_{\sigma,2}=\mathcal W_{\sigma,2}=I$.

The actual tangential action on the free pro-two group is
unramified at three.  One explicit source for exactly this
automorphism statement is Hain--Matsumoto,
\emph{Weighted completion of Galois groups and Galois actions on
the fundamental group of $\mathbb P^1-\{0,1,\infty\}$},
\emph{Compositio Mathematica} 139 (2003), Section 8.5,
the paragraph before Theorem 8.7 (take $N=1$ and $\ell=2$).
Thus $\Phi_\sigma$ acts identically on $F^{(2)}$.
The image $q_2(F)$ lies in the principal congruence subgroup
modulo two of $\mathrm{SL}_2(\mathbb Z_2)$, which is pro-two.
Hence $q_2\Phi_\sigma=q_2$.  The covariance says that the
upper unipotent $\mathcal V_{\sigma,2}$ commutes with $U^2,L^2$,
forcing it to be $I$.  The corrected-frame definitions then give
$q_2(A_\sigma)=q_2(\beta_\sigma)=I$ as well.

The identity element has all those same matrices and parity zero,
which proves the asserted failure of factorization.  For the final
assertion, the displayed adelic pair belongs to
$N_3(4)\times N_4(4)$: its component over two is the identity,
and $S$ has norm one at three.  D23.C6 supplies an actual $h$
with that value.  Its parities and discrepancy have the stated
values by their definitions.\hfill$\square$

\paragraph{D23.C8 (proved; an explicit smaller gate for B.40).}
Let
\[
 N=\ker\bigl((\mathcal T,\mathcal W):H\longrightarrow
                   N_3(4)\times N_4(4)\bigr),
 \qquad
 Q=G/\overline{[G,G]G^2G_{\mathbb Q}}.
\]
Then, as a set product,
\[
 \ker\Delta=N G_{\mathbb Q}.
\]
If $\bar\Delta:Q\to\mathbb F_2$ is the induced nonzero
character, its kernel is exactly the image of $N$.  Consequently
the assertion that all relative continuous quadratic characters
are $0,\Delta$ is equivalent to the concrete vanishing condition
\[
 N\subseteq\overline{[G,G]G^2G_{\mathbb Q}}.
\]
This computes the entire remaining kernel in terms of elements
whose two full adelic frames are identity.  It does not assert
that this subgroup has trivial image in $Q$.

For the two particular characters $d$ and $\Delta\circ D$,
the known actual kernel witness of D23.C3 gives the further
equivalences
\[
 \begin{split}
 d\ne0&\quad\Longleftrightarrow\quad
      \exists n\in N:\ a_2(Dn)=1,\\
 \Delta\circ D\ne0&\quad\Longleftrightarrow\quad
      \exists n\in N:\ \Delta(Dn)=1.
 \end{split}
\]
In particular both characters vanish if and only if
$(a_2,a_3)(D(N))=\{(0,0)\}$.

\emph{Proof.}
Write $g=h\sigma$ with $h\in H$ and $\sigma\in G_{\mathbb Q}$.
If $g\in\ker\Delta$, then $h\in\ker\delta$, so its two
adelic frame values have equal parity.  D23.C6 gives an actual
$\tau\in G_{\mathbb Q^{\rm ab}}$ with exactly the same entire
frame pair as $h$.  The two frame maps are homomorphisms on
$H$, hence $n=h\tau^{-1}\in N$ and $g=n(\tau\sigma)$.
The converse follows from $\Delta|_N=\Delta|_{G_{\mathbb Q}}=0$.
No normality assertion about $N$ in all of $G$ is needed.

The closed subgroup in the definition of $Q$ is the inverse
image of the image of $G_{\mathbb Q}$ in the maximal elementary
abelian pro-two quotient, so the quotient is well defined.
The preceding set equality shows
$\ker\bar\Delta=\operatorname{im}(N\to Q)$.  Continuous
quadratic characters separate points of an elementary abelian
pro-two group.  Thus they are exactly $0,\bar\Delta$ if and
only if this kernel is trivial.  Pullback gives the asserted
equivalence on $G$.  In the integrated relative-abelianization
statement P.12, $Q=K_G/2K_G$.

On $N$, one has $a_2=a_3=0$, so
$d(n)=a_2(Dn)$.  Any relative character vanishing on $N$
vanishes on $\ker\Delta=N G_{\mathbb Q}$ and is therefore
$0$ or $\Delta$.  The actual $k$ of D23.C3 has
$d(k)=(\Delta\circ D)(k)=0$ but $\Delta(k)=1$;
it excludes the second possibility for both characters.
This proves the two equivalences.  Their right sides both
fail exactly when $a_2(Dn)=0$ and
$a_2(Dn)+a_3(Dn)=0$ for every $n\in N$, equivalently when
both descended parities vanish there.\hfill$\square$


\section{Finite comparison gates and failed shortcuts}
% DELTA23 agent /root/pentagon. New finite-quotient obstruction, not A.84 decision.
\subsection{Independent central flips in all four-point forgetting quotients}

\paragraph{D23.P1 (independent central coordinates).}
Let $P_n=\widehat\Gamma_{0,n}$, $n\geq4$. For each four-element
subset $I$ of the labelled punctures, let $\pi_I:P_n\to P_4$ be
forgetting the other punctures, with any orientation-preserving marking
of the target. Use the signed modular quotient of LM7\_PUSH E.58.11,
\[
 q:P_4\cong\widehat F_2\longrightarrow
 V_4\times\operatorname{SL}_2(\mathbb F_3),
\]
whose level-twelve projective notation sends $(1,-I)$ to $[5I]$.
Write $z=(1,-I)$. Then the simultaneous map has exact image
\[
 (q\pi_I)_I(P_n'')=\prod_{|I|=4}\langle z\rangle.
 \tag{D23.P1.1}
\]
Here the primes denote closed derived subgroups. In particular the
central signs in distinct four-point quotients can be prescribed
independently by one actual element of $P_n''$.

\begin{proof}
The second derived subgroup of $V_4\times\operatorname{SL}_2(\mathbb F_3)$
is $\{1\}\times\{I,-I\}$. The relevant explicit matrices are
\[
 x=\begin{pmatrix}1&2\\0&1\end{pmatrix},\qquad
 y=\begin{pmatrix}1&0\\1&1\end{pmatrix}.
\]
They generate $\operatorname{SL}_2(\mathbb F_3)$; its commutator
subgroup is the quaternion group of order eight and that group's
commutator subgroup is $\{I,-I\}$. These finite assertions, as well
as the matrix calculation below, are enumerated in the certificate
\path{agents/pentagon/check_four_point_central_flips.py}.
This gives the inclusion from left to right in (D23.P1.1).

Fix $I=\{i_1,i_2,i_3,i_4\}$. Choose a disk neighbourhood of a
chain joining these four punctures, with no other punctures in it.
Let $a,b,c\in P_n$ be the positive twists about the consecutive
pairs $(i_1,i_2),(i_2,i_3),(i_3,i_4)$ in that disk, and set
\[
 w_I=[[a,b],a[b,c]a^{-1}],\qquad [u,v]=uvu^{-1}v^{-1}.
 \tag{D23.P1.2}
\]
This lies in $P_n''$. If any puncture of $I$ is forgotten, one of
the two inner commutators is trivial: forgetting $i_1$ kills $a$,
forgetting $i_4$ kills $c$, and forgetting $i_2$ or $i_3$ kills
both twists in at least one inner commutator. A twist about a disk
containing at most one retained puncture is trivial in the punctured
sphere mapping class group. Therefore $\pi_J(w_I)=1$ for every
four-element subset $J\ne I$.

After forgetting all punctures outside $I$, the curves defining $a$
and $c$ are complementary and isotopic in the four-punctured sphere.
Thus $\pi_I(c)=\pi_I(a)$, and in the chosen standard marking their
signed modulo-three images are $q_3(a)=q_3(c)=x$, $q_3(b)=y$.
Direct multiplication gives
\[
 [x,y]=\begin{pmatrix}1&2\\2&2\end{pmatrix},\quad
 x[y,x]x^{-1}=\begin{pmatrix}1&1\\1&2\end{pmatrix},\quad
 [[x,y],x[y,x]x^{-1}]=-I.
\]
The modulo-four coordinate is trivial because $w_I\in P_n''$ and
the modulo-four target is abelian. Hence $q\pi_I(w_I)=z$.
Any other orientation-preserving marking is obtained by a mapping
class on the four-punctured sphere. Its conjugation action preserves
the modular quotient and its central element $z$; equivalently one
uses the quotient $\Gamma_{0,[4]}\to C_2\amalg C_3$ and its normal
level-twelve kernel. Thus the value $z$ does not depend on the
chosen chain marking. The images of the $w_I$ are the coordinate
basis of the product in (D23.P1.1). Their finite products give every
prescribed tuple. This proves the reverse inclusion.
\end{proof}

\paragraph{D23.P2 (finite simultaneous frames do not remove the cubic flip).}
Let a finite subgroup $H$ of $\widehat\Gamma_{0,[n]}$ act on the
four-element subsets by its label permutation, and retain the finite
quotient obtained from all the preceding four-point modular maps
and the label permutation map. For every $H$-orbit $\mathcal O$ of
four-element subsets there is an element of the image of $P_n''$
that centralizes the image of $H$ and whose four-point coordinates
are $z$ on $\mathcal O$ and $1$ elsewhere.
Consequently, whenever a finite simultaneous frame for $H$ is used to
read an equianharmonic frame sign at a subset $I\in\mathcal O$,
this centralizer changes that cubic sign by one. At any four-point
harmonic coordinate it leaves the harmonic parity unchanged. It also
leaves the abelianization of the simultaneous pure frame unchanged.

\begin{proof}
Conjugation by an element of the full mapping class group transports
forgetting maps according to its permutation of the labels. On the
second-derived central coordinates its effect is exactly that
permutation: $z$ is the distinguished nonidentity element of a group
of order two and is fixed by changes of target marking. By D23.P1,
the orbit-indicator tuple is the image of a product of the $w_I$.
It is fixed by conjugation by every element of $H$, hence centralizes
its image in the stated finite quotient. Multiplying any chosen
simultaneous frame by this element preserves all its finite
conjugator equations.

On the cubic coordinate, $z$ is $-I$ in
$C_{\operatorname{SL}_2(\mathbb F_3)}(R_0)\cong C_6$; its sign
on the rational fixed line is $-1$, so $a_3$ changes by one.
On the harmonic coordinate, $z=-I=S^2$ in
$\langle S\rangle\cong C_4$; the parity of the exponent of $S$
does not change. Trivial coordinates do not change that parity
either. Since the multiplying element lies in $P_n''$, the pure
frame's abelianization is also unchanged.
\end{proof}

\paragraph{Exact scope.}
D23.P2 is a statement about the specified finite quotient and finite
conjugator equations. It does not construct a different actual
unique frame in the full profinite group, nor a different gentle
or GT element. In particular, the additional affine five-point
quotient used in DELTA\_ON\_GT FP.32 records more than these
four-point projections and can exclude this ambiguity by its
relation $a_3=\kappa_d\bmod2$. No claim is made that D23.P2
survives that additional quotient. This computation therefore
blocks one simultaneous-frame shortcut but does not decide A.84.

% DELTA23: /root/pentagon, 2026-11-21 CET. Full proof draft.
\subsection{A small affine evaluation detecting the cubic frame sign}

\paragraph{D23.P3 (an order-1944 free-group detector).}
Put $k=\mathbb F_9=\mathbb F_3[i]$, $i^2=-1$, and use affine
multiplication $(v,M)(w,N)=(v+Mw,MN)$ on $k^2\rtimes\operatorname{GL}_2(k)$.
Retain the modular generators $s,r$, $x=(sr)^{-2}$, $y=sxs$ of
LM7\_PUSH E.58 and E.65. Set
\[
 S=\begin{pmatrix}0&1\\-1&0\end{pmatrix},\quad
 R=\begin{pmatrix}0&-1\\1&-1\end{pmatrix},\quad
 U=\begin{pmatrix}1&1\\0&1\end{pmatrix},\quad v=e_1.
\]
There is a continuous homomorphism
\[
 \eta:\widehat B=C_2\amalg C_3\longrightarrow E,
 \qquad \eta(s)=(0,iS),\qquad \eta(r)=(v,R),
 \tag{D23.P3.1}
\]
onto the finite group
\[
 E=k^2\rtimes H,\qquad H=\langle iI\rangle\operatorname{SL}_2(\mathbb F_3),
 \quad |E|=3888.
\]
Its restriction to the pure free group $F=\widehat F(x,y)$ has
image $E_F=k^2\rtimes\operatorname{SL}_2(\mathbb F_3)$, of order
$1944$. The image of $C_E(\eta(r))$ in $H$ is exactly $\langle R\rangle$,
and its order is $27$.

For an actual multiplier-one gentle element $h=(1,f)$, write
\[
 \eta(f)=(d,M_\rho),\quad d=(d_1,d_2)^{\mathsf t},\quad
 M_\rho=\begin{pmatrix}1&\rho\\\rho&1+\rho^2\end{pmatrix},
 \qquad\rho\in\mathbb F_3.
\]
Here $\rho$ is the reduction of the frame parameter in E.65; it is
also read directly as the lower-left entry of the linear part.
Then
\[
 \boxed{\quad (-1)^{a_3(h)}
       =(1-\rho)(1-d_1)-d_2\quad\text{in }k.\quad}
 \tag{D23.P3.2}
\]
The right-hand side belongs to $\{1,-1\}\subset\mathbb F_3$ for
every actual gentle element. Moreover, if $\epsilon=(-1)^{a_3(h)}$,
its full affine evaluation is exactly
\[
 d_1=\frac{1-\rho-\epsilon}{1+i},\qquad
 d_2=(\rho+i)d_1.
 \tag{D23.P3.3}
\]
Thus every such $f$-evaluation lies in an explicitly specified set of six elements of $E_F$.
For the E.58 witness one obtains
\[
 \eta(f_{58})=((1-i,1+i)^{\mathsf t},I),
 \tag{D23.P3.4}
\]
not merely the value of its sign.

\begin{proof}
The matrix $iS$ squares to $I$. In characteristic three,
$(R-I)^2=0$, so $I+R+R^2=0$. These identities prove the two
torsion relations in (D23.P3.1). Hence the homomorphism exists
by the defining free-product presentation, and continuity follows
because its target is finite.

The linear image is $H$: indeed $iSR=iU^{-1}$ and
$(iU^{-1})^3=-iI$, so it contains $iI,S,R$, which generate $H$.
The intersection $\langle iI\rangle\cap\operatorname{SL}_2(\mathbb F_3)$
is $\{I,-I\}$; hence $|H|=4\cdot24/2=48$.
The translation subgroup $N$ of the image is an $H$-invariant
$\mathbb F_3$-subspace of $k^2$. Scalar multiplication by $iI$
makes it a $k$-subspace. The natural representation of
$\operatorname{SL}_2(\mathbb F_3)$ over $k$ is irreducible:
a common invariant line for the two nonidentity unipotent
matrices $U$ and $L=\left(\begin{smallmatrix}1&0\\-1&1\end{smallmatrix}\right)$
would have to be both $ke_1$ and $ke_2$. Thus $N$ is zero or $k^2$.

If $N=0$, the affine image is the graph of a cocycle $c:H\to k^2$.
The central element $z=-I$ acts as $-1$. Comparing $c(gz)$ and
$c(zg)$ gives $2c(g)=(1-g)c(z)$, so
$c(g)=(1-g)t$ with $t=c(z)/2$. In particular this would imply
$v=c(R)\in(1-R)k^2$. But $(1-R)k^2=k(1,-1)^{\mathsf t}$
and $v=e_1$ is outside that line. This contradiction proves
$N=k^2$ and the asserted full image.

The linear images of $x,y$ are $-U^2,-L^2$. Cubing either
produces $-I$, so they generate $\operatorname{SL}_2(\mathbb F_3)$.
The pure image is normal in $E$, since $F$ is normal in $B$.
Its intersection with $k^2$ is again zero or all of $k^2$.
The zero alternative would make the quotient of $E$ by the pure
image have order $3888/24=162$, whereas it is a quotient of
$B/F\cong S_3$. It follows that the pure image is exactly $E_F$.

An element of the centralizer of $R$ in $H$ has the form
$i^jR^m$, $j\in\mathbb Z/4$, $m\in\mathbb Z/3$.
Indeed $C_{\operatorname{SL}_2(\mathbb F_3)}(R)
=\langle-I\rangle\times\langle R\rangle$, and the scalar
extension adds $iI$. For an affine element $(t,i^jR^m)$ to
centralize $(v,R)$ one must have
$(1-R)t=v-i^jR^m v$. In the one-dimensional coinvariant quotient
$k^2/(1-R)k^2$, the nonzero class of $v$ is fixed by $R$.
Consequently $i^j=1$. Conversely all powers of $(v,R)$ centralize
it, so the linear image of its centralizer is exactly
$\langle R\rangle$. Each of its three linear values has
$|\ker(1-R)|=9$ translation lifts. This proves the order $27$.

We next justify the signed comparison between the linear part of
$\eta$ and the supplied $q_{12}$, rather than identifying the
representations without a sign convention. There is a character
$\chi_4:\operatorname{SL}_2(\mathbb Z/4)\to\langle i\rangle$
with $\chi_4(U)=\chi_4(L)=-i$, $\chi_4(S)=i$,
$\chi_4(R)=1$, and $\chi_4(-I)=-1$. Its finite definition is
verified by enumerating all $48$ matrices and propagating the
values along multiplication by $U,L$; the certificate checks
consistency on every such edge. Thus
\[
 [(g_4,g_3)]\longmapsto\chi_4(g_4)g_3
\]
is a well-defined homomorphism from
$(\operatorname{SL}_2(\mathbb Z/4)\times
\operatorname{SL}_2(\mathbb F_3))/\langle(-I,-I)\rangle$ to $H$.
It sends the modular $s,r$ to $iS,R$, precisely the linear part
of $\eta$. For the actual pure frame $A_h$, E.65 proves
$q_4(A_h)=I$, so its linear image under $\eta$ is exactly
$q_3(A_h)$. For $f\in F'$ the modulo-four coordinate is also
trivial. E.65.9 therefore gives its displayed linear part
$M_\rho$, and $V=U^{2\rho}$ is the reduction of E.65's shear.

Write $\eta(A_h)=(t,A_0)$ and $T=A_0^{-1}V$. By E.65,
$T\in\langle-I\rangle\times\langle R\rangle$, and its sign
on the fixed line is $\epsilon=(-1)^{a_3(h)}$. The exact frame
identity $h(r)=A_hrA_h^{-1}=f^{-1}r$ gives, on translations,
\[
 M_\rho^{-1}(v-d)
       =A_0v+(1-A_0RA_0^{-1})t.
\]
Multiplying by $V^{-1}$ and passing to $R$-coinvariants makes
the last term zero. The row vector $\ell=(1,1)$ realizes this
coinvariant quotient; it satisfies $\ell R=\ell$, $\ell v=1$,
and $\ell T^{-1}=\epsilon\ell$. We obtain
\[
 \ell V^{-1}M_\rho^{-1}(v-d)=\epsilon.
\]
Direct multiplication gives
$V^{-1}M_\rho^{-1}=\left(\begin{smallmatrix}1&0\\-\rho&1\end{smallmatrix}\right)$,
which proves (D23.P3.2).

Finally $h(s)=sf$ has order two. The translation part of
$\eta(sf)^2=1$, with the displayed matrices, is equivalent to
$d_2=(\rho+i)d_1$. Substitution in (D23.P3.2) gives
$\epsilon=1-\rho-(1+i)d_1$, proving (D23.P3.3).
E.58 has $\rho=0$ and $a_3=1$, which gives (D23.P3.4).
The exact finite program \path{agents/pentagon/affine_cubic_detector.py}
independently enumerates both images, the centralizer, and all
six affine evaluations, and checks the coinvariant calculation
for all $486$ combinations of sign, cubic power, and translation.
\end{proof}

\paragraph{D23.P4 (an explicit finite evaluation gate for A.84).}
Retain $p=(123)$ and $t=(1234)$ in $S_4$. The simultaneous map
\[
 F\longrightarrow E_F\times S_4,\qquad
 x\longmapsto(\eta(x),p),\quad
 y\longmapsto(\eta(y),t)
 \tag{D23.P4.1}
\]
is onto a group of order $46656$. For an actual gentle element
$h=(1,f)$, write its $S_4$-evaluation as
$f(p,t)=t^{2\nu_4}p^{a_f}$, as in DELTA\_ON\_GT FP.37.
For actual full GT elements, the existing pentagon consequences
FP.33 and FP.37 give $a_f=\rho$ and $a_2=\nu_4$.
Consequently A.84 is equivalent to the following concrete finite
comparison on actual full-GT source residues:
\[
 \boxed{\quad
 (1-\rho)(1-d_1)-d_2=(-1)^{\nu_4}.
 \quad} \tag{D23.P4.2}
\]
After the two previously proved pentagon consequences, there are
only twelve possible joint evaluations: $\rho\in\mathbb F_3$
and $(a_3,\nu_4)\in\mathbb F_2^2$. Six obey (D23.P4.2) and
six fail it. The remaining question is whether any of the latter
six is the image of an actual full-GT element. This is not an
assertion that any particular finite candidate lifts to GT.

\begin{proof}
Since $-I$ acts as $-1$ on $k^2$ and $2$ is invertible, commutators
with $-I$ generate the translation subgroup. Hence
$E_F^{\mathrm{ab}}\cong\operatorname{SL}_2(\mathbb F_3)^{\mathrm{ab}}
\cong C_3$. On the other hand $S_4^{\mathrm{ab}}\cong C_2$.
Both groups are finite solvable. Any nontrivial common quotient
would be solvable and would have a nontrivial abelian quotient;
its abelianization would be a quotient of both $C_3$ and $C_2$,
which is impossible. Both projections of (D23.P4.1) are onto,
so Goursat's lemma proves that its image is the full product.
Its order is $1944\cdot24=46656$.

D23.P3 computes the cubic bit from its first coordinate, while
FP.37 computes the harmonic bit on full GT from its second
coordinate. Thus their equality is exactly (D23.P4.2).
The six possible first coordinates are (D23.P3.3); the second
coordinates have the six values $t^{2\nu_4}p^{a_f}$.
Imposing $a_f=\rho$ leaves the stated twelve candidates and
partitions them according to the displayed sign equality.
The reduction from arbitrary multiplier to multiplier one is
LM7\_PUSH E.65.6, using Galois vanishing, exactly as in
DELTA\_ON\_GT AO.1. This proves the equivalence with A.84.
\end{proof}

\subsection{The cubic bit in an order-216 marked affine quotient}

\paragraph{D23.P10 (a nine-fold smaller exact main gate).}
Let $V=\mathbb F_3^2$, $E_0=V\rtimes\operatorname{SL}_2(\mathbb F_3)$,
and use affine multiplication $(v,M)(w,N)=(v+Mw,MN)$.
There is a continuous surjection from the same marked pure free
group $F=\widehat F(x,y)$ as in E.65,
\[
 \xi:F\twoheadrightarrow E_0,\qquad
 \xi(x)=\left(0,\begin{pmatrix}-1&1\\0&-1\end{pmatrix}\right),\qquad
 \xi(y)=\left(e_1,\begin{pmatrix}-1&0\\-1&-1\end{pmatrix}\right).
\tag{D23.P10.1}
\]
The two linear matrices are $-U^2,-L^2$ in the source's positive
convention. For every actual multiplier-one gentle element
$h=(1,f)$ one has the exact formula
\[
 \boxed{\quad
 \xi(f)=\left(
   \begin{pmatrix}e\\\rho+(\rho+1)e\end{pmatrix},
   \begin{pmatrix}1&\rho\\\rho&1+\rho^2\end{pmatrix}
  \right),\qquad e=a_3(h)\in\{0,1\}\subset\mathbb F_3.
 \quad}
\tag{D23.P10.2}
\]
Thus the cubic bit is literally the first translation coordinate
of this order-$216$ finite evaluation of $f$.

The joint marked map
\[
 F\longrightarrow E_0\times S_4,\qquad
 x\longmapsto(\xi(x),(123)),\qquad
 y\longmapsto(\xi(y),(1234))
\tag{D23.P10.3}
\]
is onto, of order $5184$. Its image on $F'$ has order $864$ and
is exactly
\[
 (V\rtimes Q_8)\times A_4.
\tag{D23.P10.4}
\]
For actual full-GT elements, the supplied pentagon consequences
$a_f=\rho$ and $a_2=\nu_4$ leave precisely the same twelve
candidate states as D23.P4, now in a nine-fold smaller source
quotient. In these coordinates A.84 is exactly the assertion
that none of the six states with
\[
                  e\ne\nu_4
\tag{D23.P10.5}
\]
is attained by an actual full-GT element. The number of possible
mixed states has not decreased; the finite group realizing the
exact comparison has.

Among quotients of $E_0$, no proper quotient can retain this
detector on the supplied E.58 witness: every nontrivial normal
subgroup of $E_0$ contains $V$, whereas
\[
       \xi(f_{58})=((1,1)^{\mathsf T},I)\ne1.
\tag{D23.P10.6}
\]
This is a minimality statement within the quotient family of
$E_0$, not a claim of global minimality among all finite detectors.

\begin{proof}
First construct the map directly from the free presentation.
Write $X=\xi(x)$ and $Y=\xi(y)$. Their linear images generate
$\operatorname{SL}_2(\mathbb F_3)$: cubing either gives $-I$,
after which $U^2,L^2$ and hence $U,L$ are available. Moreover
\[
              [X^3,Y]=(e_1,I).
\]
Indeed $X^3=(0,-I)$, and conjugation by it negates the
translation of $Y$ while leaving its linear part unchanged;
the resulting commutator translation is $-2e_1=e_1$ in
characteristic three. Conjugating this translation by $Y$
gives the translation $(-1,-1)^{\mathsf T}$. Together these
span $V$. Thus the image is all of $E_0$, of order $9\cdot24=216$.

We next compare the marked map with D23.P3; this step fixes
the actual cubic sign rather than guessing from an isomorphic
finite group. On $E_F=\mathbb F_9^2\rtimes\operatorname{SL}_2(\mathbb F_3)$,
real-part projection of the translation is a group homomorphism,
because every linear matrix has entries in $\mathbb F_3$.
D23.P3's two pure generators have real translation $(1,1)$
and linear parts $-U^2,-L^2$. Change affine origin by
$t=(0,2)^{\mathsf T}$, which changes a translation $d$ of
linear part $M$ into $d-(I-M)t$. The two resulting translations
are $0,e_1$, precisely (D23.P10.1).

For an actual gentle $f$, D23.P3 gives
\[
 d_1=\frac{1-\rho-\epsilon}{1+i},\qquad
 d_2=(\rho+i)d_1,\qquad \epsilon=(-1)^{a_3(h)}.
\]
Since $(1+i)^{-1}=2+i$, writing
$c=1-\rho-\epsilon\in\mathbb F_3$ gives
\[
 \operatorname{Re}d_1=2c,\qquad
 \operatorname{Re}d_2=(2\rho-1)c.
\]
The origin correction is $(I-M_\rho)t=(\rho,\rho^2)^{\mathsf T}$.
If $(e,b)^{\mathsf T}$ denotes the corrected translation, then
\[
 e=2(1-\rho-\epsilon)-\rho,\qquad \epsilon=1+e.
\]
Thus $e=0$ when $a_3=0$ and $e=1$ when $a_3=1$.
Substituting $c=-\rho-e$ into the second coordinate gives
\[
 b=(2\rho-1)(-\rho-e)-\rho^2=\rho+(\rho+1)e.
\]
This proves (D23.P10.2), and E.58's $\rho=0,a_3=1$ gives
(D23.P10.6). The use of E.58 here is the supplied actual
witness; its absent infinite lifting input is not independently
reproved by this finite calculation.

Because $-I$ acts as $-1$ on $V$, commutators with it generate
all translations. Therefore $E_0'=V\rtimes Q_8$ and
$E_0^{\mathrm{ab}}\simeq C_3$. Since
$S_4^{\mathrm{ab}}\simeq C_2$ and both groups are solvable,
they have no nontrivial common quotient: such a quotient would
have a nontrivial abelian quotient common to $C_3,C_2$.
Goursat's lemma proves surjectivity of (D23.P10.3). Taking
derived subgroups proves (D23.P10.4). The further derived orders,
also verified in the exact certificate, are
\[
 |E_0^{(j)}|:216,72,18,9,1;\qquad
 |(E_0\times S_4)^{(j)}|:5184,864,72,9,1.
\]
On full GT, supplied FP.33 and FP.37 identify the marked
$S_4$ value as $t^{2\nu_4}p^\rho$ and identify $a_2=\nu_4$.
Together with (D23.P10.2), their equality criterion is exactly
(D23.P10.5). Reduction from arbitrary multiplier to multiplier
one is the supplied E.65.6 normalization, as in D23.P4.

Finally let $N\lhd E_0$ be nontrivial. If $N\cap V\ne1$,
the irreducibility of the natural $\operatorname{SL}_2(\mathbb F_3)$
module implies $V\subset N$. If $N\cap V=1$, then
$[N,V]\subset N\cap V$ forces $N$ to centralize $V$.
The centralizer of all translations in $E_0$ is $V$ itself,
since the linear action is faithful. Hence $N\subset V$,
contradicting $N\cap V=1$ and $N\ne1$. Every proper quotient
therefore kills the nonidentity translation in (D23.P10.6),
as claimed.

The exact program \path{agents/pentagon/cubic_detector_216.py}
checks the marked commutator, enumerates the affine image,
checks all six projected values from D23.P3, and independently
computes the order and derived series of the joint action on
nine affine points plus four permutation points.
\end{proof}

\subsection{The cubic detector does not extend through one natural rank-four module}
\paragraph{D23.P5 (a checked obstruction to the twisted symplectic affine route).}
Let $k=\mathbb F_9=\mathbb F_3[i]$, and let $t_1,t_2,t_3,t_4$ be
D23.T1's rank-four symplectic transvections. On the standard subgroup
$B_4\subset B_5$, put $A_j=-it_j$, $1\leq j\leq3$. Every
continuous affine lift of this linear representation has a common
fixed point; equivalently every continuous cocycle
$B_4\to k^4$ is a coboundary. In particular the cubic affine detector
D23.P3, on the nondegenerate span $W=\langle e_1,e_2\rangle$ of
the first two braid generators, does not extend through this
rank-four module. The analogous extension to $B_5$ also fails.

\begin{proof}
Write $z_j$ for the translation of an affine lift of $A_j$.
The matrix $I-A_1$ is invertible, because $t_1$ is unipotent and
$1+i\ne0$. A change of affine origin therefore makes $z_1=0$.
The distant commutation of the first and third generators now gives
$(I-A_1)z_3=0$, hence $z_3=0$. The two remaining braid equations
on translations give
\[
 (A_1-I-A_2A_1)z_2=0,\qquad
 (I+A_2A_3-A_3)z_2=0.
\]
The matrices are, respectively,
\[
 K_1=\begin{pmatrix}
 -i&1-i&0&0\\
 -1&-1-i&1&0\\
 0&0&-i&0\\
 0&0&0&-i
 \end{pmatrix},\qquad
 K_2=\begin{pmatrix}
 i&0&0&0\\
 1&1+i&-1&-1\\
 0&1-i&i&-1+i\\
 0&0&0&i
 \end{pmatrix}.
\]
Rows three and four of $K_1$ force the third and fourth coordinates
of $z_2$ to vanish; row one of $K_2$ then forces its first coordinate
to vanish, and row one of $K_1$ forces the second to vanish.
All translations are zero in the new origin, which proves the
coboundary assertion.

The restriction of $k^4$ to $B_3=\langle\sigma_1,\sigma_2\rangle$
splits as $W\oplus W^\perp$, where
$W^\perp=\langle e_1+e_3,e_4\rangle$. Projection of a coboundary
to $W$ is again a coboundary. D23.P3's cocycle is not a coboundary
on $W$: its value on $r$ is $e_1\notin(1-R)W$.
This excludes the proposed extension. The exact solver in
\path{agents/pentagon/twisted_affine_extension.py} checks the
matrices and the full affine braid equations independently.

For $B_5$, its central full twist acts on this twisted module by
$(-i)^{20}C^5=-I$, using D23.T2. For any cocycle $c$, comparison
of its value on $gz$ and $zg$, with this central $z$, gives
$2c(g)=(1-g)c(z)$, so it is again a coboundary. This gives an
independent conceptual obstruction in the five-strand case.
\end{proof}

\paragraph{Scope.}
This rules out an affine lift through the specified rank-four
linear representation. It does not rule out higher-dimensional,
non-affine, relative, or Brunnian representations, and it is not
a counterexample to the pentagon or to the desired comparison.

% DELTA23 /root/five_level_six. An attempted main-gate detector that collapses.

\subsection*{D23.F11. The new symplectic covariance has only the old shear}
Let $t_h\in\widehat{\mathbb Z}$ be the full signed
three-strand shear of DELTA\_ON\_GT AD.1 for
$h\in\mathcal H=\ker\lambda\subset\widehat{GT}_0$.
Use the chain form and the representation $q$ of D23.T1, and
write $w=e_1+e_3$. Define
\[
 V(t)=\begin{pmatrix}
 1&t&0&t\\0&1&0&0\\0&0&1&t\\0&0&0&1
 \end{pmatrix}.
 \tag{D23.F11.1}
\]
Even without assuming the pentagon, the four Drinfeld
generator formulas, evaluated in this symplectic representation,
are exactly conjugation by $V(t_h)$.
For an actual full-GT element, this is the $V_h$ of D23.T1.
In particular
\[
 V_he_1=e_1,\quad V_he_3=e_3,\quad V_hw=w.
\]
The putative sign on the line $\widehat{\mathbb Z}w$ is
identically positive; it cannot detect $\Delta$.

\begin{proof}
The integral lattice is the orthogonal sum of the symplectic
planes
\[
 E=\langle e_1,e_2\rangle,\qquad
 E^\perp=\langle w,e_4\rangle,
\]
with each displayed ordered pair having pairing one.  The
matrix $V(t)$ is $\operatorname{diag}(U^t,U^t)$ in this
integral symplectic basis, so it preserves the chain form.

Put $t_i=q(\sigma_i)$ and
$z_i=(\sigma_1\cdots\sigma_{i-1})^i$.
The exact Drinfeld formula can be written
\[
 \sigma_1\mapsto\sigma_1,\qquad
 \sigma_i\mapsto f_h(z_i,\sigma_i^2)^{-1}
                   \sigma_i f_h(z_i,\sigma_i^2)\quad(i\geq2).
 \tag{D23.F11.2}
\]
See Nakamura, RIMS-1691, Section 5.4, equation (5.4.2),
printed p.~28.  This is equivalent to the usual formula
with $y_i=\sigma_{i-1}\cdots\sigma_1^2\cdots\sigma_{i-1}$:
$z_i=y_i z_{i-1}$, and $z_{i-1}$ commutes with both $y_i$
and $\sigma_i^2$.  Multiplying the first argument of a
commutator word by such a commuting central factor has no
effect, because its exponent sum in that argument is zero.

For $i=1$, $V(t)$ fixes $e_1$ and hence its transvection.
For $i=2$, the two arguments act on $E$ as $U^2,L^2$ and
fix $E^\perp$ pointwise.  AD.1 says precisely that the
resulting image of $t_2$ is its conjugate by $U^{t_h}$ on
$E$.  Thus it is conjugation by $V(t_h)$ on the full lattice.

For $i=3$, direct multiplication gives
\[
 q(z_3)=
 \begin{pmatrix}-1&0&2&0\\0&-1&0&0\\0&0&1&0\\0&0&0&1\end{pmatrix}.
\]
It preserves the plane $\langle e_1,w\rangle$, while
$t_3^2$ is identity on that plane.  Their restricted actions
commute.  Therefore the image of every commutator word,
in particular $f_h$, is identity on that plane and fixes
$e_3=w-e_1$.  The evaluated conjugate of $t_3$ is exactly
$t_3$, which is also its conjugate by $V(t_h)$.

For $i=4$, the integral identity
\[
 q(z_4)=(t_1t_2t_3)^4=t_w^2
\]
holds, where $t_w=I+ww^{\mathsf t}J$.
Both $t_w^2$ and $t_4^2$ act on $E^\perp$ as exactly
$U^2,L^2$, in the basis $(w,e_4)$, and fix $E$ pointwise.
The identical AD.1 calculation therefore gives the shear
$U^{t_h}$ on this second plane.  This is again conjugation
by $V(t_h)$.
All four formulas have now been checked.  For a full-GT
element they give the actual braid action, so uniqueness
in D23.T1 identifies $V_h$ with (D23.F11.1).  Its action
on $w$ is visibly identity.
\end{proof}

\paragraph{Scope of the failed detector.}
All the evaluated braid relations and this line sign are
automatic for gentle elements, including the extra-CM
element.  This does not assert that every separate lift
of the five-cycle word has trivial symplectic evaluation.
The representation does not descend to the five-point
sphere group, and any such lifted-pentagon assertion
requires its own calculation with its central factors retained.

\subsection{A five-point Krammer quotient retaining the cubic detector}

\paragraph{D23.P6 (a concrete extension outside the low-derived tests).}
Let $k=\mathbb F_3[i]$, $i^2=-1$. In the ordered basis
$(F_{12},F_{13},F_{14},F_{23},F_{24},F_{34})$, put
\[
A=\begin{pmatrix}
-i&0&0&-1-i&-1-i&0\\
0&0&0&i&0&0\\
0&0&0&0&i&0\\
0&1&0&1-i&0&0\\
0&0&1&0&1-i&0\\
0&0&0&0&0&1
\end{pmatrix},\quad
B=\begin{pmatrix}
0&i&0&0&0&0\\
1&1-i&0&0&0&0\\
0&0&1&0&0&0\\
0&1-i&0&-i&0&-1-i\\
0&0&0&0&0&i\\
0&0&0&0&1&1-i
\end{pmatrix},
\]
\[
C=\begin{pmatrix}
1&0&0&0&0&0\\
0&0&i&0&0&0\\
0&1&1-i&0&0&0\\
0&0&0&0&i&0\\
0&0&0&1&1-i&0\\
0&0&1-i&0&1-i&-i
\end{pmatrix}.
\tag{D23.P6.1}
\]
The assignment $(\sigma_1,\sigma_2,\sigma_3)\mapsto(A,B,C)$
is a continuous finite representation of $\widehat B_4$ and
$(ABC)^4=-I$. It therefore defines an actual projective quotient
of $L=B_4/\langle z_4\rangle$, the label-five stabilizer in
$\Gamma_{0,[5]}$. Induction from that subgroup gives a finite
quotient of the full five-point group if desired.

The subspace $W=\langle F_{12},F_{13},F_{23}\rangle$ is
invariant under $A,B$. Its representation is exactly the Frobenius
contragredient of the affine detector D23.P3. More precisely, with
\[
P=\begin{pmatrix}
1-i&1-i&-1\\
-1-i&1-i&i\\
-1&-1-i&1
\end{pmatrix},
\quad
P^{-1}=\begin{pmatrix}
0&-1&i\\
-i&-1&0\\
1-i&1-i&1+i
\end{pmatrix},
\tag{D23.P6.2}
\]
and $\overline{i}=-i$, one has
\[
 A|_W\,P=P\,\overline{\eta(\sigma_1)}^{-\mathsf T},\qquad
 B|_W\,P=P\,\overline{\eta(\sigma_2)}^{-\mathsf T}.
\tag{D23.P6.3}
\]
Here affine matrices are written homogeneously in dimension three,
$\sigma_1=(sr)^{-1}$ and $\sigma_2=s\sigma_1s$.
Consequently $f(A^2,B^2)|_W$ determines the cubic sign $a_3(h)$
for every actual gentle $h=(1,f)$ by D23.P3. This is an actual
five-point quotient retaining that sign; it is not merely a
four-point forgetful product.
The same recovery works from the projective evaluation: the
homogeneous affine image $E_F$ contains no nonidentity scalar
matrix, so its affine normalization selects the unique lift on $W$.

The generated first-pair image
$P_{12}=\langle A^2,B^2\rangle$ has order $46656$, with derived
series orders
\[
46656,\quad3888,\quad486,\quad243,\quad3,\quad1.
\tag{D23.P6.4}
\]
In particular $P_{12}''$ is not nilpotent. Its quotient on $W$
is D23.P3's $E_F$, whose second derived group is
$k^2\rtimes\langle-I\rangle$, a nonnilpotent group of order $162$.
Thus the nilpotent-second-derived exclusion criterion of supplied
FP.36 does not dispose of this representation.

\begin{proof}
The matrices are the Krammer matrices at $q=i,t=-i$ in the
corrected six-case formula of Bigelow, \emph{The Lawrence--Krammer
representation}, arXiv:math/0204057v1, section4.2, printed page13.
No assertion identifying this specialized lattice with the
homological Lawrence lattice is used. Alternatively, the three
braid/distant relations and $(ABC)^4=-I$ follow by direct
multiplication of (D23.P6.1); the exact script
\path{agents/pentagon/krammer_f9_probe.py} verifies them.
All matrices are invertible, as follows either from the formula
or from their explicitly invertible diagonal blocks after the
corresponding basis permutation. Hence they define the asserted
finite representation, which extends continuously to the profinite
completion. Killing the scalar full twist is precisely the
projective passage needed for $L$; no spherical relation is
assumed for an unsupported direct $B_5$ representation.

Inspection of the first, second and fourth columns shows the
invariance of $W$. Its two matrices are
\[
K_1=\begin{pmatrix}-i&0&-1-i\\0&0&i\\0&1&1-i\end{pmatrix},
\qquad
K_2=\begin{pmatrix}0&i&0\\1&1-i&0\\0&1-i&-i\end{pmatrix}.
\]
Direct multiplication verifies both $PP^{-1}=I$ and
(D23.P6.3). In encoded $\mathbb F_9$ arithmetic these are
checked independently in
\path{agents/pentagon/krammer_affine_conjugacy.py}.
Since Frobenius followed by contragredience is an invertible
operation on matrix groups, these identities identify the entire
marked $B_3$ representation on $W$, rather than merely its order
or composition factors. Thus its pure image is exactly the
Frobenius contragredient of $E_F$, and its value on $f$ recovers
$\eta(f)$ by
\[
 \eta(f)=\overline{\bigl(P^{-1}f(A^2,B^2)|_W P\bigr)^{-\mathsf T}}.
\]
The cubic-sign formula now follows from D23.P3.

For completeness, the finite-order statements are certified in two
ways. The script \path{agents/pentagon/krammer_f9_probe.py}
uses a faithful permutation action on a union of vector orbits
containing a basis, and computes its stabilizer chain and derived
series. The independent matrix enumeration
\path{agents/pentagon/krammer_first_pair.py} enumerates exactly
$46656$ distinct matrices closed under the two generators and
exactly $3888$ matrices generated by the normal closure of their
commutator. The latter commutator has $486$ conjugates and three
of them generate the derived subgroup. The associated finite data
are retained in the JSON/NPZ certificates. These checks agree.

Finally, $E_F'=k^2\rtimes Q_8$ and
$E_F''=k^2\rtimes C_2$, since $-I$ acts as $-1$ on $k^2$.
The latter group is nonnilpotent: its Sylow two subgroup does not
commute with its normal Sylow three subgroup. A quotient of a
nilpotent finite group is nilpotent. The surjection
$P_{12}''\twoheadrightarrow E_F''$ therefore proves the final
assertion independently of the complete derived-series counts.
\end{proof}

\paragraph{Current scope.}
D23.P6 constructs and checks a new finite five-point test. It does
not prove the pentagon in this quotient for every gentle element,
and it does not prove $\Delta|_{\widehat{GT}}=0$ or exhibit a
full-GT counterexample. The exact first-pair evaluations are further
reduced to six matrices in D23.F13. The correlations with the other
two pentagon pair types remain the substantive calculation.

\subsection{The Krammer quotient and the disappearing central phase}

\paragraph{D23.F13 (six necessary first-pair values).}
Let $k=\mathbf F_3[i]$, $i^2=-1$, and use the specialized
Lawrence--Krammer matrices $K_1,K_2,K_3$ of D23.P6, in the
ordered basis $(12,13,14,23,24,34)$, with $q=i,t=-i$.
Set $X=K_1^2$, $Y=K_2^2$, $G=\langle X,Y\rangle$, and
\[
 W=\langle12,13,23\rangle,\qquad Q=k^6/W.
\]
The following statements hold.
\begin{enumerate}
\item The image of $G$ on $Q$ is dihedral of order eight.
For every actual gentle multiplier-one pair $(1,f)$,
$f(X,Y)|_Q=I$.
\item The map from $G$ to its two diagonal blocks has image
$E_F\times D_8$ of order $15552$ and kernel of order three.
Here $E_F$ is the cubic affine image of D23.P3, transported
to $W$ by the Frobenius contragredient comparison in D23.P6.
The kernel is central even in $\langle K_1,K_2\rangle$ and is
generated by
\[
 n=((K_1K_2)^3)^4=I+N,
 \qquad
 N_{W,Q}=\begin{pmatrix}
 -1+i&-1+i&-1+i\\
 1+i&1+i&1+i\\
 1-i&1-i&1-i
 \end{pmatrix},
\tag{D23.F13.1}
\]
all other blocks of $N$ being zero.
\item The actual first-pair matrix $f(X,Y)$ is uniquely
determined by $(\rho_h,a_3(h))\in\mathbf F_3\times\mathbf F_2$.
There are six explicit necessary matrices, one for each such
pair, in \texttt{krammer\_quotient\_extension.json} and
\texttt{krammer\_six\_first\_pair.npz}.  This assertion does
not assert a profinite lift for an arbitrary finite candidate.
In particular neither the quotient block nor the central
order-three factor supplies an additional independent state
of the first pair.
\end{enumerate}

\begin{proof}
The quotient matrices on $(14,24,34)$ are
\[
 A=\begin{pmatrix}0&i&0\\1&1-i&0\\0&0&1\end{pmatrix},
 \qquad
 B=\begin{pmatrix}1&0&0\\0&0&i\\0&1&1-i\end{pmatrix}.
\]
Direct multiplication gives $(A^2)^2=(B^2)^2=I$ and
$A^2B^2$ of order four.  Two involutions whose product has
order four generate the dihedral group of order eight:
each word reduces to $(A^2B^2)^j$ or $(A^2B^2)^jA^2$,
and these eight words are distinct.  Its commutator subgroup
is central of order two.

For completeness, the consequence of the gentle relation
uses no congruence-covariance assumption.  In a two-generator
group of nilpotency class at most two, a word from the closed
commutator subgroup of the free profinite group evaluates
as $[x,y]^c$ for its commutator coefficient
$c\in\widehat{\mathbf Z}$ modulo the order of $[x,y]$.
Writing $z=(xy)^{-1}$, the class-two identities give
$[y,z]=[z,x]=[x,y]$.  The multiplier-one three-cycle relation
therefore gives $[x,y]^{3c}=1$.  In $D_8$ the commutator has
order two, and hence $[x,y]^c=1$.  This proves the first
assertion for every actual gentle $f$.

By D23.P3 and the explicit Frobenius contragredient
intertwiner, the image on $W$ is
$E_F=k^2\rtimes\operatorname{SL}_2(\mathbf F_3)$, of order
$1944$.  Its abelianization is $C_3$: the central element
$-I$ makes the translation subgroup a commutator subgroup,
and $\operatorname{SL}_2(\mathbf F_3)^{\rm ab}=C_3$.
A nontrivial finite $2$-group quotient of $E_F$ would have
a further quotient $C_2$, which this abelianization excludes.
Thus $E_F$ and $D_8$ have no common nontrivial quotient.
Goursat's lemma applied to the subdirect diagonal image
gives the full product, of order $1944\cdot8=15552$.

Here is the exact finite certificate used for the remaining
kernel count; its scope is finite matrix arithmetic, not a
formal proof-assistant verification.  The program
\texttt{krammer\_quotient\_extension.py} starts from the two
literal six-dimensional matrices over $k$, represents
$a+bi$ by the block
$\left(\begin{smallmatrix}a&-b\\b&a\end{smallmatrix}\right)$
over $\mathbf F_3$, and enumerates $\langle X,Y\rangle$.
It starts with the identity and repeatedly adjoins right
products by $X,Y$ until no new matrix appears; full exact
matrix byte strings are used as keys, with equality tested
by equality of their entries.  Every enumerated element is
a word in $X,Y$, and at termination the finite set is closed
under right multiplication by these generators.  Since
the generators are invertible, the set is exactly $G$.
This exhaustive enumeration has $46656$ distinct matrices.
It is rebuilt independently of the matrix array supplied
by D23.P6, and the two sets are compared for equality.
The program checks the diagonal block image has $15552$
distinct values, and that exactly three matrices have both
blocks the identity.  Thus the kernel has order three.
The explicit product $((K_1K_2)^3)^4$ is exactly (D23.F13.1).
It is nonidentity, $N^2=0$, and $(I+N)^3=I$; it consequently
generates the kernel.  The braid relation makes
$(K_1K_2)^3$ central in $\langle K_1,K_2\rangle$, proving
the asserted centrality without relying on the enumeration.

It remains to remove the apparent choice of a central phase.
Put $S=K_1K_2K_1$.  The braid relation gives
$SXS^{-1}=Y$ and $SYS^{-1}=X$.  The two-cycle relation on
an actual $f$ is therefore the literal matrix equation
\[
 S f(X,Y)S^{-1}=f(X,Y)^{-1}.
\tag{D23.F13.2}
\]
If two possible matrices $F,F'$ have the same $W$ block
and trivial $Q$ block, then $F'=Fn^j$, $j\in\mathbf F_3$.
Since $n$ is central, substituting into (D23.F13.2) gives
$n^j=n^{-j}$.  As $n$ has order three, this forces $j=0$.
Thus the $W$ block uniquely determines an actual value.

For an explicit description, D23.P3 gives the six affine
matrices
\[
 \eta(f)=\begin{pmatrix}
 1&\rho&d_1\\
 \rho&1+\rho^2&d_2\\
 0&0&1
 \end{pmatrix},\quad
 d_1=\frac{1-\rho-\epsilon}{1+i},\quad
 d_2=(\rho+i)d_1,\quad \epsilon=(-1)^{a_3}.
\]
Their $W$ matrices are
$P\,\overline{\eta(f)}^{-T}P^{-1}$, for
\[
 P=\begin{pmatrix}1-i&1-i&-1\\-1-i&1-i&i\\-1&-1-i&1\end{pmatrix}.
\]
The certificate constructs these six blocks directly from
the displayed formulas, searches the enumerated group for
matrices with each block and trivial quotient block, and
finds three in each fiber.  It then tests (D23.F13.2) and
finds precisely one in each fiber.  These are the six
matrices in the stated files.  All restrictions used were
necessary identities of actual gentle elements.  No finite
matrix solution is being promoted to an actual profinite
pair, and no conclusion about $a_2=a_3$ follows from this
first-pair computation alone.
\end{proof}

\paragraph{Source and marking.}
The polynomial matrix convention is Stephen Bigelow,
\emph{The Lawrence--Krammer representation}, arXiv:math/0204057,
Section 4.2, as used and checked in D23.P6.  The passage
from the $W$ block to the affine cubic detector is the
explicit intertwiner in
\texttt{krammer\_affine\_conjugacy.py}; it is Frobenius
contragredient, not an ordinary conjugacy of the affine
homogeneous matrices.

% DELTA23 independent mixed Krammer analysis, /root/characters.
% Matrix certificates use F9=F3[i], i^2=-1, integer encoding a+3b.

\subsection{The two mixed Krammer pairs and their exact gentle correlation}

Use the six-dimensional Krammer representation of $B_4$ at
$q=i,t=-i$ over $k=\mathbb F_9$.  Its three literal matrices
$K_1,K_2,K_3$ are constructed in the companion pentagon workstream
from Bigelow's displayed Krammer formula.  Set
\[
 X_1=K_1^2,\quad X_2=K_2^2,\quad X_3=K_3^2,
 \quad X_4=(K_1K_2K_1)^2,\quad X_5=(K_3K_2K_3)^2.
\]
The exact matrices, all the bases below, and the full six-dimensional
intertwiners are recorded in
\path{agents/characters/records/krammer_mixed_correlation.json}; the executable
\path{agents/characters/code/check_krammer_mixed_correlation.py} verifies every
displayed matrix identity in the literal original basis.

In the matrix arrays in this subsection, an integer from zero to
eight means $a+bi\in\mathbb F_9$ when its integer value is $a+3b$,
with $0\le a,b\le2$.  Thus $3=i$, $6=-i$, and $2=-1$.

\paragraph{D23.C9 (proved by exact finite linear algebra and group certificates).}
Each mixed pair $(X_3,X_4)$ and $(X_4,X_5)$ preserves a direct
sum of two three-dimensional spaces.  On the first its image is
a solvable group $H$ of order $108$, with derived orders
\[
 108,27,3,1.
\]
On the second its image is a full unitary group $U(3,3)$ of order
$24192$, with perfect derived group $SU(3,3)$ of order $6048$.
The orders and successive derived orders of the full mixed images are
\[
\begin{array}{c|rrrr}
 \text{pair}&G&G'&G''&G'''=G''''\\\hline
 (X_3,X_4)&1306368&163296&18144&6048\\
 (X_4,X_5)&653184&163296&18144&6048.
\end{array}
\]
Their full linear images inject into the projective images obtained
by quotienting by the scalar $-I$.

There are explicit $C,Q,D\in\mathrm{GL}_6(k)$ such that, writing
$z=(X_4X_5)^{-1}$,
\[
 X_3C=CDz,\qquad X_4C=CX_4,\qquad
 [D,X_4]=[D,X_5]=1,
\]
\[
 zQ=Qz,\qquad X_4Q=QX_5.
\]
Consequently, for every actual multiplier-one gentle word $f$, if
\[
 R=f(X_3,X_4),\qquad r=C^{-1}RC,
\]
then its two mixed values obey the literal six-dimensional identity
\[
 \boxed{f(X_4,X_5)=Q^{-1}rQr^{-1}.}
\]
In particular the second mixed value is not an independent variable
in a finite necessary pentagon test.  The first mixed value belongs
to a group of order $163296$.

\emph{Proof.}
Let $B_{34}$ and $B_{45}$ be the splitting matrices from the
certificate, with the solvable summand first.  Their columns are
the indicated bases of the two invariant spaces, and have rank six.
Writing the corresponding three-dimensional generator pairs as
$(x_H,y_H),(x_U,y_U)$ and $(u_H,v_H),(u_U,v_U)$, the certificate
checks
\[
 B_{34}^{-1}(X_3,X_4)B_{34}
    =((x_H\oplus x_U),(y_H\oplus y_U)),
\]
\[
 B_{45}^{-1}(X_4,X_5)B_{45}
    =((u_H\oplus u_U),(v_H\oplus v_U)).
\]
The invariant-space discovery enumerated all $101$ projective
eigenvectors of one generator in each pair.  Discovery is not
needed as an assumption: the displayed bases and intertwining
identities independently certify the direct sums.

The two solvable block groups were computed as permutation groups
on unions of vector orbits containing a full $\mathbb F_3$-basis
of $k^3$.  Such actions are faithful, since a matrix fixing those
vectors is the identity.  Exact stabilizer chains give orders
$108,27,3,1$ in the derived series.  The first derived group has
exponent three.  The determinant maps the whole group onto
$\mu_4\subset k^\times$, so its kernel, of order $27$, is exactly
the derived group.

For the unitary blocks, the two respective invariant Hermitian
forms are
\[
 J_{34}=\begin{pmatrix}0&1&5\\1&0&3\\8&6&0\end{pmatrix},
 \qquad
 J_{45}=\begin{pmatrix}0&3&6\\6&0&4\\3&7&0\end{pmatrix}.
\]
They are Hermitian for $i\mapsto-i$, have nonzero determinants
$2$ and $1$, and the exact identities $g^*Jg=J$ hold for both
generators.  A faithful union of vector orbits has $224$ vectors;
the stabilizer-chain calculation gives order $24192$ and derived
order $6048$, with the derived group perfect.  The full unitary
group has order
\[
 3^3(3+1)(3^2-1)(3^3+1)=24192.
\]
For a direct count at this small field, there are $252$ vectors of
norm one in dimension three, $24$ in their orthogonal complement
of dimension two, and $4$ in dimension one, giving the same product.
Thus containment and the computed order identify the full group.
The determinant image has order four and the derived group lies
in its kernel; the kernel has order $6048$, so it is exactly the
special unitary group.  This identification does not need a
classification or simplicity theorem for finite groups.

For completeness, the full group orders do not require enumerating
the large six-dimensional image.  Each full mixed image is a
subdirect product of its solvable block $H$ and unitary block $U$.
Any common quotient of $H$ and $U$ is solvable, hence kills the
perfect group $U'=SU(3,3)$ and is abelian.  It therefore also kills
$H'$.  Goursat's lemma gives $H'\times U'$ inside the mixed image;
the image is the full preimage of its determinant-pair image.
For pair $34$, the generator determinant pairs are
\[
 (-1,1),\quad(-i,-i),
\]
generating the equal-parity subgroup of $\mu_4^2$ of order eight.
For pair $45$ they are both $(-i,-i)$ and generate the diagonal
subgroup of order four.  Multiplication by
$|H'|\,|U'|=27\cdot6048$ gives the two stated group orders.
Applying the same solvable/perfect Goursat argument to the derived
subgroup gives $G'=H'\times U'$, and then the remaining orders
follow by taking derived subgroups of a direct product.  Each
solvable block has a common fixed line.  Thus the full mixed image
does not contain scalar $-I$, proving the assertion about the
projective quotient.

Here are small matrices constructing the required full intertwiners:
\[
 C_H=\begin{pmatrix}4&1&0\\3&0&0\\0&0&1\end{pmatrix},
 \quad C_U=\begin{pmatrix}1&3&2\\3&5&0\\0&1&0\end{pmatrix},
\]
\[
 Q_H=\begin{pmatrix}4&0&6\\2&1&0\\1&0&0\end{pmatrix},
 \quad Q_U=\begin{pmatrix}2&0&2\\2&5&0\\6&2&1\end{pmatrix}.
\]
Put $z_H=(u_Hv_H)^{-1}$ and $z_U=(u_Uv_U)^{-1}$.  Direct
multiplication gives
\[
 x_HC_H=C_Hz_H^3,\quad y_HC_H=C_Hu_H,
 \quad x_UC_U=C_U(-z_U),\quad y_UC_U=C_Uu_U,
\]
and $z_H^2$ commutes with $u_H,v_H$.  For each block one also has
$z_jQ_j=Q_jz_j$ and $u_jQ_j=Q_jv_j$.  Define
\[
 \begin{split}
 C&=B_{34}(C_H\oplus C_U)B_{45}^{-1},\\
 Q&=B_{45}(Q_H\oplus Q_U)B_{45}^{-1},\\
 D&=B_{45}(z_H^2\oplus(-I_3))B_{45}^{-1}.
 \end{split}
\]
All three matrices are invertible; the stated literal
six-dimensional identities follow by conjugating the block
identities.  The companion computation also checks them directly
in the original basis, independently of this deduction.

If a profinite word $f$ belongs to $F'$, multiplying an input by
an element centralizing both inputs does not change its value:
the extra factor is raised to the corresponding abelian exponent,
which is zero.  Consequently
\[
 C^{-1}f(X_3,X_4)C=f(Dz,X_4)=f(z,X_4)=r.
\]
The other intertwiner gives
$f(z,X_5)=Q^{-1}rQ$.  The multiplier-one two-cycle equation says
$f(X_5,z)=f(z,X_5)^{-1}=Q^{-1}r^{-1}Q$.  In the source's exact
order, its three-cycle equation is
\[
 f(z,X_4)f(X_5,z)f(X_4,X_5)=I
\]
(see DELTA\_ON\_GT P.9, the displayed two- and three-cycle
equations in the normalized lift proof).  Substitution and inversion
give $f(X_4,X_5)=Q^{-1}rQr^{-1}$, as asserted.  Finally
$f\in F'$ puts its first mixed value in $G_{34}'$, whose order
has already been computed.\hfill$\square$

\paragraph{D23.C10 (proved; the nilpotent residual does not erase this block).}
Under either mixed-pair evaluation, the image of $F''$ is
$H''\times SU(3,3)$, of order $3\cdot6048$.  The image of
$R_{\rm nil}(F'')$ is exactly
\[
 \{1\}\times SU(3,3).
\]
These assertions concern the images of all words in the indicated
subgroups.  They do not assert that actual gentle words fill those
images or satisfy the pentagon there.

\emph{Proof.}
Surjectivity of the free evaluation and D23.C9 give the claimed
image of $F''$ as the second derived subgroup.  The $H''$ factor
has order three, so it is a finite nilpotent quotient and is
annihilated by $R_{\rm nil}(F'')$.  Let $N$ be the image of that
residual in the perfect $SU(3,3)$ factor.  The finite group
$SU(3,3)/N$ is a quotient of the pronilpotent group
$F''/R_{\rm nil}(F'')$, hence is nilpotent.  It is also perfect,
being a quotient of a perfect group.  A finite perfect nilpotent
group is trivial, so $N=SU(3,3)$.\hfill$\square$

\subsection{Exact reduced pentagon assignments in the new quotient}

\paragraph{D23.P7 (a finite calculation with a precise scope).}
Use D23.P6's matrices and put
\[
 X_1=A^2,\quad X_2=B^2,\quad X_3=C^2,\quad
 X_4=(ABA)^2,\quad X_5=(CBC)^2,\qquad J=ABCABA.
\]
Let $F_{\rho,e}$, $\rho\in\mathbb F_3,e\in\mathbb F_2$, be
the six exact first-pair matrices of D23.F13. Let $C_0,Q_0$ be
the full six-dimensional intertwiners of D23.C9. In the finite
derived subgroup $G_{34}'=\langle X_3,X_4\rangle'$, of order
$163296$, consider the equation
\[
 R(JR^{-1}J^{-1})(JF_{\rho,e}^{-1}J^{-1})
 \bigl(Q_0^{-1}rQ_0r^{-1}\bigr)F_{\rho,e}=I,
 \qquad r=C_0^{-1}RC_0.
\tag{D23.P7.1}
\]
For each of the six pairs $(\rho,e)$, exactly $36$ matrices $R$
satisfy this equation. There are no solutions with the right side
replaced by $-I$. Thus the reduced matrix equation admits both
cubic-sign values at every $\rho$.

Every actual multiplier-one full-GT element gives one of these
$216$ assignments, because D23.F13 fixes its first-pair value,
the two-cycle gives the two reversing-half-twist factors, and
D23.C9 gives the displayed fourth factor. This is a necessary
finite assignment theorem. The count by itself does not assert
that every listed assignment comes from a common word, an actual
gentle element, or an actual full-GT element.

\begin{proof}
The reversing half twist satisfies $J^2=-I$ and interchanges
$X_1,X_3$ and $X_4,X_5$ by conjugation while fixing $X_2$.
These are direct matrix identities. Therefore the two-cycle
expresses the five-cycle exactly as in (D23.P7.1), using the
actual inverse $J^{-1}$ rather than silently treating the lift
$J$ as an involution.

The group decomposition D23.C10 gives
\[
 G_{34}'=B_{34}\bigl(H'\times U'\bigr)B_{34}^{-1},
 \qquad |H'|=27,\quad |U'|=6048,
\]
where the two three-dimensional blocks and the splitting matrix
$B_{34}$ are literal matrices in D23.C9's certificate. The
enumeration program reconstructs each of $H',U'$ as the normal
closure of the commutator of its two displayed generators. It
then runs through all $27\cdot6048$ block-diagonal pairs, without
a sampling assumption.

For efficiency the program uses the following equivalent form of
(D23.P7.1). Put
\[
 A_R=RJR^{-1},\qquad
 B_R=rQ_0^{-1}r^{-1}Q_0J.
\]
Multiplying (D23.P7.1) on the indicated sides yields
\[
 F_{\rho,e}A_R=B_RF_{\rho,e}.
\tag{D23.P7.2}
\]
The program tests equality of all entries for all six $F_{\rho,e}$
and every enumerated $R$. It obtains $36$ each; testing the
negative equality gives zero each. Every retained assignment was
then checked again by multiplying the original five factors in
(D23.P7.1), rather than by reusing (D23.P7.2).

All arithmetic is exact over $\mathbb F_3$, after the faithful
restriction-of-scalars embedding of $\mathbb F_9$ matrices. The
reproducible certificate is
\path{agents/pentagon/krammer_pentagon_search.py}; its counts are
in the matching JSON file, and all $216$ matrices with their
$(\rho,e)$ labels are retained in
\path{agents/pentagon/krammer_pentagon_solutions.npz}.

The actual full-GT pentagon lies in the label-five stabilizer
$B_4/\langle z_4\rangle$. The displayed word is a product of
evaluations of $f\in\widehat F_2'$, so its braid exponent sum is
zero. If its lift to $\widehat B_4$ is $z_4^a$, abelianization
gives $12a=0$ in $\widehat{\mathbb Z}$, hence $a=0$. Thus the
literal identity in (D23.P7.1), not merely an arbitrary scalar,
is necessary. In this finite calculation the only additional
possible scalar lift would in any event be $-I$, and it yields
no solutions.
\end{proof}

\paragraph{Unresolved comparison.}
D23.P7 is stronger than discarding the new representation by a
nilpotence criterion: its first-pair cubic sign enters a genuinely
new pentagon calculation. Its outcome nevertheless leaves both
signs. The common-word correlations and the distinction between
finite assignments and actual profinite lifts must still be
respected. In particular no inference $e=\nu_4$ is contained in
the count $36$.

\subsection{Common-word realization on the three independent pair types}

\paragraph{D23.P8 (the three-pair commutator image is a direct product).}
For the matrices of D23.P6, the simultaneous evaluation map
\[
 \widehat F'\longrightarrow
 P_{12}'\times G_{34}'\times G_{45}'
\tag{D23.P8.1}
\]
is surjective. Here the factors are the derived groups of the
three corresponding pair images. In particular each of the $216$
reduced assignments in D23.P7 has a common commutator-word
preimage for these three pair values.

This statement does not yet prescribe the reversed mixed-pair
value, and it does not assert that the common word satisfies
the profinite gentle identities. These qualifications matter:
the reductions in D23.P7 used those identities for actual gentle
elements, whereas surjectivity in (D23.P8.1) concerns unrestricted
commutator words.

\begin{proof}
Write $G_{34}'=H_{34}'\times S_{34}$ and
$G_{45}'=H_{45}'\times S_{45}$ in the exact splitting bases of
D23.C9, where $|H_{34}'|=|H_{45}'|=27$ and
$S_{34}\simeq S_{45}\simeq\operatorname{SU}_3(3)$.

First, the two solvable full pair blocks $H_{34},H_{45}$ have
abelianization $C_4$ via determinant. With $-i$ as generator,
their two marked determinant phase pairs are $(2,1)$ for $x$
and $(1,1)$ for $y$. In any common abelian quotient of the two
blocks, if $a,b$ denote the images of their respective cyclic
abelianization generators, the common marking would impose
$2a=b$ and $a=b$. Hence $a=b=0$. If their simultaneous image
were a proper subdirect product, Goursat's lemma would give a
nontrivial common solvable quotient, whose abelianization would
be nontrivial. The displayed calculation excludes that.
Thus the full simultaneous image is $H_{34}\times H_{45}$,
and its derived image is $H_{34}'\times H_{45}'$.

Next the two simple factors have different kernels even on
$\widehat F'$. The following ordinary commutator word is a
small exact witness:
\[
 w=y^4x^2yx^{-1}y^{-1}x^{-1}y^{-4}xy^2x^2yxy^{-1}x^{-3}y^{-2}x^{-1}.
\tag{D23.P8.2}
\]
Its exponent sums are both zero. In the unitary block matrices
of D23.C9 it has values
\[
 w_{34}=I,\qquad
 w_{45}=\begin{pmatrix}5&7&4\\8&8&6\\5&7&5\end{pmatrix}\ne I,
\tag{D23.P8.3}
\]
using the same encoding $a+3b$ for $a+bi$ as the certificates.
The script \path{agents/pentagon/krammer_commonword_witness.py}
verifies the two values by literal word multiplication.
Since each factor is nonabelian simple, a subdirect product
of the two is either the whole product or the graph of an
isomorphism. The latter would give equal kernels, contrary
to (D23.P8.3). Thus the perfect joint image is
$S_{34}\times S_{45}$.

The standard simplicity of $\operatorname{SU}_3(3)$ is the
primary-source input used here. Its centre is trivial since
$\gcd(3,4)=1$; the GAP Character Table Library's
\emph{Simple Groups by Group Order}, row $U_3(3)$, lists its
order $6048$, and the corresponding group record identifies
it with $\operatorname{PSU}(3,3)$. This is the same cited
input used in D23.F14, not an inference from the order alone.

Finally $P_{12}'$ is solvable, and its abelianization has
order $3888/486=8$, by D23.P6. The group
$H_{34}'\times H_{45}'$ is a three-group. Any nontrivial
common quotient would be solvable and would have a nontrivial
abelian quotient, simultaneously a quotient of a two-group
and a three-group. This is impossible. Hence Goursat gives
the full product of these three solvable factors. It has no
nontrivial common quotient with the perfect product
$S_{34}\times S_{45}$, because a finite group both solvable
and perfect is trivial. A final application of Goursat proves
(D23.P8.1).

The image of the ordinary free commutator subgroup in any
finite target equals the image of its profinite closure.
Therefore ordinary commutator-word preimages can be chosen.
This proves the stated finite-word realization and its scope.
\end{proof}

\paragraph{Primary-source locators.}
\url{https://www.math.rwth-aachen.de/homes/Thomas.Breuer/ctbllib/ctbltoc/views/simplebyorder.html},
row $U_3(3)$;
\url{https://www.math.rwth-aachen.de/homes/sam/ctbllib/ctbltoc/data/U3(3).html},
``Group constructions in GAP'' entry $\operatorname{PSU}(3,3)$.

\subsection{The full common-word audit for the Krammer test}

\paragraph{D23.P9 (all reduced pentagon assignments have finite-word realizations).}
Retain D23.P6--P8's notation, and write
$G_{43}=\langle X_4,X_3\rangle$ with the indicated reversed
marking. The simultaneous map
\[
 \widehat F'\longrightarrow
 P_{12}'\times G_{34}'\times G_{45}'\times G_{43}'
\tag{D23.P9.1}
\]
is onto. Every one of the $216$ assignments of D23.P7 can
therefore be realized by an ordinary free commutator word whose
five actual Krammer evaluations multiply to $I$ in the displayed
pentagon order. No unrelated-word substitution is involved.

This is a statement about ordinary finite word images. It does
not prove that one of these words is a gentle parameter. In
particular it does not construct an actual full-GT element with
nonzero $\Delta$.

\begin{proof}
Use D23.C9's faithful decompositions of each of the three mixed
pair groups into a solvable block of order $108$ and a unitary
block of order $24192$. On the reversed pair use the same basis
as pair $34$ and exchange its two generators. Their derived
groups are respectively $H_j'$ of order $27$ and simple
$S_j=\operatorname{SU}_3(3)$ of order $6048$.

The exact finite calculation
\path{agents/pentagon/krammer_allfive_commonword.py} forms the
three block-diagonal solvable generator pairs and computes the
normal closure of their joint commutator. It contains exactly
$19683=27^3$ matrices. The conjugacy orbit of that commutator
has $432$ elements and six of them generate the entire closure.
The image is contained in the product of the three order-$27$
derived groups, so this count proves equality with that full
product. This is a complete finite enumeration, not a random
test. Normal closure of the commutator computes the derived
subgroup of a two-generator group because the resulting
quotient makes the two generators commute.

The simple factors are independent as well. D23.P8's word
$w_1$ of length $28$ has identity pattern
\[
  (w_1(S_{34}),w_1(S_{45}),w_1(S_{43})):
       \quad(I,\ne I,\ne I).
\]
A second explicit word is
\[
 w_2=yxy^2xyx^{-1}y^{-3}x^{-1}y^{-1}
       xyx^2yxy^{-1}x^{-3}y^{-1}x^{-1}.
\tag{D23.P9.2}
\]
It has exponent sums zero and satisfies
$w_2(S_{45})=I$, $w_2(S_{43})\ne I$.
The program verifies these literal word values in the stored
three-dimensional matrices after restriction of scalars to
$\mathbb F_3$. Hence the kernels on $F'$ of the three simple
projections are pairwise distinct. A subdirect product of
nonabelian simple groups is a product of diagonal factors;
distinct kernels exclude every identification. The joint image
is therefore $S_{34}\times S_{45}\times S_{43}$.
The standard simplicity input and primary locators are those
recorded explicitly in D23.P8 and D23.F14.

By D23.P6, the abelianization of the solvable group $P_{12}'$
has order eight. The joint solvable mixed image is a three-group.
They have no nontrivial common quotient, because any nontrivial
finite solvable quotient has a nontrivial abelian quotient.
Goursat therefore separates $P_{12}'$ from the three solvable
mixed factors. A further application of Goursat separates this
solvable product from the perfect product of the three simple
factors. This proves (D23.P9.1).

Now take a listed solution $(F,R)$ of D23.P7 and put
\[
 H=Q_0^{-1}(C_0^{-1}RC_0)Q_0(C_0^{-1}R^{-1}C_0).
\]
Surjectivity provides a common commutator word $f$ with
\[
 f(X_1,X_2)=F,\quad f(X_3,X_4)=R,\quad
 f(X_4,X_5)=H,\quad f(X_4,X_3)=R^{-1}.
\]
The first-pair half twist $S=ABA$ exchanges $X_1,X_2$.
Thus, for every word, not merely for gentle words,
$f(X_2,X_1)=Sf(X_1,X_2)S^{-1}$. The selected six first-pair
matrices satisfy $SFS^{-1}=F^{-1}$ by D23.F13. The reversing
half twist $J=ABCABA$ gives
\[
 f(X_2,X_3)=JF^{-1}J^{-1},\qquad
 f(X_5,X_1)=JR^{-1}J^{-1}.
\]
The actual five values of this one word are consequently the
five factors in D23.P7.1. Their product is $I$ by the checked
assignment. The finite image of the ordinary free commutator
subgroup equals that of its profinite closure, so $f$ can be
chosen as an ordinary word.

The argument prescribes only finitely many evaluations and their
displayed relations. It does not establish the identities in
the full profinite free group required for a gentle pair, which
is the stated limitation.
\end{proof}

\subsection{Complete outer three-puncture closure on the mixed perfect blocks}

\paragraph{D23.C15 (the mixed perfect closure retains all 216 candidates).}
In the unitary block of the pair $(X_4,X_5)$ from D23.C9 put
$a=X_{4,U}$, $b=X_{5,U}$, $c=(ab)^{-1}$. For a pair of these
matrices, its kernel below means the kernel of evaluation restricted
to $\widehat F'$. The six oriented pairs have exactly three kernel
classes:
\[
 \{(a,b),(b,a)\},\qquad
 \{(b,c),(a,c)\},\qquad
 \{(c,a),(c,b)\}.
\tag{D23.C15.1}
\]
Every coincidence is certified by an explicit linear intertwiner,
and every noncoincidence by a literal commutator word whose value
is identity in one pair and nonidentity in the other. The six pairs
are the full orbit modulo simultaneous inner conjugation under
permuting the three peripheral conjugacy classes.

All 216 necessary Krammer pentagon solutions of D23.P7 satisfy
all value correlations arising from (D23.C15.1). They also satisfy
the opposite-orientation three-cycle equation in the full original
six-dimensional representation. Thus this complete closure on the
mixed perfect blocks removes none of D23.P7's 36 candidates for
each $(\rho,a_3)$, including the cases $(0,1)$.
This is a finite necessary-equation result, not an assertion that
any candidate lifts to an actual gentle or full-GT element.

\begin{proof}
For $g\in U(3,3)$ put $g^\circ=(\det g)g$. The determinant is in
$\mu_4$, so $\det(g^\circ)=(\det g)^4=1$. Central multipliers of
inputs do not change evaluations of words in $F'$. We can therefore
use the pairs of determinant-one matrices without changing the
kernels or values under discussion. Their images are $SU(3,3)$:
each pair generates $U(3,3)$ before normalization, and its derived
image is already $SU(3,3)$ by D23.C9.

In the field encoding $a+3b\leftrightarrow a+bi$ used in D23.C9,
intertwiners for the three displayed coincidences, in the same
order, are
\[
 J_0=\begin{pmatrix}2&7&2\\0&3&0\\2&6&1\end{pmatrix},\quad
 J_1=\begin{pmatrix}2&8&2\\8&6&8\\5&4&1\end{pmatrix},\quad
 J_2=\begin{pmatrix}2&0&2\\2&5&0\\6&2&1\end{pmatrix}.
\tag{D23.C15.2}
\]
The certificate verifies $p_iJ=Jq_i$ for both generators of each
normalized source pair $(p_1,p_2)$ and target pair $(q_1,q_2)$.
Hence $f(q_1,q_2)=J^{-1}f(p_1,p_2)J$ for every $f\in F'$.
No semilinear automorphism or classification of automorphisms of
$SU(3,3)$ is used.

For all twelve pairs of distinct classes the certificate records a
word in the alphabet $\{1,2,-1,-2\}$, with zero exponent sum in
each generator, and checks its two evaluations exactly. These
words were obtained by a Cayley-graph collision and then taking
a commutator; the verification needs only the recorded words and
the displayed matrices. Thus they prove noncoincidence on $F'$
itself, rather than merely a failed search for linear intertwiners.
Together these checks prove (D23.C15.1).

Let $R=f(X_3,X_4)$ and $r=C^{-1}RC$ in the full six-dimensional
notation of D23.C9. The three oriented values on $(a,b,c)$ are the
unitary blocks of
\[
 A=Q^{-1}rQr^{-1},\qquad B=Q^{-1}r^{-1}Q,\qquad C_0=r.
\tag{D23.C15.3}
\]
They satisfy $C_0BA=1$ by construction. Two-cycle gives the three
reversed values $A^{-1},B^{-1},C_0^{-1}$. Restricting each of
D23.P7's exact matrix solutions to the pair-$45$ splitting basis,
the certificate checks all three intertwiner relations (D23.C15.2)
on these six values. Every solution passes. In particular the
self-duality of $A$ under the first reflection, which is not just
the original displayed cycle, is retained by every candidate.

For completeness the three-cycle at the reversed input $(b,a)$
has third element $c'=(ba)^{-1}=a^{-1}ca=bcb^{-1}$. Its exact
form after two-cycle and simultaneous conjugation is
\[
              bB^{-1}b^{-1}a^{-1}C_0^{-1}aA^{-1}=1.
\tag{D23.C15.4}
\]
The separate certificate checks this identity directly in the full
six-dimensional matrices for every one of the 216 candidates,
with no projective scalar discarded. Cyclic rotation supplies the
other cycles of this orientation; the original orientation is
already (D23.C15.3). The six representatives suffice because
swapping the first two generators changes the third to an inner
conjugate, as the displayed formulas for $c'$ show.

D23.C9 identifies the pair-$34$ perfect evaluation with $(c,a)$,
up to a central input scalar; its reverse is $(a,c)$. Hence the
same closure covers both mixed perfect blocks, not just pair $45$.
No additional statement about the solvable blocks is needed or
claimed here.

The executable records are
\path{agents/characters/code/krammer_nielsen_perfect.py} and
\path{agents/characters/code/check_krammer_reverse_cycle.py};
their JSON outputs record the three intertwiners, all twelve
commutator witnesses, all 216 value tests, and the six counts
$(0\text{ failures},36\text{ passes})$.
\end{proof}

\subsection{The two solvable mixed blocks have independent word images}

\paragraph{D23.F15.}
Let $(x_H,y_H)$ and $(u_H,v_H)$ be the solvable
three-dimensional block pairs of D23.C9 for pairs $34$
and $45$, respectively.  The simultaneous free-group
evaluation has image
\[
 \operatorname{im}\bigl(\widehat F\longrightarrow
       H_{34}\times H_{45}\bigr)=H_{34}\times H_{45},
\]
and hence
\[
 \operatorname{im}\bigl(\widehat F'\longrightarrow
       H_{34}'\times H_{45}'\bigr)=H_{34}'\times H_{45}',
 \qquad |H_{34}'\times H_{45}'|=729.
\tag{D23.F15.1}
\]
In particular the cyclic Nielsen automorphism
$\phi(x)=(xy)^{-1}$, $\phi(y)=x$ does not preserve
the kernel of the $45$-block evaluation on $\widehat F'$.
There is no internal conjugacy, even allowing central
factors in the two arguments, that makes this Nielsen
operation an automorphism of this marked derived image.

\begin{proof}
Both block groups are solvable of order $108$, and
their determinant maps identify their abelianizations
with $C_4$, by D23.C9.  Use $-i$ as generator of $C_4$.
The determinant exponents of the two simultaneous
generators are
\[
 x\longmapsto(2,1),\qquad y\longmapsto(1,1).
\]
Let $J\subset H_{34}\times H_{45}$ be the simultaneous
image. It is subdirect. By Goursat's lemma, a proper
such subgroup would give a nontrivial common quotient
$T$ of the two groups in which these two marked
generators agree.  The group $T$ is solvable and
nontrivial, so $T^{\rm ab}$ is nontrivial.  If $a,b$
are the images in this abelianization of the two
determinant generators, marked agreement gives
$2a=b$ and $a=b$.  Subtraction yields $a=b=0$.
Both maps onto $T^{\rm ab}$ would consequently be
trivial, a contradiction.  Thus $J$ is the full
direct product.  Images commute with taking the
closed derived subgroup under a finite surjection,
giving (D23.F15.1).

There is also an explicit short witness to failure
of kernel invariance.  The word
\[
 w=[x^2,y]\in\widehat F'
\]
is identity in the $34$ block because $x_H^2=I$.
In the $45$ block direct multiplication of the
matrices in D23.C9 gives
\[
 [u_H^2,v_H]=
 \begin{pmatrix}
 -i&0&-1+i\\
 0&1&1\\
 -1+i&0&-1+i
 \end{pmatrix}\ne I.
\tag{D23.F15.2}
\]
The central correction and the intertwiner $C_H$
from D23.C9 give, for every commutator word $f$,
\[
 C_H^{-1}f(x_H,y_H)C_H
       =f((u_Hv_H)^{-1},u_H).
\]
Consequently the $34$-block kernel on $\widehat F'$
is the inverse image under $\phi$ of the $45$-block
kernel.  The two kernels differ by (D23.F15.2), so
the latter is not $\phi$-invariant.  An internal
conjugacy of the marked pair, up to factors
centralizing both inputs, would imply exactly that
invariance on commutator words, which is impossible.
This does not discard the actual gentle correlation
of D23.C9: that correlation uses the two-cycle and
three-cycle identities, which arbitrary commutator
words need not satisfy.
\end{proof}

\subsection{The simultaneous Krammer word image cannot recover the harmonic phase}

\paragraph{D23.F14 (no $S_4$ quotient of the joint Krammer image).}
Use the same specialized matrices as D23.F13 and put
\[
 (X_1,X_2,X_3,X_4,X_5)
 =(K_1^2,K_2^2,K_3^2,(K_1K_2K_1)^2,(K_3K_2K_3)^2).
\]
Let $\Gamma$ be the image of the free profinite group
$\widehat F(x,y)$ under the simultaneous five evaluations
at the ordered pairs
\[
 (X_1,X_2),\quad(X_2,X_3),\quad(X_3,X_4),\quad
 (X_4,X_5),\quad(X_5,X_1).
\tag{D23.F14.1}
\]
Then $S_4$ is not a quotient of $\Gamma$.
Adjoin the marked evaluation $x\mapsto p=(123)$,
$y\mapsto t=(1234)$, and let
$\widetilde\Gamma\subset\Gamma\times S_4$ be the joint
image.  One has
\[
 \{1\}\times V_4\ \subset\ \widetilde\Gamma',
\tag{D23.F14.2}
\]
where $V_4$ is the group of double transpositions.
Consequently, for every commutator word $f$ whose marked
$S_4$ value is $t^{2\nu}p^\rho$, there is another
commutator word $f_*$ with exactly the same five Krammer
values and marked value $t^{2(\nu+1)}p^\rho$.
This is an exact statement about finite word images.
It does not assert that the altered word is an actual
gentle parameter or that it satisfies all profinite
two-cycle and three-cycle identities.

\begin{proof}
We first record exactly which previously computed matrix
facts enter the argument.  The first two pair images are
conjugate, up to exchanging generators, to the group $G$
of D23.F13.  The three remaining pair images have faithful
three-plus-three block decompositions into a solvable
group $H$ of order $108$, with $|H'|=27$, and the full
three-dimensional unitary group $U=\operatorname U_3(3)$
of order $24192$.  In each case the projection onto $U$
is surjective.  These are the literal block decompositions
and exact matrix group computations of D23.C9, in
\texttt{characters/records/krammer\_mixed\_correlation.json}
and the associated mixed-pair certificate. The proof
below needs only these stated facts, not an unverified
identification of a matrix order with a group structure.
The third mixed pair, not separately listed in D23.C9,
is accounted for by the reversing half twist
\[
 D_4=K_1K_2K_3K_1K_2K_1.
\]
The braid identities give
$D_4K_1D_4^{-1}=K_3$, $D_4K_2D_4^{-1}=K_2$ and
$D_4K_3D_4^{-1}=K_1$. Hence
$D_4X_3D_4^{-1}=X_1$ and
$D_4X_4D_4^{-1}=X_5$. Thus pair $(X_5,X_1)$ is
the conjugate by $D_4$ of the reversed pair $(X_4,X_3)$,
which has the same image group as pair $(X_3,X_4)$.
These literal matrix identities have also been checked
directly from the three stored Krammer matrices.
The same half twist proves the asserted relationship
between the first two pair groups.

By D23.F13, $G$ has a normal subgroup $N$ of order
$3^5=243$ whose quotient is
\[
 G/N\simeq\operatorname{SL}_2(\mathbf F_3)\times D_8.
\]
Indeed $N$ is the inverse image of the translation
subgroup of $E_F$ under the diagonal-block map, with
the $D_8$ coordinate constrained to be identity.
It is an extension of a group of order $3^4$ by the
central kernel of order three and is therefore a
three-group.  The group
$\operatorname{SL}_2(\mathbf F_3)$ has normal quaternion
Sylow two-subgroup $Q_8$ and quotient $C_3$.
Taking the inverse image of $Q_8\times D_8$ gives a
normal series for $G$ whose successive factors are,
in order, a three-group, a two-group and a three-group.
Each solvable block $H$ has the shorter series
$1\leq H'\leq H$, with factors a three-group and a
two-group; adjoining a central $C_4$ leaves this property
unchanged.

The determinant map on $U$ takes values in the cyclic
group of norm-one scalars $C_4$.  On scalar matrices
its value is $c^3$, which is an automorphism of $C_4$.
Thus scalar multiplication splits determinant and gives
the direct product
\[
 U=C_4\times S,\qquad S=\operatorname{SU}_3(3).
\]
Here $S$ is the nonabelian simple group of order $6048$;
its scalar centre is trivial because $\gcd(3,4)=1$.
For the standard simplicity fact, see W.~L.~Edge,
\emph{The simple group of order 6048}, Proc. Cambridge
Philos. Soc. \textbf{56} (1960), 189--204, introduction
on p.~189; see also the primary GAP Character Table
Library's \emph{Simple Groups by Group Order}, row
$U_3(3)$ of order $6048$, and its $\operatorname{PSU}(3,3)$
identification in the corresponding character-table record.

After these splittings, the simultaneous image embeds
in $A\times S^3$, where $A$ is a product of two copies
of $G$ and three copies of $H\times C_4$.
Write $A_0$ and $D$ for its projections onto $A$ and
$S^3$.  The first projection is solvable, and $D$ is
subdirect in the three nonabelian simple factors.
Such a subdirect product is a product of diagonal
copies of $S$, and in particular is perfect.  Here is
the short argument in the precise form used.  Induct
on the number of factors.  Goursat's lemma for a
subdirect subgroup of $S^r\times S$ either adds the
last factor independently or identifies it with a
simple quotient of $S^r$.  Every surjection from
$S^r$ to $S$ is projection onto one factor followed
by an automorphism: the factor images are commuting
normal subgroups of the simple centreless group $S$,
so exactly one image is nontrivial.  This proves the
claim by induction.

Apply Goursat again to the subdirect embedding
$\Gamma\subset A_0\times D$.  Any common quotient
would be both solvable, as a quotient of $A_0$, and
perfect, as a quotient of $D$.  A finite group with
both properties is trivial.  Therefore
\[
 \Gamma=A_0\times D.
\]
Every solvable quotient of $\Gamma$ consequently
factors through $A_0$.

Groups admitting a normal series with factors a
three-group, a two-group and a three-group are closed
under finite direct products, subgroups and quotients:
take respectively products, intersections or images
of the three terms of the series.  Hence $A_0$ and
all its quotients have such a series.  But $S_4$ does
not.  Its largest normal three-subgroup is trivial,
and its largest normal two-subgroup is $V_4$; the
quotient $S_4/V_4\simeq S_3$ is not a three-group.
No smaller normal two-subgroup can have three-power
index either.  It follows that $S_4$ is not a quotient
of $A_0$, proving the first assertion.

The map to $S_4$ in the statement is surjective, since
the displayed $p,t$ generate $S_4$.  Goursat's lemma
for $\widetilde\Gamma\subset\Gamma\times S_4$ says
that the quotient of $S_4$ by
\[
 N_4=\{v\in S_4:(1,v)\in\widetilde\Gamma\}
\]
is a quotient of $\Gamma$.  The already proved result
excludes $N_4=1$.  Every nontrivial normal subgroup
of $S_4$ contains $V_4$, so
$\{1\}\times V_4\subset\widetilde\Gamma$.
Moreover $[S_4,V_4]=V_4$: the action permutes the
three nonidentity elements transitively and has no
nontrivial quotient on which it acts trivially.
Taking commutators with lifts of $S_4$ proves
(D23.F14.2).

The image of $\widehat F'$ under a finite surjection
is exactly the derived subgroup of its image.
By (D23.F14.2) choose $u\in\widehat F'$ whose five
Krammer images are identity and whose $S_4$ image
is $t^2$.  Then $f_*=uf$ has unchanged Krammer
values and $S_4$ value
$t^2t^{2\nu}p^\rho=t^{2(\nu+1)}p^\rho$, as claimed.
All images are finite, so a word in the ordinary free
commutator subgroup can equally well be chosen.
The argument never supplies the other identities
needed for an actual gentle pair, which explains the
scope stated above.
\end{proof}

\paragraph{Primary-source URLs checked in DELTA23.}
Edge's original paper record:
\url{https://doi.org/10.1017/S0305004100034472}.
GAP Character Table Library, simple-group table:
\url{https://www.math.rwth-aachen.de/homes/Thomas.Breuer/ctbllib/ctbltoc/views/simplebyorder.html}.
The corresponding primary group record is
\url{https://www.math.rwth-aachen.de/homes/sam/ctbllib/ctbltoc/data/U3(3).html}.
The proof distinguishes the notation $U$ for the full
unitary group of order $24192$ from the ATLAS notation
$U_3(3)$ for its simple projective special factor of
order $6048$.

\subsection{The whole alternating phase is independent of the Krammer word data}

\paragraph{D23.F17 (strengthening D23.F14).}
Let $\Gamma$ be the simultaneous five-pair word image
of D23.F14. Then $S_3$ is not a quotient of $\Gamma$.
For the joint image $\widetilde\Gamma\subset
\Gamma\times S_4$ with the marked map
$x\mapsto(123)$, $y\mapsto(1234)$, one has
\[
 \{1\}\times A_4\subset\widetilde\Gamma'.
\tag{D23.F17.1}
\]
Thus all twelve alternating $S_4$ values can be
prescribed independently of any fixed value of a
commutator word in the five Krammer evaluations.
In particular both of the finite $S_4$ coordinates
occurring in FP.37 can vary at the word level.
No claim of independent variation on actual gentle
or full-GT elements is made.

\begin{proof}
We use the full simultaneous derived-image theorem
D23.P9.  After the reversing-half-twist identifications,
it identifies the image of $F'$ with
\[
 P'\times H_{34}'\times H_{45}'\times H_{43}'
       \times S^3,
\tag{D23.F17.2}
\]
where $P$ is the first-pair group, each $H_j$ is the
solvable order-$108$ block, and $S=SU_3(3)$ is
perfect.  This is a statement about the joint image,
not merely its separate surjective projections.

The abelianization of $P'$ has order $8$, by D23.P6.
Each $H_j'$ has order $27$, exponent three, and
commutator subgroup of order three.  Its abelianization
is therefore the two-dimensional vector space over
$\mathbf F_3$. The explicit conjugation actions of
the marked generators $x,y$ on these three spaces are
\[
 \begin{array}{c|cc}
 &x&y\\\hline
 H_{34}'/H_{34}''&-I&J\\
 H_{45}'/H_{45}''&J&J\\
 H_{43}'/H_{43}''&J&-I
 \end{array},\qquad
 J=\begin{pmatrix}0&-1\\1&0\end{pmatrix}.
\tag{D23.F17.3}
\]
Here $43$ means the reversed $34$ pair, so its row
is obtained by interchanging the marked generators.
To give a complete certificate for the two independent
rows, put $A=[x_H,y_H]$, $B=y_HAy_H^{-1}$ and
$Z=[A,B]$ in the literal block matrices of D23.C9.
The $27$ products $A^aB^bZ^c$, $a,b,c\in\mathbf F_3$,
are distinct and form the entire derived block.
The program \texttt{krammer\_solvable\_action.py}
checks this by exact matrix multiplication, checks
that $Z$ is central of order three, and expresses
the conjugates of $A,B$ by each generator in those
coordinates.  The resulting first two coordinates
are exactly the matrices in (D23.F17.3); all
underlying three-dimensional matrices are retained
in \texttt{krammer\_solvable\_action.json}.
In particular $J^2=-I$ and its characteristic
polynomial $T^2+1$ has no root in $\mathbf F_3$.

Suppose that $\pi:\Gamma\twoheadrightarrow S_3$
were surjective. Then its restriction to $\Gamma'$
would surject onto $S_3'=C_3$, and would be
$\Gamma$-equivariant for conjugation.  The factors
$P'$ and $S^3$ of (D23.F17.2) must be killed: the
former has two-primary abelianization, while the
latter is perfect. On each of the three remaining
factors the induced map is a linear functional
$\ell:\mathbf F_3^2\to\mathbf F_3$.
Conjugation by any element of $S_3$ on $C_3$ is
either $+1$ or $-1$.  Each row of (D23.F17.3) has
a generator acting by $J$. Equivariance would
therefore require $\ell J=\ell$ or
$\ell J=-\ell$. Since $J$ has no eigenvalue $1$
or $-1$ over $\mathbf F_3$, this forces $\ell=0$.
All factor restrictions vanish, contradicting
surjectivity onto $C_3$. Hence no such quotient
$S_3$ exists, even without retaining a marking.

In the Goursat argument of D23.F14, the normal
subgroup $N_4\triangleleft S_4$ now cannot be
$1$ or $V_4$: the corresponding common quotients
would be $S_4$ or $S_3$, both excluded. Thus
$N_4$ contains $A_4$. Furthermore
$[S_4,A_4]=A_4$. One can see this from
$A_4'=V_4$ and the fact that an odd permutation
acts by inversion on $A_4/V_4=C_3$.
Taking commutators with lifts of $S_4$ proves
(D23.F17.1).  Finally the image of $F'$ is exactly
the derived joint image, so multiplication by
a preimage of any $(1,a)$, $a\in A_4$, changes
the $S_4$ value arbitrarily while retaining all
five Krammer values.  This constructs finite
word images only and does not supply the full
gentle identities for the modified word.
\end{proof}

% DELTA23, /root/characters.  Source overlap is stated explicitly below.

\subsection{An induced harmonic quotient and the exact finite ambiguity}

\paragraph{D23.C11 (proved; an audit of the proposed induced $S_4$ quotient).}
Induce the marked permutation representation
\[
 F\longrightarrow S_4,\qquad x\longmapsto(123),\quad
 y\longmapsto(1234),
\]
from the index-six pure subgroup to $B=C_2\amalg C_3$.
The induced representation has degree $24$.  Its $B$ image has
order $41472$ and its $F$ image has derived orders
\[
 6912,1728,64,1.
\]
In particular the abelianization of the first derived image is
a three-group of order $27$.

Join this representation to the signed projective congruence quotient
$Q_{12}$, represented faithfully on its $48$ primitive row vectors
modulo global sign.  The joint $B$ image has order $995328$ and
its pure $F$ image has order $165888$.  The centralizer of the
image of $s$ within that pure group has order $192$ and contains
an element $b$ with positive signed residue
\[
 q_{12}(b)=\begin{pmatrix}9&4\\8&9\end{pmatrix}.
\]
This residue is $I$ modulo four and $S$ modulo three, so its
harmonic parity is odd.  The same $b$ belongs to the first derived
pure image, of order $13824$, whose $s$-centralizer has order $96$.
Thus the finite conjugacy equation for a pure $s$-frame in this
quotient has an odd-parity ambiguity even with zero abelianized
frame.  This is a finite ambiguity assertion, not an assertion
that the ambiguous frame choices lift to actual gentle elements.

\emph{Proof.}
The companion script \path{agents/characters/code/induced_s4_harmonic.py}
uses the literal words
$x=(sr)^{-2}$ and $y=sxs$.  Its six coset representatives are
$1,s,r,r^2,sr,sr^2$.  Every Schreier cocycle is reduced in the
free product and rewritten exactly as a word in $x,y$; the full
rewriting table and the two induced permutations are recorded in
\path{agents/characters/records/induced_s4_harmonic.json}.
The checks $s^2=r^3=1$ and the restrictions
$x|_{\text{identity block}}=(123)$,
$y|_{\text{identity block}}=(1234)$ verify the source marking.
Exact permutation-group stabilizer chains give the displayed orders
and derived subgroups.  The joint image uses the same right-action
convention on primitive row vectors, so its generators are the
same literal $s,r$ in both factors.
For comparison with FP.37's left-action convention, inversion of
the two four-point permutations gives $(p^{-1},t^{-1})$.
This pair is simultaneously conjugate to $(p,t)$ by $(13)$;
thus it has exactly the same marked free-group kernel.  This
convention conversion does not alter the induced quotient used here.

The displayed centralizer witness has induced permutation
\[
 (0\ 1)(2\ 3)(4\ 5)(6\ 7)(12\ 15)(13\ 14)
             (16\ 18)(17\ 19),
\]
with the other points fixed, using zero-based labels.
Its full $72$-point permutation is in the certificate.  Exact
membership in the derived pure group and commutation with the
joint image of $s$ are checked.  The matrix calculation modulo
four and three gives the asserted parity.\hfill$\square$

The Ferrari extension of this quotient contributes no additional
pentagon equation by itself.  Under $a=c=u,b=v$ one has
$(X_1,X_2,X_3,X_4,X_5)=(x,y,x,1,1)$; the two surviving terms
are $f(y,x)f(x,y)=1$.  This is precisely the general forgetful
obstruction already proved in source FP.30, not a new result.

\subsection{A graph proof of the existing relative homology detector}

\paragraph{D23.C12 (proved; independent graph reproof and exact implementation).}
Let $p$ be prime, $q\ge2$, and
$\Gamma=C_p\amalg C_q=\langle a,b\rangle$ be the profinite
free product.  Let $\pi:\Gamma\twoheadrightarrow Q$ be a finite
quotient injective on both factors, and set $U=\ker\pi$.
Then
\[
 \operatorname{pr}_Q C_{\Gamma/([U,U]U^p)}(a)
       =\langle\pi(a)\rangle.
\]
For the literal quotient $Q_{12}$ of $B=C_2\amalg C_3$, the
kernel $U$ is free profinite of rank $97$.  The modulus-two
relative homology extension separates the harmonic frame residues,
and the modulus-three extension separates the cubic ones.  Their
combined quotient
\[
 \Gamma_6=B/([U,U]U^6),\qquad |\Gamma_6|=576\,6^{97},
\]
therefore determines both pure $Q_{12}$ frame residues from the
actual multiplier-one action, and hence determines $a_2,a_3$.

\emph{Proof.}
Use the Bass--Serre tree of the free product, with edges indexed
by $\Gamma$, and vertices indexed by the cosets of its two
factors.  The torsion-free kernel $U$ acts freely.  Its quotient
is the finite connected graph with edges $Q$ and vertices
$Q/\langle\pi(a)\rangle$ and $Q/\langle\pi(b)\rangle$.
Its cycle space over $\mathbb F_p$ identifies with
$U/[U,U]U^p$.  All statements may first be made for the discrete
free product and its finite-index free kernel, then completed;
the finite quotient and its finite homology are unchanged.

Choose as base vertex $o$ the identity coset of the first factor.
An element with $Q$-image $g$ has a path-chain coordinate $v$
with boundary $g o-o$.  Changing the lift by an element of $U$
adds an arbitrary cycle.  The generator $a$ fixes the base vertex
and has zero path-chain coordinate.  Consequently a lift of
$g\in C_Q(\pi(a))$ centralizes $a$ in the relative extension
only if some path chain $v$ with that boundary is fixed by $a$.
This follows directly from the path cocycle rule
$v_{gh}=v_g+g v_h$; the commutator condition is
$(a-1)v=0$.

The action of $a$ on the edge set $Q$ is free, with all orbits
of length $p$.  Every invariant edge chain is a linear combination
of their orbit sums.  At a fixed vertex, the $p$ contributions
of such an orbit have the same incidence sign, and sum to zero
over $\mathbb F_p$.  Hence the boundary of an invariant chain
has coefficient zero at every $a$-fixed vertex.

Both $o$ and $g o$ are $a$-fixed when $g$ centralizes $a$.
If they are distinct, the prescribed boundary $g o-o$ has a
nonzero coefficient at both, a contradiction.  Thus $g o=o$,
or $g\in\langle\pi(a)\rangle$.  Powers of $a$ themselves
centralize $a$ and project onto that factor, proving equality.

For $Q_{12}$, the quotient graph has $576$ edges and
$288+192$ vertices, hence rank $576-288-192+1=97$.
Apply the first assertion with $a=s,p=2$ and, after interchanging
the factors, with $a=r,p=3$.  The finite images of the two
factors meet the pure subgroup $q(F)$ trivially.  Therefore two
pure conjugators of $s$ with the same finite conjugate have the
same $Q_{12}$ residue; the analogous statement holds for $r$.
Actual frames supply existence.  Their positive signed residues
are selected by their identity reduction modulo four.  R3 then
recovers the shear from the frame abelianization modulo three,
and the corrected frame parities follow.\hfill$\square$

The independent computation
\path{agents/characters/code/check_relative_frame_homology.py}
constructs the actual graph, a spanning tree, and all $97$ cycle
coordinates.  It solves every lifting equation
$(a-1)m=-(a-1)v$ over the relevant field.  At two the original
centralizer has order $16$, its liftable image has order $2$,
and $\operatorname{rank}(s-1\mid H_1)=45$.
At three the corresponding numbers are $18,3,63$.
It also checks the fixed-vertex boundary obstruction directly
on every edge-orbit sum.  Full data are in
\path{agents/characters/records/relative_frame_homology.json}.

\paragraph{Exact overlap and remaining scope.}
The source already proves the same frame-separation mechanism:
FP.23, in \path{sources/DELTA_ON_GT/agent_five/five_point_results.tex},
uses $\mathbb F_2[Q/H_2]\oplus\mathbb F_3[Q/H_3]$ and fixed
cosets; RH.2, in \path{sources/DELTA_ON_GT/sections/A_homology.tex},
explicitly says that detector factors through the same $\Gamma_6$.
D23.C12 is a direct graph proof, its general $C_p*C_q$ formulation,
and an exact rank-$97$ implementation of the lifting obstruction.
It is not claimed as a new detector or a new pentagon constraint.
The rank-$129$ and rank-$4097$ Ferrari-kernel tests concern different
cover homologies and retain five-point information.  As FP.30 proves,
simply applying Ferrari and inducing any of the present four-point
detectors leaves the pentagon word identically one for every
two-cycle word.  A comparison with a representation that retains
the relevant five-point kernel remains necessary.

\subsection{A smaller quotient when only the two parities are needed}

\paragraph{D23.C13 (proved; strict compression of the supplied detector).}
There is an explicit quotient of the FP.23 source image with order
\[
 576\cdot2^{34}\cdot3^{52}
\]
such that the value of an actual multiplier-one gentle word $f$
in this quotient determines both corrected frame parities
$a_2,a_3$.  Its ambient permutation modules have dimensions
$72$ over $\mathbb F_2$ and $64$ over $\mathbb F_3$.
This is a quotient for reading actual normalized action values,
or equivalently the word value $f$; invariance of its kernel
under all gentle automorphisms is not asserted.  The stable
quotient $\Gamma_6$ of D23.C12 maps onto it, so an actual action
on $\Gamma_6$ supplies the required readout.

The original FP.23 source image has exact order
$576\cdot2^{97}\cdot3^{96}$.  Thus the present quotient is
strictly smaller; the proved readout here concerns the two
parities, with no full frame-residue claim.  It still gives no additional
pentagon equation after Ferrari extension alone.

\emph{Proof.}
Write $C_s=C_Q(q(s))$ and $C_r=C_Q(q(r))$, of orders $16$ and
$18$.  A signed lift of an element of $C_s$ commutes with $S$
at both primes: a possible simultaneous anticommutation is
excluded modulo four, since it would require a sum of two
squares to be $-1$ modulo four.  At three, such a lift is
$S^j$; at four it is $aI+bS$ with $a^2+b^2=1$.
Define
\[
 \eta_2(g)=j+b\pmod2.
\]
This is independent of the global sign of the lift: changing that
sign adds two to $j$ and leaves $b$ unchanged modulo two.
Both summands are homomorphisms, the second being reduction to
the two-element centralizer of $S$ over $\mathbb F_2$.
It is a nonzero character and $\eta_2(q(s))=1+1=0$.
On the positive pure centralizers with residue $I$ modulo four,
it is exactly the harmonic exponent parity.

For an element of $C_r$, choose a signed lift.  At both moduli
four and three its centralizer has the form
$\{\pm R_0^j: j\in\mathbb Z/3\}$.
The lift commutes with $R_0$ genuinely: anticommutation would
negate its trace $-1$, already impossible modulo three.
Let $e_4,e_3\in\mathbb F_2$ be the respective signs and put
\[
 \eta_3(g)=e_4+e_3.
\]
Changing the global lift changes both signs and preserves their
sum.  This is again a nonzero character, trivial on $q(r)$;
on a positive pure centralizer it is exactly the cubic sign.

Put $H_2^+=\ker\eta_2\subset C_s$ and
$H_3^+=\ker\eta_3\subset C_r$.  Their orders are $8,9$ and
they contain $q(s),q(r)$, respectively.  Form
\[
 V_2=\mathbb F_2[Q/H_2^+],\qquad
 V_3=\mathbb F_3[Q/H_3^+],
\]
of dimensions $72$ and $64$, and let $e_2,e_3$ denote the
identity-coset basis vectors.  As in FP.23 the elements
\[
 S^+=(e_2,0,q(s)),\qquad R^+=(0,e_3,q(r))
\]
in $(V_2\oplus V_3)\rtimes Q$ have orders two and three.
They define a finite homomorphism $\theta^+:B\to E^+$.
Since $\langle q(s)\rangle\subset H_2^+$ and
$\langle q(r)\rangle\subset H_3^+$, the two coset projections
give a surjective module map from FP.23's ambient modules to
these modules, compatible with both generators.  In particular
$E^+$ is a quotient of the actual FP.23 image.

If $(m,g)$ centralizes $S^+$, its projection centralizes $q(s)$
and satisfies $(1-q(s))m_2=(1-g)e_2$.  Passing to coinvariants
shows that the two $q(s)$-fixed cosets $H_2^+$ and $gH_2^+$
must coincide; hence $g\in H_2^+$.  Thus every centralizer
ambiguity has $\eta_2(g)=0$.  The same argument with $R^+$
gives $g\in H_3^+$ and $\eta_3(g)=0$.

Given the value $\theta^+(f)$ for an actual gentle element,
the normalized formulas give
\[
 \theta^+\Phi(s)=S^+\theta^+(f),\qquad
 \theta^+\Phi(r)=\theta^+(f)^{-1}R^+.
\]
Solve the finite conjugacy equations for frames in $\theta^+(F)$
whose positive signed $Q$-residues are $I$ modulo four.
Actual frames guarantee solutions.  The ratio of two harmonic
solutions centralizes $S^+$ and has positive residue $I$
modulo four, so its $\eta_2$ is precisely its possible change
of harmonic parity; that change is zero.  The corresponding
argument with $\eta_3$ proves uniqueness of the cubic parity.
The shear is also known: the positive signed value of $f$
has $I$ residue modulo four, and its upper-right entry modulo
three is $\rho$ by R3.  Hence $V=U^{8\rho}$ modulo twelve
is determined before the corrected parities are read.
This proves the asserted readout without assuming every finite
admissible conjugacy solution lifts to an actual gentle element.

For the exact size, the implementation uses $576$ Schreier
representatives for $U=\ker q$.  If $v_g$ is the affine
translation of a chosen representative and $t_a$ that of a
generator $a\in\{s,r\}$, the Schreier kernel generator has
translation
\[
 v_g+g t_a-v_{ga}.
\]
These $1152$ vectors span exactly the image of $U$ in each
module.  Exact row reduction gives ranks $34$ and $52$.
The joint image of $U$ is the full product of these elementary
abelian groups, because their orders are coprime and the joint
image surjects onto both.  Surjectivity onto $Q$ gives the
stated total order.  The same calculation in FP.23's original
modules gives ranks $97$ and $96$, proving strictness.
Every character equation, every coset obstruction, the Schreier
spans, and the image ranks are checked by
\path{agents/characters/code/check_compressed_parity_detector.py},
with exact records in
\path{agents/characters/records/compressed_parity_detector.json}.
Finally FP.30 applies to every quotient factoring through the
forgetful map, so compression does not change its automatic
pentagon conclusion.\hfill$\square$

\subsection{Inducing the marked $S_4$ quotient does not remove the harmonic frame ambiguity}

\paragraph{D23.F16.}
Combine the first-pair Krammer representation of D23.F13
with the induced $24$-point representation of
$\widehat B=C_2\amalg C_3$ from D23.C11, and retain
the signed projective level-twelve quotient as well.
This combination is formed as a representation of
$\widehat B_3$, pulling the latter two factors back
along $\widehat B_3\to\widehat B$. In its Krammer
coordinate, ``centralizing $s$'' means centralizing
the half-twist lift $s_P=K_1K_2K_1$, whose square
is central; the six-dimensional representation is
not asserted to factor through $\widehat B$.
The finite centralizer of $s$ in the image of $F'$
still contains an element $\beta$ with
\[
 q_4(\beta)=I,\qquad q_3(\beta)=S
 =\begin{pmatrix}0&1\\-1&0\end{pmatrix}.
\tag{D23.F16.1}
\]
Its induced $24$-point coordinate can be the identity.
Thus this combined quotient does not determine the
harmonic parity from the finite harmonic frame equation,
even after imposing $\beta\in\operatorname{im}(F')$.
This is a finite ambiguity, not an assertion that the
nonidentity finite conjugator is an actual pure frame.

\begin{proof}
Let $P$ be the Krammer first-pair group.  The exact
derived orders of D23.P6 are
$46656,3888,486,243,3,1$. Thus $P'$ is solvable
and its abelianization has order $8$.
Let $E$ be the restriction to $F$ of the induced
representation in D23.C11. Its exact derived orders
are $6912,1728,64,1$, so $E'$ is solvable and its
abelianization has order $27$.
The simultaneous image of $F'$ in $P'\times E'$
is subdirect. A nontrivial common quotient of these
solvable groups would have a nontrivial abelianization,
which would be a quotient of both a two-group and
a three-group. This is impossible. Goursat's lemma
therefore gives the full product $P'\times E'$.

Under the diagonal-block projection of D23.F13,
\[
 P'\longrightarrow E_F'\times D_8'
\]
is surjective, with central kernel $\langle n\rangle$
of order three.  In the affine coordinates of D23.P3
choose the first coordinate
$\gamma=(0,S)\in k^2\rtimes Q_8=E_F'$,
and choose the second coordinate to be identity.
It follows that there is $\beta_P\in P'$ with
these diagonal values.

Put $s_P=K_1K_2K_1$. Its affine $W$ coordinate is
$(0,iS)$, which commutes with $\gamma$; its quotient
coordinate commutes with identity. Thus
$[s_P,\beta_P]\in\langle n\rangle$.
This kernel is central in $\langle K_1,K_2\rangle$,
and $s_P^2=(K_1K_2)^3$ is also central there.
Writing $s_P\beta_Ps_P^{-1}=\beta_Pn^j$ and
conjugating once more gives $n^{2j}=1$.
Since $n$ has order three, $j=0$. Every such lift
already centralizes $s_P$.

By the full-product statement, the pair
$(\beta_P,1)$ is the image of an element of $F'$.
It centralizes $s$ in both representations.  The
signed level-twelve coordinate of an element of
$F'$ has identity component modulo four: the two
signed generators $U^2,L^2$ commute modulo four,
so their commutator subgroup is killed there.
On $F'$ the linear part of the affine detector is
exactly $q_3$, since the modulus-four character
occurring in D23.P3 is trivial. Our choice of
$\gamma$ therefore gives $q_3(\beta_P)=S$.
This also shows that adjoining the level-twelve
coordinate imposes no further restriction on the
chosen pair: on $F'$ it is already determined by
the Krammer linear coordinate.  The resulting
matrix is
\[
 q_{12}(\beta)=\begin{pmatrix}9&4\\8&9\end{pmatrix},
\]
whose two reductions are those of (D23.F16.1).

For a literal finite certificate, the program
\texttt{krammer\_induced\_frame\_ambiguity.py}
finds the three lifts of $(\gamma,1)$ inside the
enumerated $P'$, checks that all commute with
$s_P$, and chooses the unique lift of order four.
Its complete six-dimensional matrix is in the
companion JSON. The existence proof above does
not depend on a guessed lift or on an assumed
congruence-kernel invariance.

At the identity action the harmonic frame equation
is simply that the frame conjugator centralizes $s$.
Both the identity and the constructed finite
conjugator solve this equation and satisfy the
positive modulus-four normalization. The latter
has odd harmonic residue modulo three, whereas
the former has even residue. Thus the finite
equation does not recover the actual parity.
The actual pure centralizer in the profinite
source is trivial; our assertion concerns only
the specified finite quotient.
\end{proof}

\subsection{A two-primary harmonic cover still separates from the cubic Krammer state}

\paragraph{D23.F19.}
Let $\Gamma_K$ be the common-word Krammer image of
D23.F14 and D23.P9. Its maximal two-group quotient
on the derived subgroup is
\[
 \Gamma_K'\twoheadrightarrow Q_8\times C_2.
\tag{D23.F19.1}
\]
On the six actual-gentle necessary first-pair values
of D23.F13, this quotient takes the value
\[
 \left(M_\rho,1\right),\qquad
 M_\rho=\begin{pmatrix}1&\rho\\\rho&1+\rho^2\end{pmatrix},
\tag{D23.F19.2}
\]
independently of the cubic bit $a_3$.

Let $E_2$ be any finite quotient of $\widehat B$
mapping onto $Q_{12}$ with two-group kernel; this
includes the relative mod-two harmonic graph
cover of D23.C12. Put $T_2=\operatorname{im}(F'\to E_2)$.
In the simultaneous image of $F'$ in
$\Gamma_K'\times T_2$, every Krammer value in the
kernel of (D23.F19.1) can be paired with the
identity of $T_2$.
Consequently all the reduced Krammer pentagon
solutions with $\rho=0,a_3=1$ can be combined
with the identity harmonic-cover word value.
By D23.P9 these combinations have a common
commutator-word realization. This refutes a
finite certificate obtained merely by adjoining
this harmonic cover by product and the Ferrari
forgetful map. It does not produce an actual
mixed-sign gentle solution of the full pentagon.

\begin{proof}
For the first-pair group $P$, D23.F13 gives
a central order-three kernel over
$E_F\times D_8$. On derived groups the quotient
is $E_F'\times D_8'$, where
$E_F'=k^2\rtimes Q_8$ and $D_8'=C_2$.
The inverse image of $k^2\times1$ is a normal
three-subgroup $N$ of order $3^5=243$ in $P'$,
and
\[
 P'/N\simeq Q_8\times C_2.
\]
Every homomorphism from $P'$ to a two-group
kills $N$, since the order of an image of a
three-group must be both a power of three
and a power of two. Conversely the displayed
quotient is itself a two-group. It is therefore
the maximal two-group quotient of $P'$.

D23.P9 identifies the remaining factors of
$\Gamma_K'$ as three groups $H_j'$ of order
$27$ and three perfect groups $SU_3(3)$.
Every homomorphism to a two-group kills the
former by order and the latter because a
finite two-group is solvable and cannot
contain a nontrivial perfect image. Thus
the maximal two-group quotient of the
entire common-word image is precisely
(D23.F19.1).

In the affine $W$ coordinate, the six values
have linear part $M_\rho$, by D23.P3; their
translation part is where $a_3$ occurs.
Their $D_8$ quotient coordinate is identity,
by D23.F13. Killing $N$ removes exactly
these translations and the central
order-three lift, which proves (D23.F19.2).
In particular both values with $\rho=0$
are in the kernel of the maximal two-group
quotient.

The signed level-twelve image of $F'$
is $Q_8$: modulo four it is identity,
and modulo three it is
$\operatorname{SL}_2(\mathbf F_3)'=Q_8$.
The first assertion follows directly
because $U^2,L^2$ commute modulo four;
the exact image modulo three is the
standard pure-pair image. Thus $T_2$ is
an extension of a subgroup of a two-group
by a subgroup of the assumed two-group
kernel of $E_2\to Q_{12}$, and is itself
a two-group.

Let $J\subset\Gamma_K'\times T_2$ be
the joint image. It is subdirect. The
common quotient in Goursat's lemma is
a two-group because it is a quotient
of $T_2$. The corresponding map from
$\Gamma_K'$ factors through its maximal
two-group quotient. Hence if $g$ belongs
to the kernel of (D23.F19.1), its value
in that common quotient is identity;
the fibre-product description gives
$(g,1)\in J$.

D23.P7 supplies $36$ reduced solutions
for each pair $(\rho,a_3)$, and D23.P9
supplies common commutator words for
their full five-factor assignments.
For $\rho=0,a_3=1$ their first value
has trivial maximal two-group image,
and none of the mixed factors contributes
to that quotient. The just-proved
Goursat assertion therefore supplies
common words realizing those same
Krammer assignments and the identity
in $T_2$.

Under the Ferrari forgetful map the
five cyclic inputs are exactly
$(x,y,x,1,1)$, as in FP.30. Consequently
a commutator word that is identity in
the harmonic cover has identity in
each of its five evaluation factors:
the reversed pair is conjugate by
the half twist, and factors with an
identity argument are identity.
The combined finite pentagon is
therefore satisfied. If the harmonic
frame is separated by the cover of
D23.C12, the identity normalized
action has the even harmonic residue.
Thus the added cover cannot discard
the unequal-bit finite assignments.
The common words need not satisfy
the full profinite gentle relations;
no claim of such a lift is used.
\end{proof}

\paragraph{D23.F20 (the whole relative pro-two Ferrari family).}
Let $L=\widehat B_4/\langle z_4\rangle$, let
$q_L:L\to\widehat B$ be the Ferrari forgetful map,
and put $N=\ker(q_{12}q_L)$.
The actual mixed-sign gentle witness $k$ of
DELTA\_ON\_GT FP.36, supplied by the strengthened
relative lifting input LI.8, satisfies every
five-cycle factor separately in
\[
 L/O^2(N).
\]
This holds for the entire relative pro-two
completion, not just its first homology layer.

\begin{proof}
The kernel of the map from this quotient
to $Q_{12}$ is pro-two. For each of
the five adjacent evaluation pairs,
its image in $Q_{12}$ is the ordinary
pure pair or a cyclic group with
one generator trivial, by FP.30.1--2
and the reversing half twist.
The pure pair has derived series
of orders $96,8,2,1$; in particular
its second derived group is $C_2$.
The other base pairs have trivial
second derived group. Hence the
second derived subgroup of every
pair in $L/O^2(N)$ is pro-two:
it is an extension of a pro-two
kernel by a finite two-group.

FP.36 gives one actual gentle
$k=(1,f_k)$ with $\Delta(k)=1$
and $f_k\in R_{\rm nil}(F'')$.
Every continuous homomorphism
from $F''$ into a pro-two group
kills this nilpotent residual,
because all its finite quotients
are nilpotent. Apply this to
each evaluation $F''\to P_j''$.
Each of the five values of $f_k$
is identity, proving the claim.
This uses the stated actual
lifting input; it does not
independently reconstruct LI.8.
\end{proof}

\subsection{The joint cubic--quartic cyclic cover}

\begin{statement}{D23.F21}{Proved}{The degree-twelve cyclic cover is blind at every two-adic and three-adic homology level}
\end{statement}
Let $C$ be the connected cyclic cover of the five-punctured sphere with
deck group $C_{12}$ and local monodromies
\[
 (9,9,10,4,4)\in(\mathbb Z/12)^5.
\]
Fill all punctures upstairs, and let $T$ denote the positive deck
transformation on the resulting integral homology. The compact curve has
genus $11$. Its degree-four and degree-three quotient covers have branch
data $(1,1,2,0,0)$ and $(0,0,1,1,1)$, respectively, and each has genus one.
Thus this single cover is a natural geometric attempt to retain the
quartic and cubic elliptic structures together.

Use the five pure braid elements
\[
 X_1=\sigma_1^2,\quad X_2=\sigma_2^2,\quad
 X_3=\sigma_3^2,\quad X_4=(\sigma_1\sigma_2\sigma_1)^2,
 \quad X_5=(\sigma_3\sigma_2\sigma_3)^2,
\]
and descend their homology action from $PB_4$ to the pure five-point
mapping class group by the linking-number normalization specified below.
For $p=2$ and for $p=3$, the closed subgroup generated by any adjacent
pair $(X_j,X_{j+1})$, with indices modulo five, has pro-$p$ second
derived subgroup in its action on $H_1(C,\mathbb Z_p)$.

Consequently the supplied actual gentle element $k$ of
\textup{FP.36}, with
\[
 f_k\in R_{\rm nil}(\widehat F(x,y)'')\qquad\hbox{and}\qquad
 \Delta(k)=1,
\]
has all five evaluations $f_k(X_j,X_{j+1})$ equal to the identity on
both $H_1(C,\mathbb Z_2)$ and $H_1(C,\mathbb Z_3)$, hence on every
finite quotient of these lattices. This assertion uses the supplied
existence statement for $k$; it does not reconstruct that lifting input.
It proves no claim about other primes or about nonabelian covering
invariants.

\begin{proof}
The sum of the branch weights is zero modulo $12$, and the weights
generate $\mathbb Z/12$, proving that the cover is well-defined and
connected. A puncture of weight $w$ has $\gcd(12,w)$ preimages, each
with ramification index $12/\gcd(12,w)$. The five numbers of preimages
are $3,3,2,4,4$. Removing the five branch values gives Euler
characteristic $12(2-5)=-36$; adding the $16$ points upstairs gives
$\chi(C)=-20$, so its genus is $11$. The same count at degrees four
and three gives Euler characteristic zero for both stated quotient
covers.

Here is a complete finite cellular model, including the marking used
by the certificate. Remove the fifth point and represent the other four
inertias by the free generators $x_1,x_2,x_3,x_4$. On the degree-twelve
cover of the resulting four-petal rose put
\[
 C_0=\mathbb Z^{12},\qquad C_1=\mathbb Z^{48},\qquad
 \partial e_{j,g}=v_{g+w_j}-v_g,
 \quad(w_1,w_2,w_3,w_4)=(9,9,10,4).
\]
Here $g\in\mathbb Z/12$, and $T(e_{j,g})=e_{j,g+1}$.
The graph is connected, so $\ker\partial$ has rank $37$.
The compact surface is obtained by filling all lifted peripheral
circles: for the first four punctures these are the translates of
$x_j^{12/\gcd(12,w_j)}$, and for the fifth they are the translates of
$(x_1x_2x_3x_4)^3$. Choose one translate for every coset of
$\langle w_j\rangle$ or $\langle8\rangle$, respectively.
There are sixteen such circles, with one relation, so their image has
rank fifteen and the quotient homology has rank $22$.
Equivalently, filling the boundary circles of a connected orientable
surface gives precisely their sole total-boundary relation. This also
shows that the quotient lattice is torsion-free.

The positive Artin automorphism is
\[
 \sigma_i(x_i)=x_ix_{i+1}x_i^{-1},\qquad
 \sigma_i(x_{i+1})=x_i,
\]
and fixes the other generators. The usual braid relations follow by
substitution in this free group; the certificate also checks all three
relations relevant to $B_4$ as literal reduced words. Each pure braid
preserves the weight map, acts on the displayed graph chains, preserves
the peripheral submodule, and commutes with $T$.

Write $P=x_1x_2x_3x_4$. The full twist
$z_4=(\sigma_1\sigma_2\sigma_3)^4$ acts by conjugation by $P$.
On the homology of the kernel of the weight map, this is the deck
action $T^8$. This identity follows because $P$ has weight $8$ and
changing a lift by an element of the kernel changes the action only
by an inner automorphism, which is trivial on its abelianization.
Let $\ell_{12}:PB_4\to\mathbb Z$ be the first-two-strands linking
number. Since $\ell_{12}(z_4)=1$, the corrected representation
\[
 \gamma\longmapsto T^{-8\ell_{12}(\gamma)}q(\gamma)
\]
is trivial on $z_4$ and descends to $PB_4/\langle z_4\rangle$.
For the five displayed elements the linking numbers are
$(1,0,0,1,0)$, so the correction multiplies precisely $X_1$ and
$X_4$ by $T^4$. Because the deck action is central, this correction
does not change the evaluation of any word $f\in\widehat F'$ at a
pair of elements.

Reduce the homology modulo $p$. The semisimple submodule, or socle,
for the deck action is
\[
 S_p=\bigoplus_h\ker h(T),
\]
where $h$ runs through the distinct monic irreducible factors of
$T^{12}-1$ in $\mathbb F_p[T]$. For $p=2$ these are
$T-1,T^2+T+1$; for $p=3$ they are $T-1,T+1,T^2+1$.
The chain calculation just described gives the following table. Each
entry is the sequence of orders of the successive derived groups of
the indicated adjacent-pair image on the indicated socle summand;
the last group is trivial.
\[
\begin{array}{c|c|c|ccccc}
 p&h&\dim_{\mathbb F_p}\ker h(T)&12&23&34&45&51\\\hline
2&T-1&1&1&1&1&1&1\\
2&T^2+T+1&6&(9,1)&(36,4,1)&(36,4,1)&(36,4,1)&(9,1)\\
3&T-1&1&1&1&1&1&1\\
3&T+1&3&(2,1)&(2,1)&1&(2,1)&(2,1)\\
3&T^2+1&6&(36,9,1)&(72,9,1)&(2,1)&(4,1)&(72,9,1)
\end{array}
\]
In particular every adjacent-pair image on $S_p$ is metabelian.

For reproducibility, \texttt{cyclic12\_cover.py} builds the boundary
matrix and all sixteen peripheral columns over the integers, computes
the compact quotient separately over $\mathbb F_2$ and $\mathbb F_3$,
and derives each pure braid matrix by tracing the literal Artin words
through the covering graph. It checks that every matrix preserves the
relation subspace, that the compact dimension is $22$, that $T^{12}=1$,
that the pure matrices commute with $T$, and that the full twist acts
as $T^8$. It then restricts to every irreducible deck socle summand.
For each pair it closes the orbit of a basis under the pair's matrices;
the resulting permutation action is faithful, since a linear map
fixing that orbit fixes the basis. Exact permutation derived-group
calculations produce the table. The outputs are
\texttt{cyclic12\_cover.json} and the two \texttt{.npz} matrix records.
Thus the computation starts with the specified covering and free-word
action, rather than assuming guessed braid matrices or a quotient
identification.

We next justify why this finite table controls all $p$-adic depths.
For a finite-dimensional $\mathbb F_p[T]$-module $V$, the kernel of
\[
 \operatorname{Aut}_{\mathbb F_p[T]}(V)
 \longrightarrow\operatorname{Aut}_{\mathbb F_p[T]}(\operatorname{soc}V)
\]
is a $p$-group. Indeed, if the kernel had an element of prime order
$q\ne p$, extend scalars to an algebraic closure. That element is
semisimple and, if nonidentity, has a nonzero eigenspace with eigenvalue
different from one. The eigenspace is $T$-stable because the two actions
commute. A nonzero $T$-stable space contains a $T$-eigenvector, which
lies in the scalar extension of the socle. This contradicts trivial
action on the socle. Cauchy's theorem now proves the claim.

Every adjacent-pair second derived group acts trivially on $S_p$ by
the table, so its image on $H_1(C,\mathbb F_p)$ is a $p$-group.
The kernel of reduction from $\operatorname{GL}_{22}(\mathbb Z_p)$
to $\operatorname{GL}_{22}(\mathbb F_p)$ is pro-$p$, as is seen from
the successive congruence kernels. Therefore the closed second
derived group of each $p$-adic pair image is pro-$p$.

Finally the restriction of each evaluation homomorphism
$\widehat F\to\operatorname{GL}_{22}(\mathbb Z_p)$ to
$\widehat F''$ has pro-$p$ image. Every finite continuous quotient
of that image is a finite $p$-group and hence nilpotent. It follows
directly from the definition of $R_{\rm nil}(\widehat F'')$ as the
intersection of kernels of finite nilpotent quotients that the
restriction kills $R_{\rm nil}(\widehat F'')$. Apply this to the
supplied $f_k$ for each of the five pairs and for $p=2,3$.
\end{proof}

\paragraph{Scope and source locators.}
The five pure braid elements and their pentagon order are those in
\texttt{DELTA\_ON\_GT/agent\_five/five\_point\_results.tex},
\textup{FP.30}. The sole actual-element existence input is
\textup{FP.36} in the same file, equivalently the strengthened
relative lifting witness used in \textup{A.94} and \textup{A.109}.
The cyclic covering, its chain action, the socle calculation, and the
all-depth implication above are new calculations in this run.

\subsection{Two noncentral Hurwitz tests and an explicit finite escape}

\paragraph{D23.P11 (genuine five-point Hurwitz quotients).}
For a finite group $G$ and $p\in G$, the Hurwitz moves
\[
 (g_i,g_{i+1})\longmapsto(g_ig_{i+1}g_i^{-1},g_i)
\tag{D23.P11.1}
\]
act on four-tuples of product $p$, modulo simultaneous conjugation
by $C_G(p)$. Every orbit is a permutation representation of
$B_4/\langle z_4\rangle$, where $z_4=(\sigma_1\sigma_2\sigma_3)^4$.
In the following two explicit cases the product is nonidentity:

\begin{enumerate}
\item Let $G=\mathbb F_3^2\rtimes\operatorname{GL}_2(\mathbb F_3)$,
of order $432$, and let all four entries lie in the conjugacy
class of $(0,\operatorname{diag}(1,-1))$, of size $36$.
Take $p=((0,1)^{\mathsf T},U)$. Its centralizer has order $9$.
The quotient set of such tuples has $891$ elements and four
$B_4$-orbits of sizes $9,72,162,648$. Use the orbit of size
$648$, containing
\[
 \left(
  (0,\begin{psmallmatrix}0&1\\1&0\end{psmallmatrix}),
  (0,\begin{psmallmatrix}0&1\\1&0\end{psmallmatrix}),
  (0,\begin{psmallmatrix}1&1\\0&2\end{psmallmatrix}),
  ((1,2)^{\mathsf T},\begin{psmallmatrix}1&2\\0&2\end{psmallmatrix})
 \right).
\]
Call this representation $q_{648}$.
\item Let $G=E=\mathbb F_9^2\rtimes H_{48}$ be D23.P3,
with $s=(0,iS)$ and $r=(e_1,R_0)$. The tuple $(s,s,r,r)$
has product $p=r^2\ne1$, whose centralizer has order $27$.
Its Hurwitz orbit modulo this centralizer has size $20736$.
Call this representation $q_{20736}$.
\end{enumerate}

Both tests contain mixed evaluation pairs with nonnilpotent
second derived subgroups; thus neither of these selected tests
is dismissed by the supplied FP.36 nilpotence criterion.
For example, a degree-$144$ component of
$\langle X_3,X_4\rangle$ in $q_{648}$ has derived orders
\[
408146688,\quad102036672,\quad34012224,\quad531441,\quad1.
\tag{D23.P11.2}
\]
It admits the literal marked quotient
$X_3\mapsto(123),X_4\mapsto(1234)$ onto $S_4$.
A degree-$96$ mixed component of $q_{20736}$ has derived orders
\[
3265173504,\quad272097792,\quad4251528,\quad531441,\quad1.
\tag{D23.P11.3}
\]
Here, as throughout this branch,
\[
 X_1=A^2,\quad X_2=B^2,\quad X_3=C^2,\quad
 X_4=(ABA)^2,\quad X_5=(CBC)^2.
\]

\begin{proof}
The elementary Hurwitz moves satisfy the braid and distant
relations, preserve the product, and commute with simultaneous
conjugation. The full twist acts by simultaneous conjugation
by the product (or its inverse if the opposite action convention
is chosen); this is trivial after quotienting by $C_G(p)$.
Thus these are quotients of the label-five stabilizer used by
the actual five-cycle, not the product-one forgetful family.

All displayed orbit sizes and component orders are finite exact
computations. The program
\path{agents/pentagon/hurwitz_affine_reflections.py}
enumerates the first set using the first three entries of a tuple,
the fourth being determined by its product. It quotients by the
explicit centralizer and checks all braid and full-twist relations.
The program \path{agents/pentagon/hurwitz_cubic_affine.py}
builds $E$ from its literal affine generators, enumerates the
single orbit of $(s,s,r,r)$, and checks the same relations.
There is no random orbit sampling in these constructions.

For (D23.P11.2)--(D23.P11.3), the exact permutation groups are
computed on a component of the specified evaluation pair, and
their derived subgroups are computed successively. The order-$531441$
third derived subgroup is elementary abelian of exponent three;
the preceding derived subgroup is checked to be nonnilpotent.
Adjoining the marked four-point action to the degree-$144$
component leaves its group order unchanged, proving the stated
marked $S_4$ quotient. See
\path{hurwitz_affine_affine_unipotent_3_components.json},
\path{hurwitz_cubic_ssrr_components.json}, and
\path{hurwitz_affine_mixed_marked_quotients.json}.
The first of these analyzes all components of the selected
$648$-point orbit; the second leaves larger components unanalyzed,
and no assertion about those components is made here.
\end{proof}

\paragraph{D23.P12 (one ordinary word defeats both finite tests).}
Use $[a,b]=aba^{-1}b^{-1}$ and let $\iota$ interchange $x,y$.
Define ordinary words
\[
\begin{aligned}
 c&=[x,y],&d&=[c,xcx^{-1}],&u&=[d,xdx^{-1}],\\
 v&=u(c^4uc^{-4}),&w&=[c^6,v]^{-1},\\
 b&=(xy)w^{-1}(xy)^{-1},& f_*&=\iota(b)b^{-1}.
\end{aligned}
\tag{D23.P12.1}
\]
Then $f_*\in F'''$ and $\iota(f_*)=f_*^{-1}$ as an exact
identity of ordinary free words. Its expanded length is $904$
(freely reduced length $588$). In D23.P3 and D23.P10 it has
the exact values
\[
 \eta(f_*)=((1-i,1+i)^{\mathsf T},I),\qquad
 \xi(f_*)=((1,1)^{\mathsf T},I),
\tag{D23.P12.2}
\]
while its marked $S_4$ value is $1$. In both $q_{648}$ and
$q_{20736}$, all ten orientations of the five pentagon pairs
have identity $f_*$-value. In particular the entire literal
five-cycle equals $1$ in both representations.
The same identity values hold on the entire Nielsen $S_3$
closure of these pairs, modulo inner conjugacy. Thus all gentle
two- and three-cycle identities in that finite closure hold.

Moreover, its complete first-pair Krammer value is the matrix
$F_{0,1}$ of D23.F13, not merely the projection of that matrix
to the cubic detector. Thus these finite tests, even taken
together with the exact first-pair Krammer state and the marked
$S_4$ state, do not discard the mixed state
$(\rho,a_3,\nu_4)=(0,1,0)$.

This statement concerns an ordinary word and specified finite
images. It does not claim the free three-cycle for $f_*$,
gentle admissibility, the Krammer pentagon, or a full-GT lift.

\begin{proof}
The derived-depth assertion follows directly from the nested
commutators: $c\in F'$, $d\in F''$, $u\in F'''$, and
normality of $F'''$ puts $v,w,b,f_*$ in $F'''$. The identity
\[
 \iota(f_*)=b\iota(b)^{-1}=f_*^{-1}
\]
is formal. The $S_4$ value is therefore trivial already because
$S_4'''=1$; this part does not use an unproved gentle identity.

The remaining assertions are exact evaluations of the explicitly
given ordinary word. To make the computation reproducible,
\path{agents/pentagon/hurwitz_escape_word.py} constructs the
expanded word from (D23.P12.1), freely checks the two-cycle,
checks the saved Hurwitz generators' braid and full-twist
relations again, and evaluates the word on the five ordered pairs
\[
 (X_3,X_4),(X_5,X_1),(X_2,X_3),(X_4,X_5),(X_1,X_2)
\tag{D23.P12.3}
\]
and on their reversals. All twenty resulting permutations are
identity. In fact the intermediate $448$-letter word $w$ has
identity values on these same twenty pairs.
For each ordered pair $(a,b)$ and its reversal the certificate
also sets $z=(ab)^{-1}$ and evaluates $f_*$ on
$(a,b),(b,z),(z,a)$. All sixty values across the two families
are identity. These representatives suffice for the entire
Nielsen $S_3$ closure: the transformations
$\tau:(x,y)\mapsto(y,(xy)^{-1})$ and
$\iota:(x,y)\mapsto(y,x)$ satisfy
$\tau^3=\iota^2=1$ and $(\iota\tau)^2=\operatorname{Inn}(x^{-1})$.
Inner conjugation merely conjugates a word value. In particular
each product $f_*(z,a)f_*(b,z)f_*(a,b)$ is identity in both tests.

The original affine evaluation, with the encoding $a+3b$ for
$a+bi\in\mathbb F_9$, is
\[
 \eta(w)=((1,1-i)^{\mathsf T},I),\qquad
 \eta(b)=\eta(f_*)=((1-i,1+i)^{\mathsf T},I).
\]
The change from $w$ to $b$ is literal affine conjugation; the
change from $b$ to $f_*$ is also evaluated, rather than inferred
from an assumed three-cycle. The program independently checks
the compressed $3\times3$ homogeneous matrix
\[
 \xi(f_*)=\begin{pmatrix}1&0&1\\0&1&1\\0&0&1\end{pmatrix}.
\]
Finally it evaluates $f_*$ at the actual $12\times12$ first-pair
Krammer matrices and compares all entries with the saved
D23.F13 matrix labelled $(0,1)$. This equality also follows
from D23.F13: the quotient $D_8$ is killed by derived depth,
the affine block has the required value, and the exact two-cycle
fixes the remaining central phase uniquely.

The exact JSON certificate stores the entire word, all zero
support counts, both affine values, the Krammer label, and hashes
of the two finite Hurwitz actions. No existence of an unspecified
word or a profinite lift is used in these finite evaluations.
\end{proof}

\paragraph{Scope of the failure.}
The two tested Hurwitz representations are genuinely new mixed
tests: one sees the marked harmonic $S_4$ quotient, and both
have a nonnilpotent second derived mixed pair. Their failure is
therefore established by an explicit word, not by misapplying
the earlier nilpotence criterion. The full Nielsen closure was
checked as a further test and is also annihilated by this one
word. The outstanding finite issue on this branch is the full
Nielsen closure of the mixed Krammer pairs, jointly with their
pentagon correlations; D23.P12 does not decide it.


\section{The cubic dessin and distribution comparisons}
\subsection{An explicit rational quotient of the cubic detector dessin}

\paragraph{D23.C14 (the marked genus-zero quotient and its cyclic pullback).}
Let $\xi:F\twoheadrightarrow\mathbb F_3^2\rtimes\operatorname{SL}_2(\mathbb F_3)$
be the marked cubic detector of D23.P10. Its degree-nine action has a
canonical extension to
\[
 \Gamma_+=\ker(C_2*C_3\longrightarrow C_2)
       =\langle u=r,v=srs\mid u^3=v^3=1\rangle.
\]
In the same affine origin as D23.P10, the images are
\[
 u=\left((2,2)^{\mathsf T},
       \begin{pmatrix}0&2\\1&2\end{pmatrix}\right),\qquad
 v=\left((1,1)^{\mathsf T},
       \begin{pmatrix}2&2\\1&0\end{pmatrix}\right).
\tag{D23.C14.1}
\]
Here $x=(vu)^{-1}$, and the branch-cycle types of $(x,v,u)$ are
\[
                   (621),\qquad(33111),\qquad(333).
\tag{D23.C14.2}
\]
The corresponding degree-nine rational Belyi map is
\[
 \Psi(z)=\frac{432z^6(z-1)(4z-1)^2}
                   {(19z^3-9z^2-3z+1)^3}.
\tag{D23.C14.3}
\]
Put
\[
 \begin{split}
 C(z)&=53z^3+15z^2-3z-1,\\
 Q(z)&=z^2-4z+1,\qquad H(z)=19z^3-9z^2-3z+1.
 \end{split}
\]
Its cyclic cubic pullback has the rational model
\[
                  w^3=C(z),
\tag{D23.C14.4}
\]
which is explicitly birational over $\mathbb Q$ to
$E:Y^2=X^3-4$. The formulas in both directions are
\[
 \begin{split}
 G&=-X^2-6X-3Y+10,\\
 z&=(2X+Y-2)/G,\qquad w=(X^2-2X+2Y+4)/G,\\
 X&=(4+16z-2w)/(1+z+w),\\
 Y&=6(w^2-19z^2-8z-1)/(1+z+w)^2.
 \end{split}
\tag{D23.C14.5}
\]
Let $n=(1+\sqrt{-3})/2$, $\bar n=1-n$, and let $\phi$ be the
normalized first-orbit degree-nine map in the retrieved primary BelyiDB data,
with $x_{\rm old}=9X/1600$ and $y_{\rm old}=27Y/64000$.
Then its compact expression and the precise composition identity are
\[
 \phi=\bar n\,\frac{H+wQ}{H-nwQ},\qquad
 \boxed{\ \Psi(z)=1-\left(\frac{\phi-\bar n}{\phi-n}\right)^3.\ }
\tag{D23.C14.6}
\]
Thus the square formed by $E\xrightarrow{\phi}\mathbb P^1_\lambda$,
$E\xrightarrow{z}\mathbb P^1_z$, $\mathbb P^1_z\xrightarrow{\Psi}\mathbb P^1_b$,
and $b=1-((\lambda-\bar n)/(\lambda-n))^3$ commutes with the stated
marking, not just with the stated passports.

The map $E\to\mathbb P^1_z$ is cyclic of degree three over
$\mathbb Q(\sqrt{-3})$ and has three ramification points, the points
$w=0$. It is therefore the quotient by a geometric order-three
symmetry with fixed points. In particular, the fact that $u$ has no
fixed sheet in (D23.C14.2) does \emph{not} make this symmetry a
fixed-point-free translation: the fixed sheets of $v$ account for
its three fixed points.
More explicitly, $w\mapsto\zeta^2w$, with $\zeta=n^2$, is
\[
       \alpha(Q)=\zeta^2Q+(1-\sqrt{-3},2\sqrt{-3})
\]
in the group law on $E$. It cycles the sixfold ramification points
$O$, $(1-\sqrt{-3},2\sqrt{-3})$, and
$(1+\sqrt{-3},-2\sqrt{-3})$.

\begin{proof}
The subgroup presentation is the index-two Reidemeister--Schreier
presentation of $C_2*C_3$. Restrict D23.P3's affine representation to
$\Gamma_+$. Its linear matrices then lie in
$\operatorname{SL}_2(\mathbb F_3)$, so taking the real part of the
translation is a homomorphism. The origin shift $(0,2)^{\mathsf T}$
from D23.P10 gives (D23.C14.1), and direct multiplication gives its
marked $x,y$. Their action on the nine affine points has the cycles
(D23.C14.2). In the labeling $(a,b)\mapsto3a+b$, the simultaneous
conjugator to the BelyiDB rational branch pair $(\sigma_0,\sigma_1)$ is
\[
                     [4,0,5,2,1,6,3,8,7].
\]
The certificate checks the two literal conjugacy equations. The
rational database passport has size one, and its map, after the
source coordinate change $x_{\rm DB}=7z$, factors as (D23.C14.3).
Equivalently, the formula itself has the required branch data since
\[
                 \Psi(z)-1=\frac{Q(z)^3C(z)}{H(z)^3},
\tag{D23.C14.7}
\]
all indicated factors are separable and coprime, and the partitions
in (D23.C14.2) have total ramification $16=2\cdot9-2$.

For the orbifold topology, let $N$ be the inverse image of a point
stabilizer in $\Gamma_+$, and put $U=N\cap F$. There are three
cusps, from the cycles of $x$, no elliptic points above $u$, and
three order-three elliptic points above $v$. Since
$\chi(\Gamma_+)=-1/3$ and $[\Gamma_+:N]=9$, the signature equation is
\[
 -3=2-2g(N)-3-3(1-1/3),
\]
so the compactification of $N$ has genus zero. A conjugate of $v$
fixing a sheet maps nontrivially to $\Gamma_+/F\simeq C_3$.
Consequently $[N:U]=3$ and $U\lhd N$. The three elliptic points are
exactly the ramification points of the compact cyclic cover.
The cusps split into nine cusps because $x\in F$. Therefore
$\chi(U)=-9=2-2g(U)-9$, and its compactification has genus one.
This also proves the fixed-point assertion without choosing a
potentially non-normalizing lift of $r$.

For the exact algebraic map, write
$L=2X+Y-2$ and $N_0=X^2-2X+2Y+4$. Modulo $Y^2-X^3+4$, one has
\[
 N_0^3=53L^3+15L^2G-3LG^2-G^3.
\]
This proves (D23.C14.4) under (D23.C14.5). Substitution verifies both
inverse formulas and the equation of $E$; hence these are
birational maps of smooth projective curves. Since $C$ is squarefree,
with discriminant $-11664$, $w^3=C(z)$ gives precisely three branch
points and no ramification at infinity.
The automorphism $w\mapsto\zeta^2w$ acts by $\zeta^2$ on the regular
differential $dz/w^2$, and its image of $O$ is
$(1-\sqrt{-3},2\sqrt{-3})$ by the inverse formulas. These two facts
give the displayed group-law expression for $\alpha$ and its orbit.

The normalized database numerator has the form $P(X)Y+R(X)$ and
denominator $D(X)$, recorded in the root's exact normalization file.
With
\[
 H_h=19L^3-9L^2G-3LG^2+G^3,\qquad
 Q_h=L^2-4LG+G^2,
\]
the following polynomial identity modulo $Y^2-X^3+4$ fixes the sign
and all scalar factors:
\[
 (PY+R)(H_h-nN_0Q_h)
       =\bar n D(H_h+N_0Q_h).
\]
It proves the first expression in (D23.C14.6). Solving for the
cross-ratio gives
\[
              \frac{\phi-\bar n}{\phi-n}
                       =n\,\frac{wQ}{H}.
\]
Because $n^3=-1$ and $w^3=C$, its cube is $-CQ^3/H^3$.
Equation (D23.C14.7) proves the boxed composition identity, with its
orientation as written. An earlier direct cross-multiplication of
the degree-27 composition checked the same identity and rejected
its opposite orientation; the displayed lower-degree identities
provide a shorter exact replay certificate.

For a normalization check, let $P_0=(2,2)$ on $E$. Then
$2P_0=(5,-11)$ and
\[
 (\phi)_0=6(O)+2(2P_0)+(-P_0),\qquad
 z(O)=0,\quad z(2P_0)=1/4,\quad z(-P_0)=1.
\]
The point with coordinates $(5,11)$ is not a zero of $\phi$.
These signs follow directly from the normalized numerator, not
from an abstract identification of the elliptic curves.

The exact scripts are
\path{agents/characters/code/check_cubic_orbifold_topology.py} and
\path{agents/characters/code/check_belyi_rational_quotient.py};
the corresponding JSON certificates are in
\path{agents/characters/records/}.
The primary database records are Michael Musty's BelyiDB,
\path{belyi_db/9/9T23-[6,3,3]-621-33111-333-g0.m} and
\path{belyi_db/9/9T23-[6,6,6]-621-621-621-g1.m},
fields \texttt{BelyiDBPermutationTriple}, \texttt{BelyiDBCurve},
and \texttt{BelyiDBBelyiMap}, first Galois orbit for the genus-one
record. The repository is
\url{https://github.com/michaelmusty/BelyiDB}.
The retrieved genus-zero record and its MIT license are retained
in the certificate directory. The normalized genus-one coefficients
are in \path{records/belyi_9T23_normalized.txt}; the marked
permutation comparison is in \path{records/belyi_9T23_match.json}.
\end{proof}

\paragraph{D23.C14a (the precise five-point comparison left by this diagram).}
In the coordinate
\[
                 t=(\lambda-\bar n)/(\lambda-n),\qquad b=1-t^3,
\]
the source of the cyclic cover is
$\mathbb P^1\setminus(\{0,\infty\}\cup\mu_3)$.
Writing $\zeta=n^2$, the points
$t=\zeta^2,\zeta,1$ are respectively $\lambda=0,1,\infty$,
and $t=0,\infty$ are $\lambda=\bar n,n$.
Thus recovering the original marked $\lambda$ sphere fills
$t=0,\infty$. The standard Enriquez distribution map instead
retains $t=0,1,\infty$ and fills the other two roots of unity.
The required relation between these two forgetting maps is an
additional marked comparison to be established; it is not being
asserted as a consequence of the existence of the diagram.

Also, $\Psi$ is indecomposable as a rational cover. Indeed its
monodromy is the affine group $\mathbb F_3^2\rtimes
\operatorname{SL}_2(\mathbb F_3)$ in its primitive degree-nine
action: a block containing zero is a translation subgroup stable
under the irreducible linear stabilizer. There is no proper such
subgroup. Consequently $\Psi$ cannot itself be a nontrivial
composition of degree-three rational Kummer covers. This rules out
that direct shortcut to applying distribution, while leaving the
five-point comparison available.

\subsection{A general limitation of fixed-target distribution tests}

\paragraph{D23.D1 (simultaneous distribution tests in a normally generated target).}
Let $F=\langle x,y\rangle$ be the free group, let $G$ be a finite
group generated by $X=q(x),Y=q(y)$, and assume that the normal
closure of $Y$ in $G$ is all of $G$. For $N\ge2$ let
\[
 K_N=\ker(F\longrightarrow\mathbb Z/N),\qquad
 x\longmapsto1,\quad y\longmapsto0.
\]
On the free basis $x^N,y_j=x^jyx^{-j}$, $0\le j<N$, define
\[
 \delta_N(x^N)=x,\quad\delta_N(y_0)=y,\quad
 \delta_N(y_j)=1\ (j>0),\qquad
 c_N(w)=\operatorname{exp}_y(\delta_N(w)).
\]
Given a finite set $\mathcal S\subset\{2,3,\ldots\}$, elements
$g_0,g_N\in G'$, and arbitrary integers $b_N$, there exists an
ordinary commutator word $f\in F'$ such that
\[
 q(f)=g_0,\qquad c_N(f)=b_N,\qquad
 q\delta_N(f)=Y^{b_N}g_N\quad(N\in\mathcal S).
\tag{D23.D1.1}
\]
For an infinite set $\mathcal S$ and arbitrary $b_N\in\widehat{\mathbb Z}$,
the same assertion holds with a profinite commutator word.

In particular, in one fixed such target the entire family of
Enriquez distribution equations can hold for a word with any
prescribed $G'$ value. This does not assert the gentle identities
for that word, or the distribution equations in other targets.

\begin{proof}
Write $L=\operatorname{lcm}(\mathcal S)$ and adjoin $L$ to
$\mathcal S$ if necessary, prescribing arbitrary extra data there.
The case of an empty set is immediate. For every positive $m$,
$q(K_m)=G$: this image is normal in $G$ and contains $Y$.
For $N\mid L$, the restriction $\delta_N(K_L)$ is exactly
$K_{L/N}$. This follows on the free basis: $x^L$ maps to
$x^{L/N}$, and $y_j$ maps to $x^{j/N}yx^{-j/N}$ if $N\mid j$
and to $1$ otherwise. Consequently
\[
 q\delta_N(K_L)=G.
\tag{D23.D1.2}
\]

First construct $b\in F'\cap K_L$ with $c_N(b)=b_N$ for all
the prescribed indices. Use the basis $x^L,y_0,\ldots,y_{L-1}$
of $K_L$. The coefficient of $y_0$ sets $c_L$ and affects every
$c_N$ by the same integer. In descending order of $N<L$ in
$\mathcal S$, the coefficient of $y_N$ sets the remaining
value of $c_N$ without altering a value already set at an index
$M>N$: indeed $c_M(y_N)=0$ for $M>N$, while $c_N(y_N)=1$.
Finally the coefficient of $y_1$ corrects the total $y$ exponent
to zero without altering any $c_N$. All these words have zero
$x$ exponent. The resulting word therefore belongs to $F'$.

Now put $V_L=K_L$. For $N<L$ put
\[
 V_N=\langle\!\langle y_N\rangle\!\rangle_{K_L}.
\]
By (D23.D1.2), the image of $V_N$ under $q\delta_N$ is the
normal closure of $XYX^{-1}$ in $G$, hence all of $G$.
For $V_L$ this follows directly from (D23.D1.2). Thus
\[
 q\delta_N(V_N')=G'.
\tag{D23.D1.3}
\]
If $M>N$ is another prescribed index, then $\delta_M(y_N)=1$
and hence $\delta_M(V_N)=1$. Every $V_N'$ lies in $F'$ and
has zero value under every homomorphism $c_M$.

Starting with $b$, proceed again in descending order of $N$.
The two elements $q\delta_N(b)$ and $Y^{b_N}g_N$ have the
same abelianization in $G$, since $\delta_N(b)$ has zero
$x$ exponent and $y$ exponent $b_N$. Their quotient therefore
belongs to $G'$. Equation (D23.D1.3) lets us correct it by
right multiplication by an element of $V_N'$. This leaves
all larger prescribed indices and all $c_M$ unchanged.
After finitely many corrections the distribution coordinates
have exactly the required values.

To set the original coordinate, put
$V_0=\langle\!\langle y_1\rangle\!\rangle_{K_L}$.
Its image under $q$ is all of $G$, because $q(K_L)=G$ and
$q(y_1)=XYX^{-1}$ normally generates $G$. Thus $q(V_0')=G'$.
But $\delta_N(V_0)=1$ for every $N\ge2$ dividing $L$.
One final factor from $V_0'$ adjusts $q(b)$ to $g_0$ without
changing any previously fixed coordinate. All constructions
are ordinary group operations with integer powers; finite
images of normal closures and their derived groups give
ordinary word representatives. This proves the finite assertion.

The same proof works in the profinite free group for finitely
many prescribed profinite integers, using their powers in the
profinite cyclic subgroups. Alternatively one can impose these
integers modulo successively larger moduli. All conditions in
(D23.D1.1) are closed conditions in the compact group
$\widehat F'$. Every finite subcollection has just been shown
to be satisfiable. Compactness therefore proves the infinite
assertion. It produces a single profinite word in this fixed
target, and does not promote that word to a gentle automorphism.
\end{proof}

\paragraph{D23.D2 (application to the compressed two-sign test).}
Let $G=E_0\times S_4$ be D23.P10's marked group of order
$5184$. Every one of its twelve necessary full-GT values
\[
 g_{\rho,e,\nu}=
 \left(\left(\binom{e}{\rho+(\rho+1)e},M_\rho\right),
               t^{2\nu}p^\rho\right),
 \quad\rho\in\mathbb F_3,\quad e,\nu\in\mathbb F_2,
\]
can be the value of one profinite word satisfying all the
Enriquez distribution equations in $G$, simultaneously for
all positive $N$, with arbitrarily prescribed $c_N$ for $N\ge2$.
In particular the fixed-target distribution equations do not
discard any of the six mixed-sign states. No actual gentle or
full-GT lift of these words is asserted.

\begin{proof}
The group is solvable. Its abelianization is $C_3\times C_2$,
and the image of $Y$ generates it: the affine linear part
$-L^2$ generates $\operatorname{SL}_2(\mathbb F_3)^{\rm ab}=C_3$,
and the four-cycle $t$ is odd. Hence the quotient of $G$ by
the normal closure of $Y$ is both solvable and perfect, so
is trivial. Every displayed state belongs to
$G'=(\mathbb F_3^2\rtimes Q_8)\times A_4$; this is also
checked in D23.P10's six-row certificate. Apply D23.D1 with
$g_0=g_N=g_{\rho,e,\nu}$ for every $N\ge2$. The case $N=1$
is the identity equation with $c_1=0$.
\end{proof}

\paragraph{Exact source and computational checks.}
The source theorem is B. Enriquez, \emph{Quasi-reflection algebras
and cyclotomic associators}, Selecta Math. 13 (2007),
Section 8, Theorem 8.1 and equation (50); arXiv:math/0408035v2,
printed pp.~30--32. It states these distribution equations
on the full profinite GT group. Remark 8.6 gives the precise
identity $c_{MN}=c_M+c_N$ under multiplication of the integer indices; that requirement
can be imposed on the arbitrary prescribed values above.
The supplied LM7\_PUSH E.6 already used this theorem in a
different finite target. The general normal-generation argument
here is a new proof of the stated fixed-target limitation.
As an independent finite check, \path{code/check_distribution_216.py}
constructs the induced wreath actions and verifies every one of
the twelve states for $N=2,3,4,6$. At each index the joint
projected commutator image has order $8957952=864^2\cdot12$;
all twelve choices of $c_N\bmod12$ occur. The proof above,
not those four finite checks, establishes the simultaneous
all-index statement.

\subsection{A cubic descent and the exact remaining cubic defect}

\begin{statement}{D23.F22}{Proved}{The index-three half-character is a descended cubic frame sign}
\end{statement}
There is a canonical continuous endomorphism
\[
 \mathcal D_E:\mathcal H\longrightarrow\mathcal H
\]
of the multiplier-one actual gentle group with the following explicit
description. In $B=\langle s,r\mid s^2=r^3=1\rangle^{\wedge}$ put
$\tau=sr$, $x=\tau^{-2}$, $y=sxs$, and
\[
 E=\langle r,srs\rangle\cong C_3\amalg C_3,
 \qquad C_f=f(x,r),\qquad t_f=\frac{3\Psi_3(f)}2.
\]
For $h=(1,f)\in\mathcal H$ first form the automorphism
\[
 \psi_f(r)=r,\qquad
 \psi_f(s)=f(r,srs)^{-1}s.
\]
Then
\[
 \boxed{\quad
 \mathcal D_E(h)=\operatorname{Inn}(x^{-t_f}C_f^{-1})\psi_f,
 \qquad A_{\mathcal D_E(h)}=x^{-t_f}C_f^{-1}.
 \quad}                                                     \tag{D23.F22.1}
\]
In particular
\[
 \boxed{\quad
 a_3(\mathcal D_E(h))=\Psi_3(h)/2\pmod2,\qquad
 \rho(\mathcal D_E(h))=\rho(h)\pmod3.
 \quad}                                                     \tag{D23.F22.2}
\]
Thus the cubic comparison defect of \textup{A.103}, on the actual
full-pentagon subgroup, has the intrinsic form
\[
 e_3=a_3+a_3\circ\mathcal D_E.
                                                               \tag{D23.F22.3}
\]
The fixed-point condition has the exact primitive formulation
\[
 \mathcal D_E(h)=h
 \quad\Longleftrightarrow\quad
 A_h^{-1}=f(x,r)x^{3\Psi_3(h)/2}.                              \tag{D23.F22.4}
\]
This is the equianharmonic primitive form with a specifically
normalized coefficient. Neither the fixed-point condition nor the
weaker equality $a_3=a_3\circ\mathcal D_E$ is asserted on all full
profinite GT here.

\begin{proof}
The kernel of the map $B\to C_2$ sending $s$ to the nonidentity
element and $r$ to one is $E=C_3\amalg C_3$, with the displayed
generators. This follows directly by the two-coset Reidemeister
rewriting, and the presentation holds profinitely as well. The
standard pure subgroup $F$ has index three in $E$; $r$ splits
$E\to C_3$, and conjugation by $r$ cyclically permutes
$x,y,(xy)^{-1}$.

The input automorphism $X\mapsto X$, $Y\mapsto f(X,Y)^{-1}Yf(X,Y)$
of the free profinite group preserves the closed normal subgroup
generated by $X^3,Y^3$, as does its inverse. Hence it induces an
automorphism of $E$, with $X=r,Y=srs$. Put $g=f(r,srs)$.
The two-cycle identity gives $sgs=g^{-1}$, so $(g^{-1}s)^2=1$.
Consequently the displayed formulas define an endomorphism $\psi_f$
of $B$. Its restriction to $E$ is exactly the induced automorphism,
and it induces the identity on $B/E$. It is therefore an
automorphism: surjectivity and injectivity follow successively from
the quotient and the restriction. The pair-composition rule
$f_{hk}=f_h\Phi_h(f_k)$ proves
$\psi_{f_{hk}}=\psi_{f_h}\psi_{f_k}$ by evaluating it in $E$.
In particular $h\mapsto\psi_f$ is a continuous homomorphism.

Both $g$ and $C_f$ belong to $F$: in $E/F=C_3$ the image of any
commutator word is one. Hence $\psi_f$ induces the identity on
$B/F=S_3$. The peripheral identity of \textup{A.100}, applied to
the pair $(r,srs)$, gives
\[
 \psi_f(y)=f(y,r)y f(y,r)^{-1},
 \qquad y=(rs)^{-2}.
\]
As $r^{-1}yr=x$ and $\psi_f(r)=r$, this gives
$\psi_f(x)=C_fxC_f^{-1}$.
It also gives $\psi_f(\tau)=C_f\tau C_f^{-1}$, with no choice
of square root. Indeed $v=C_f^{-1}\psi_f(\tau)C_f$ satisfies
$v^2=\tau^2$ and has the same image as $\tau$ in $S_3$.
The element $v\tau^{-1}\in F$ centralizes $x=\tau^{-2}$.
The pure cusp centralizer is $C_F(x)=\overline{\langle x\rangle}$;
thus $v=\tau x^b$ for some $b\in\widehat{\mathbb Z}$.
The equality of squares gives $2b=0$, whence $b=0$.

We compute the required integral normalization rather than choosing
it arbitrarily. Identify the $r$-module $F^{\rm ab}$ with
$\widehat{\mathbb Z}[\omega]$, where $[x]=1$, $[y]=\omega$,
$\omega^2+\omega+1=0$, and $r$ acts by multiplication by $\omega$.
In $F^{\rm ab}\rtimes\langle r\rangle$ the elements $r,x,srs$
are respectively
\[
 (0,\omega),\qquad(1,1),\qquad(-1,\omega^{-1}).
\]
The argument of \textup{A.90} works over the entire coefficient
ring $\widehat{\mathbb Z}[\omega]$: $1-\omega$ is a non-zero-divisor
because its norm is $3$, which is a non-zero-divisor in
$\widehat{\mathbb Z}$. The peripheral identity therefore gives
$A_f(\omega,\omega^{-1})=0$. The two-cycle and three-cycle then
give that the translation of
$f((0,\omega),(1,1))$ is a scalar $t\in\widehat{\mathbb Z}$.
Explicitly the two nontrivial translations in the three-cycle are
$-\bar t$ and $t$, so $t=\bar t$; the fixed ring is
$\widehat{\mathbb Z}$, again because multiplication by
$\omega-\omega^{-1}$ has nonzero norm $3$.

If $m_0,m_1,m_2$ are the three index-three Schreier coefficients,
the translation is $t=m_0+m_1\omega+m_2\omega^2$, with
$m_0+m_1+m_2=0$. Taking its trace gives
$2t=3m_0=3\Psi_3$. Since \textup{A.90} proves that $\Psi_3$ is
even, $t=t_f=3\Psi_3/2$.
The two-cycle now gives
\[
 C_f^{\rm ab}=-t_f[x],\qquad g^{\rm ab}=0.                    \tag{D23.F22.5}
\]
For the second assertion, evaluate the already vanishing expression
at $(0,\omega),(-1,\omega^{-1})$; changing the translation by a
scalar factor changes a commutator-word translation by the same
factor.

Let $\theta_f=\operatorname{Inn}(C_f^{-1})\psi_f$. It fixes
$\tau$, and its normalized pair coordinate is
\[
 f_0=s\theta_f(s)=(sC_f^{-1}s)gC_f.
\]
Conjugation by $s$ exchanges $[x]$ and $[y]$, so (D23.F22.5)
gives $f_0^{\rm ab}=(-t_f,t_f)$.
Conjugating once more by $x^{-t_f}$ changes its pair coordinate to
\[
 f_E=y^{-t_f}f_0x^{t_f}\in F'.
\]
The resulting automorphism fixes $\tau$, induces the identity on
$S_3$, and has commutator pair coordinate; equivalently it is an
actual multiplier-one gentle automorphism. It sends $r$ to its
conjugate by $x^{-t_f}C_f^{-1}\in F$, so uniqueness of the pure
cubic frame proves (D23.F22.1).

There is only one such normalized gentle representative in this
pure inner class. If two differ by $\operatorname{Inn}(u)$ with
$u\in F$, then fixing $\tau$ forces $u\in C_F(x)=\langle x\rangle$.
Writing $u=x^n$ changes a commutator pair coordinate by the
abelian vector $(-n,n)$, so $n=0$. The homomorphism law for
$\psi_f$ and this uniqueness prove the homomorphism law for
$\mathcal D_E$. Its formula and the continuity of $\Psi_3$ give
continuity.

The sign computation must retain the signed lift. Put
\[
 U=\begin{pmatrix}1&1\\0&1\end{pmatrix},\quad
 L=\begin{pmatrix}1&0\\-1&1\end{pmatrix},\quad
 S=ULU,\quad R=-UL
 \qquad\hbox{over }\mathbb F_3.
\]
Define $q_E:E\to\operatorname{SL}_2(\mathbb F_3)$ by
$r\mapsto R$ and $srs\mapsto SRS^{-1}$. Both images have order
three, so this is well-defined. Direct multiplication gives
\[
 q_E(x)=-U^2,\qquad q_E(y)=-L^2.
\]
Consequently $q_E(w)=(-I)^{\ell(w)}q_3(w)$ on $F$, where
$\ell(w)$ is the sum modulo two of its two abelian coordinates,
and $q_3$ is the positive signed free-subgroup lift of \textup{AD.2}.
In particular $q_E=q_3$ on the pure cubic frame in (D23.F22.1),
whose abelianization is zero by (D23.F22.5).

Write $\rho=\rho(h)\bmod3$ and
$M_\rho=\left(\begin{smallmatrix}1&\rho\\\rho&1+\rho^2\end{smallmatrix}\right)$.
The supplied signed covariance calculation gives
$f(U^2,L^2)=M_\rho$. Since $f$ has zero exponent sums, a central
sign in either input does not change its value. Using
$R=U^2L^2U^{-2}$ gives
\[
 C_E:=q_E(C_f)=f(-U^2,R)=U^2M_\rho U^{-2}.                   \tag{D23.F22.6}
\]
Put $b=\Psi_3(h)/2$, so $t_f=3b$. Because $U^3=I$,
\[
 q_3(A_{\mathcal D_E(h)})=q_E(A_{\mathcal D_E(h)})
       =(-I)^bC_E^{-1}.
\]
For $\rho_D=\rho(\mathcal D_E(h))\bmod3$, the corrected
cubic frame is therefore
\[
 T_{\mathcal D_E(h)}=(-I)^b C_E U^{2\rho_D}.
\]
It centralizes $R$ by \textup{AD.2}. Exact multiplication for
$\rho=0,1,2$ gives
$C_EU^{2\rho}=R^\rho$. Since no nonidentity power of $U$
centralizes $R$, this forces $\rho_D=\rho$, and hence
\[
 T_{\mathcal D_E(h)}=(-I)^bR^\rho.
\]
The sign character of $\langle-I\rangle\times\langle R\rangle$
is precisely $a_3$, proving (D23.F22.2). All signed identities,
all three shear residues, and both values of $b\bmod2$ are
independently checkable in \texttt{cubic\_descent.py}; its record
is \texttt{cubic\_descent.json}.

On full GT, \textup{A.76} identifies $a_3$ with
$\kappa_d\bmod2$, and \textup{A.103} defines
$e_3=\kappa_d+\Psi_3/2$. This proves (D23.F22.3).
Finally two normalized gentle automorphisms fixing $\tau$ and
having the same pure $r$-frame coincide: their images of $r$
coincide and $s=\tau r^{-1}$. Applying (D23.F22.1) proves
(D23.F22.4).
\end{proof}

\paragraph{Source boundary.}
The normalizations, actual gentle group, pure-frame uniqueness, and
the positive signed covariance are the supplied \textup{R3--R4},
\textup{LM7\_PUSH E.65}, and \textup{DELTA\_ON\_GT AD.1--AD.2}.
The character definition, evenness and peripheral identity are
\textup{A.88--A.90} and \textup{A.100}.
The new assertion here is the cubic endomorphism and its exact
parity formula. \textup{A.65}, equation \textup{(M8)}, records the
equianharmonic primitive form with an existential even coefficient;
the computation above identifies the coefficient for this descent.
The primary geometric account of the two fillings is
Lochak--Nakamura--Schneps, \emph{Eigenloci of 5 point configurations},
\S6, equations (6.3)--(6.4), author PDF pp.30--31, and its
Theorem 2, equation (HF1), author PDF p.4. That theorem has Galois
scope. Enriquez, \emph{Quasi-reflection algebras and cyclotomic
associators}, arXiv:math/0408035v2, Theorem 8.1 and equation (50),
pp.30--32, supplies the full-GT cyclic distribution relation but
is not silently substituted for the missing fixed-point or parity
comparison in (D23.F22.4).

\subsection{The harmonic sign is unchanged by cubic descent}

\paragraph{D23.P13 (an actual-source relation between the two descents).}
For every actual multiplier-one gentle element $h=(1,f)$,
the cubic descent of D23.F22 satisfies
\[
 \boxed{\quad a_2(\mathcal D_Eh)=a_2(h),\qquad
 \rho(\mathcal D_Eh)=\rho(h)\pmod3.\quad}
\tag{D23.P13.1}
\]
More precisely its normalized pair coordinate $f_E$ has exactly
the same value as $f$ in the signed affine harmonic detector of
supplied FP.34. Consequently, using D23.F22 and supplied A.107,
\[
 \boxed{\quad
 \Delta\circ\mathcal D_E
 =a_2+\Psi_3/2
 =d_2
 =a_2+a_2\circ D_L
 =a_2+a_2\circ D_P
 \quad\text{on }\mathcal H.\quad}
\tag{D23.P13.2}
\]
The half-character in this display is reduced modulo two.
This identifies an actual descended defect; it does not assert
that $d_2$ vanishes, or that either descent fixes all full-GT
elements.

\begin{proof}
We keep the exact FP.34 marking. Put
\[
 R=(\mathbb Z/4)[\omega]/(\omega^2+\omega+1),\qquad
 c=(1-\omega)^{-1}=2+3\omega,
\]
and use the affine multiplication $(a,u)(b,v)=(a+ub,uv)$.
The homomorphism $B=C_2\amalg C_3\to R\rtimes\langle-\omega\rangle$
is
\[
 s\longmapsto(c,-1),\quad r\longmapsto(c,\omega),\quad
 x\longmapsto(0,\omega),\quad y\longmapsto(1,\omega).
\tag{D23.P13.3}
\]
Write $f((0,\omega),(1,\omega))=(k,1)$.

For any commutator word $f$ and any $a,b\in R$, the exact
affine identity
\[
 f((a,\omega),(b,\omega))=((b-a)k,1)
\tag{D23.P13.4}
\]
holds. Indeed simultaneous conjugation by the translation
$ac$ changes $(0,\omega),(b-a,\omega)$ to the indicated
pair. A commutator word has linear part one, and conjugation
by a translation leaves a pure translation unchanged.
Finally multiplying the translation module by $b-a$ is an
affine-group homomorphism (no invertibility of $b-a$ is needed),
which gives (D23.P13.4). The argument applies to profinite
words by continuity, since the target is finite.

Use the notation and normalized formula from D23.F22:
\[
 C=f(x,r),\qquad g=f(r,srs),\qquad t=3\Psi_3(f)/2,
 \qquad f_E=y^{-t}(sC^{-1}s)gCx^t.
\tag{D23.P13.5}
\]
Here $C,g\in F$ and $t\in3\widehat{\mathbb Z}$; the latter
uses the already proved evenness of $\Psi_3$.
Since $srs\mapsto(-\omega c,\omega)$, formula (D23.P13.4) gives
\[
 C\longmapsto(ck,1),\qquad
 g\longmapsto(\omega^2ck,1).
\tag{D23.P13.6}
\]
Conjugation by $s$ negates translations, hence
$sC^{-1}s\mapsto(ck,1)$. Therefore
\[
 (sC^{-1}s)gC\longmapsto((2+\omega^2)ck,1)=(k,1),
\]
because $2+\omega^2=1-\omega$ and $(1-\omega)c=1$ in $R$.
Both displayed images of $x,y$ have order three. The normalizing
powers $x^t,y^{-t}$ thus have identity image, and
\[
                    q_{\rm harm}(f_E)=q_{\rm harm}(f).
\tag{D23.P13.7}
\]
This is an exact equality of finite values, not just a comparison
after forgetting a sign or reducing modulo two.

FP.34 gives, for each actual gentle normalized pair,
\[
 1-(1-\omega)k=\omega^\rho(-1)^{a_2}.
\]
The six elements $\omega^j(-1)^e$, $j\in\mathbb F_3$,
$e\in\mathbb F_2$, are distinct in $R$. Applying the formula
to $f_E$ and $f$ proves both claims in (D23.P13.1), independently
of the second congruence proved by the cubic matrices in D23.F22.

D23.F22 gives $a_3\circ\mathcal D_E=\Psi_3/2$. Adding
(D23.P13.1) proves the first two equalities of (D23.P13.2).
The remaining equalities are exactly supplied A.107, equations
(MD.8a)--(MD.8b). This deduction uses genuine endomorphisms
and actual characters on $\mathcal H$ throughout; it makes
no assertion about full-GT descent compatibility.

The exact certificate \path{agents/pentagon/cubic_descent_harmonic.py}
checks the ring, all signed peripheral images, the affine group
order, the translation computation for every $k\in R$, and
the uniqueness of the six harmonic phases. It does not replace
the cited actual-source descent construction in D23.F22 or the
supplied harmonic-frame theorem FP.34.
\end{proof}

\paragraph{Exact source locators and remaining question.}
The signed harmonic representation and recovery formula are
\path{sources/DELTA_ON_GT/agent_five/five_point_results.tex},
FP.34.1--FP.34.3. The pair formula is
\path{work/agents/five_level_six/cubic_descent.tex},
D23.F22.1 and its displayed normalization of $f_E$.
The half-character comparison with $D_L,D_P$ is
\path{sources/DELTA_ON_GT/sections/A.tex}, A.107,
(MD.8a)--(MD.8b).
The identity $\Psi_3/2\circ\mathcal D_E=\Psi_3/2$, or a different
formula for this iterated half-character, is not proved here.

\subsection{The cubic descent for all multipliers}

\begin{statement}{D23.F23}{Proved}{The whole gentle-group descent and its relative quadratic character}
\end{statement}
The endomorphism $\mathcal D_E$ of \textup{D23.F22} extends canonically
to a continuous endomorphism of the full actual gentle group
$\mathcal G=\widehat{GT}_0$, preserving its multiplier.
For $g=(\lambda,f)$ put
\[
 m=(\lambda-1)/2,\quad
 \varepsilon\in\{1,-1\},\quad\varepsilon\equiv\lambda\pmod3,
 \quad b=(1-\varepsilon)/2,
 \qquad C_g=r^{-m}f(x,r).
\]
Define
\[
 \psi_g(r)=r^\lambda,\qquad
 \psi_g(s)=f(r,srs)^{-1}s,\qquad
 \theta_g=\operatorname{Inn}(\tau^b C_g^{-1})\psi_g.
\]
Then $f_{0,g}=s\theta_g(s)$ belongs to $F$ and has abelianization
$(u_g,-u_g)$. The full descent is the explicit normalized automorphism
\[
 \boxed{\quad
 \mathcal D_E(g)=\operatorname{Inn}(x^{u_g})\theta_g,
 \quad f_{\mathcal D_E(g)}=y^{u_g}f_{0,g}x^{-u_g}\in F'.
 \quad}                                                     \tag{D23.F23.1}
\]
This construction is independent of any Galois lift of the multiplier.

Using \textup{D23.P13}, put $\gamma=\Delta\circ\mathcal D_E$.
Its restriction to $\mathcal H=\ker\lambda$ is the duplication
defect $d_2=a_2+\Psi_3/2$. There is a unique continuous character
$\chi:\widehat{\mathbb Z}^{\times}\to\mathbb F_2$ such that
$\gamma|_{G_{\mathbb Q}}=\chi\circ\lambda$.
Consequently
\[
 \widehat d_2=\gamma+\chi\circ\lambda,\qquad
 \widehat e_3=\Delta+\widehat d_2                             \tag{D23.F23.2}
\]
are canonical continuous relative quadratic characters of
$\mathcal G$, restricting respectively to
\[
 a_2+\Psi_3/2,\qquad a_3+\Psi_3/2
\]
on $\mathcal H$. In particular
$\Delta=\widehat d_2+\widehat e_3$ on the full gentle group.
This does not assert nontriviality or independence of these two
characters, and does not determine the whole relative character space.
It also does not assert that the uncorrected $\gamma$ vanishes on
all of $G_{\mathbb Q}$.
By uniqueness of relative extension, $\widehat d_2$ is the
character $d$ already constructed in D23.C1. The new assertion
is its realization through the whole-group cubic descent.

\begin{proof}
As in \textup{D23.F22}, quotient the input free profinite group by
the closed normal subgroup generated by $X^3,Y^3$. The quotient
is $E=C_3\amalg C_3$, with $X=r,Y=srs$. Both the actual input
automorphism and its inverse preserve this normal subgroup, since
$\lambda$ is a unit. They therefore induce an automorphism of $E$.
Write $G=f(r,srs)$. The two-cycle gives $sGs=G^{-1}$, so
$(G^{-1}s)^2=1$. The formulas for $\psi_g$ define an endomorphism
of $B$, whose restriction to $E$ is the induced automorphism and
whose action on $B/E=C_2$ is the identity. It is an automorphism.
The normalized pair-composition formula proves directly that
$g\mapsto\psi_g$ is a homomorphism, including arbitrary multipliers.

For clarity we derive the general peripheral conjugator. Write
$Z=(XY)^{-1}$, $G_0=f(Z,X)$, $H_0=f(Y,Z)$ and
$m=(\lambda-1)/2$. The three-cycle for $(\lambda,f)$ is
\[
 G_0 Z^m H_0 Y^m f X^m=1.
\]
Right composition with the actual element $c=(-1,1)$ gives the
pair $(-\lambda,f)$ and its three-cycle
\[
 G_0 Z^{-m-1}H_0Y^{-m-1}fX^{-m-1}=1.
\]
Eliminating $H_0$ from these equations gives
\[
 G_0Z^{-\lambda}G_0^{-1}
   =X^{m+1}f^{-1}Y^\lambda fX^m
   =X^{-m}\Phi_g(Z^{-1})X^m.
\]
Thus
\[
 \Phi_g(Z)=X^m f(Z,X)Z^\lambda f(Z,X)^{-1}X^{-m}.             \tag{D23.F23.3}
\]
Evaluate this at $(r,srs)$, where $Z=y=(rs)^{-2}$.
Since $x=r^{-1}yr$ and $\psi_g(r)=r^\lambda$, we obtain
\[
 \psi_g(x)=r^{m-\lambda}f(y,r)y^\lambda f(y,r)^{-1}r^{\lambda-m}
          =C_g x^\lambda C_g^{-1}.
\]
In the last equality use $m-\lambda+1=-m$ and
$f(y,r)r=rf(x,r)$.

The element $f(x,r)$ is pure, and $\psi_g$ induces on $S_3$
the automorphism $s\mapsto s$, $r\mapsto r^\varepsilon$.
Consequently $v=C_g^{-1}\psi_g(\tau)C_g$ has the same
permutation as $\tau^\lambda=\tau\bmod F$ and satisfies
$v^2=\tau^{2\lambda}$.
Both $v$ and $\tau^\lambda$ commute with $x^\lambda$.
Since $\lambda$ is a unit, they commute with $x$ as well.
Therefore $v\tau^{-\lambda}$ lies in
$C_F(x)=\overline{\langle x\rangle}$.
The equality of squares and injectivity of multiplication by two
in $\widehat{\mathbb Z}$ show that it is one. We have proved
\[
 \psi_g(\tau)=C_g\tau^\lambda C_g^{-1}.
                                                               \tag{D23.F23.4}
\]

The correction by $\tau^b$ is essential for the permutation
normalization. The two possible cases are
\[
\begin{array}{c|c|c|c|c}
 \lambda\bmod6&\varepsilon&m\bmod3&b&
       \tau^b C_g^{-1}\bmod F\\\hline
 1&1&0&0&1\\
 5&-1&2&1&(sr)r^2=s.
\end{array}
\]
In the first case $\psi_g$ already acts identically on $S_3$.
In the second it fixes $s$ and inverts $r$, and conjugation by
$s$ changes this into the identity on $S_3$. Thus in both cases
$\theta_g$ induces the identity on $S_3$; by (D23.F23.4) it
sends $\tau$ to $\tau^\lambda$.

It follows that $f_{0,g}=s\theta_g(s)$ belongs to $F$.
The order-two relation $\theta_g(s)^2=1$ gives the exact identity
$s f_{0,g}s=f_{0,g}^{-1}$. On $F^{\rm ab}$, conjugation by $s$
exchanges the two coordinates. Hence if
$f_{0,g}^{\rm ab}=(u,v)$, then $(v,u)=(-u,-v)$, or $v=-u$.
This proves the asserted abelian form. Conjugation by $x^{u_g}$
changes the pair coordinate to $y^{u_g}f_{0,g}x^{-u_g}$, whose
abelianization is zero, and preserves the cusp formula.

The normalized automorphism restricts to the actual free-group
automorphism
$x\mapsto x^\lambda$, $y\mapsto f_E^{-1}y^\lambda f_E$.
Its two-cycle follows from the order-two relation just used.
Its three-cycle follows from the order-three relation in $B$:
its image of $r$ is $f_E^{-1}r x^{-m}$, and expansion of its
cube, using the cyclic action of $r$ on $x,y,(xy)^{-1}$, gives
\[
 f_E(z,x)z^m f_E(y,z)y^m f_E(x,y)x^m=1.
\]
Thus (D23.F23.1) is an actual element of $\mathcal G$ of
multiplier $\lambda$.

For the composition law, each (D23.F23.1) is an inner correction
of $\psi_g$. Since $\psi$ is a homomorphism, its normalized
corrections have the same outer composition law. Two such
corrections with the same multiplier both induce the identity
on $S_3$, so an inner element relating them has trivial image
in $S_3$, whose center is trivial. It is therefore in $F$.
Since both automorphisms send $\tau$ to $\tau^\lambda$, that
element centralizes $x$ and is $x^n$. Conjugation by $x^n$
changes the pair coordinate by the abelian vector $(-n,n)$.
Both coordinates belong to $F'$, so $n=0$. This proves uniqueness
of the normalization and hence the homomorphism law.
All operations, the locally constant choice of $b$, and extraction
of $u_g$ from $F^{\rm ab}$ are continuous, proving continuity.
At $\lambda=1$ we have $b=m=0$ and, by \textup{D23.F22},
$u_g=-3\Psi_3(g)/2$. Thus the extension restricts to that of
\textup{D23.F22}.

Now $\gamma=\Delta\mathcal D_E$ is a continuous character on
$\mathcal G$. Statements \textup{D23.F22} and \textup{D23.P13}
give $\gamma|_{\mathcal H}=d_2$.
The arithmetic comparison needed here is precisely
\textup{DELTA\_ON\_GT A.103}, final paragraph after
\textup{(MD.4e)}: on $G_{\mathbb Q}\cap\mathcal H$,
Anderson duplication in the specified intrinsic convention gives
$d_2=0$. Equivalently this is the physical-descent arithmetic
compatibility of \textup{A.62}, combined with the universal
half-character formula \textup{A.107 (MD.8a)}.
We use this supplied arithmetic comparison, and do not replace
it by an unmarked assertion about a Kummer exponent.

Therefore $\gamma|_{G_{\mathbb Q}}$ vanishes on the kernel of
the cyclotomic character. The cyclotomic map
$G_{\mathbb Q}\twoheadrightarrow\widehat{\mathbb Z}^{\times}$
is a continuous surjection of compact profinite groups, hence
its target is the topological quotient by that kernel.
There is consequently a unique continuous character $\chi$
with the stated factorization. The formulas (D23.F23.2) define
continuous characters and show directly that they vanish on
$G_{\mathbb Q}$. Their restrictions to $\mathcal H$ are as
claimed. Their uniqueness as relative extensions of those
restrictions follows from
$\mathcal G=\mathcal H G_{\mathbb Q}$, the same decomposition
used in \textup{LM7\_PUSH E.65}. This also proves the last
displayed identity on all of $\mathcal G$.
\end{proof}

\paragraph{Geometric normalization.}
For the cubic eigenlocus put
$n=(1+\sqrt{-3})/2$, $\bar n=(1-\sqrt{-3})/2$ and
\[
 t=\frac{\lambda-\bar n}{\lambda-n},\qquad b=1-t^3.
\]
Direct substitution in the Legendre expression gives
\[
 \frac{256(1-\lambda+\lambda^2)^3}
      {1728\lambda^2(1-\lambda)^2}
       =\frac{4(b-1)}{b^2}.                                \tag{D23.F23.5}
\]
The pole values $n,\bar n$ correspond to $b=\infty,1$;
the three cusps $\lambda=0,1,\infty$ correspond to $b=0$.
The degree-two map on the right is the harmonic map
$4u(1-u)$ after $u=1/b$. This records the precise relation
between the $C_3\amalg C_3$ and $C_2\amalg C_3$ orbifolds.
It does not by itself identify the two GT actions: the cyclic
distribution filling retains $0,1,\infty$ in the $t$-sphere,
whereas the cubic modular filling removes the two poles and
retains the three roots of unity. The functorial comparison for
these different fillings remains the required additional step.


\part{Proof drafts and active calculations awaiting their final check}
In this part a proof draft is not yet counted as a completed
finding.  Exact computation records are retained even when the
corresponding geometric or source comparison is still open.
\section{Current five-point calculation}
The full $S_3$ Nielsen closure of the marked Krammer pair
evaluations is being analyzed. The question is whether its
additional simultaneous two-cycle and three-cycle identities
discard any of the common-word assignments proved above.
The explicit cyclic-cubic diagram of D23.C14 also leaves a
comparison between two different forgetting maps from the same
five-punctured sphere. D23.F22 constructs the corresponding cubic
descent and identifies its cubic sign, but the necessary comparison
of original and descended signs on full GT is still open.
The two Hurwitz tests of D23.P11--P12 admit an explicit mixed-sign
word even after their full Nielsen closure. No surviving finite
candidate is asserted to be an element of $\GT$.

\section{Literature and historical locators}
\small
\begin{enumerate}
\item J. Voight, \emph{Quaternion algebras}, stable version
1.0.5 (2024): (12.4.10), 12.4.12, (12.4.14), Proposition
14.2.1, Main Theorem 14.6.1, Corollary 14.6.2.
\url{https://jvoight.github.io/quat-book-v1.0.5.pdf}.
\item T. Kubota, \emph{Topological covering of SL(2) over a
local field}, J. Math. Soc. Japan 19 (1967), 114--121:
equation (2), p.~114; product formula (4), p.~115.
\url{https://doi.org/10.2969/jmsj/01910114}.
\item H. Hida, \emph{Adjoint L-value as a period integral and
the mass formula of Siegel--Shimura}, 20 September 2023:
Section 1.3, p.~7, the $s_p$ and $\kappa_p$ displays before
(1.5).  The splitting is attributed there to Gelbart,
\emph{Weil's Representation and the Spectrum of the
Metaplectic Group}, LNM 530 (1976), Proposition 2.8.
\url{https://www.math.ucla.edu/~hida/PdFinal.pdf}.
\item W. Herfort and L. Ribes, \emph{Torsion elements and
centralizers in free products of profinite groups}, J. reine
angew. Math. 358 (1985), 155--161: Theorem B and Section 4,
Theorem B', p.~159. \url{https://doi.org/10.1515/crll.1985.358.155}.
\item B. Collas, \emph{Action of a Grothendieck--Teichm\"uller
group on torsion elements of full profinite mapping class groups
in genus zero}, J. Th\'eor. Nombres Bordeaux 24 (2012),
605--622: Theorem B, p.~608.
\url{https://www.numdam.org/article/JTNB_2012__24_3_605_0.pdf}.
\item I. Marin, \emph{Characters of the Grothendieck--Teichm\"uller
group through rigidity of the Burau representation}, Ann. Fac.
Sci. Toulouse 17 (2008), 765--780: Section 4, Lemmas 4.1--4.3;
Section 5 for the general rigidity method.  D23.T1 proves the
precise integral specialization needed here.
\url{https://www.numdam.org/article/AFST_2008_6_17_4_765_0.pdf}.
\item S. Bigelow, \emph{The Lawrence--Krammer representation},
arXiv:math/0204057v1 (2002), Section 4.2, printed p.~13:
the corrected six-case polynomial Krammer matrix formula.
\url{https://arxiv.org/pdf/math/0204057}.
\item J. S. Milne, \emph{Lectures on Etale Cohomology}, v2.21
(22 March 2013), Theorem 3.4.
\url{https://www.jmilne.org/math/CourseNotes/LEC.pdf}.
\item J. S. Milne, \emph{Modular Functions and Modular Forms},
v1.10, Theorem 12.27, p.~124, and the surrounding ray-class
construction. \url{https://www.jmilne.org/math/CourseNotes/MF110.pdf}.
\item H. Nakamura, \emph{On arithmetic monodromy representations
of Eisenstein type in fundamental groups of once punctured
elliptic curves}, RIMS-1691 (2010), Section 5.4, equation
(5.4.2), p.~28.
\url{https://www.kurims.kyoto-u.ac.jp/preprint/file/RIMS1691.pdf}.
\item R. Hain and M. Matsumoto, \emph{Weighted completion of
Galois groups and Galois actions on the fundamental group of
$\mathbf P^1-\{0,1,\infty\}$}, Compositio Math. 139 (2003),
119--167: Section 8.5, paragraph before Theorem 8.7.
\url{https://arxiv.org/abs/math/0006158}.
\item J.-P. Serre and J. Tate, \emph{Good reduction of abelian
varieties}, Ann. Math. 88 (1968), 492--517: Theorem 1, p.~493.
\url{https://doi.org/10.2307/1970722}.
\item A. Grothendieck, Montpellier archive cote 145,
\emph{Printemps 81 (vieille r\'edaction)}: inspected one-based
PDF pages 7, 11, 12, 22, 28; numbered archive scans
6, 10, 11, 21, 27.  The PDF includes a cover, so the two
numbering conventions differ by one.
\url{https://grothendieck.umontpellier.fr/145.pdf}.
The exact historical interpretation and its limits are in
\path{IMPLICATIONS.md}.
\item L. Alp\"oge, \emph{Remark}, supplied four-page PDF:
the prompt's steps 3 and 5 and the Gauss-sum discussion on
p.~3, footnote 3.  Accepted as stipulated; no theorem from
its Lean archive is used here, and those Lean files were not
in the supplied ZIP.  Thus no formal-file locator is invented.
\end{enumerate}
\end{document}
